% El deporte, el ocio y el bienestar — Espagnol 1ère A — Cours signé OnBuch+
% Programme officiel MINESEC. Compiler avec : tectonic EA07-deporte-ocio-bienestar.tex
\newcommand{\DOCMATIERE}{Espagnol}
\newcommand{\DOCNIVEAU}{1ère A}
\newcommand{\DOCMODULE}{Módulo 2 : Environnement, santé et bien-être}
\newcommand{\DOCLECON}{EA07}
\newcommand{\DOCTITRE}{El deporte, el ocio y el bienestar}
% OnBuch+ — préambule « cours signés » ENRICHI (compilé avec tectonic / XeLaTeX)
% Système visuel « Pop » d'OnBuch+ : fond crème (jamais blanc), bordures encre
% épaisses, ombres « dures » décalées, titres Archivo Black en capitales.
% Version MATHÉMATIQUES : graphes (pgfplots/tikz), boîtes pédagogiques
% (définition, propriété, méthode, attention, le savais-tu, exercice résolu,
% exercice, TP, bilan), page de garde et signature OnBuch+ ancrée en bas.
%
% Macros attendues dans main.tex AVANT \input{preamble} :
%   \DOCMATIERE  \DOCNIVEAU  \DOCMODULE  \DOCLECON (numéro)  \DOCTITRE
\documentclass[11pt]{article}

\usepackage{fontspec}
\usepackage[spanish,french]{babel} % français = langue principale ; l'espagnol via \esp{..} / espagnol
\usepackage[a4paper,margin=2cm,top=3.4cm,bottom=2.8cm,headheight=1.4cm,footskip=1.3cm]{geometry}
\usepackage{amsmath,amssymb,amsfonts}
\usepackage{mathtools} % \xmapsto, \coloneqq... souvent utilisés par les modèles
\usepackage{cancel} % simplifications barrées (bilans de forces)
% \abs{z} (module, valeur absolue) et \norm{u} (norme), utilisés par les modèles
\providecommand{\abs}{}\let\abs\relax\DeclarePairedDelimiter{\abs}{\lvert}{\rvert}
\providecommand{\norm}{}\let\norm\relax\DeclarePairedDelimiter{\norm}{\lVert}{\rVert}
\usepackage[table]{xcolor} % [table] : fournit aussi \rowcolors (alternance de lignes)
\usepackage{pagecolor}
\usepackage{tikz}
\usetikzlibrary{arrows.meta,positioning,calc,shapes.geometric,shapes.misc,shapes.arrows,decorations.pathmorphing,decorations.pathreplacing,decorations.markings,patterns,shadows,mindmap,trees,fit,backgrounds,matrix,intersections,angles,quotes}
\usepackage{pgfplots}
\usepackage[version=4]{mhchem} % équations nucléaires \ce{^{235}_{92}U}
\usepackage[european,siunitx]{circuitikz} % schémas électriques
\pgfplotsset{compat=1.18}
\usepgfplotslibrary{fillbetween,groupplots} % aires entre courbes (calcul intégral)
\usepackage{enumitem}
\usepackage{fancyhdr}
% [most] charge toutes les bibliothèques tcolorbox (listings, minted, raster,
% documentation...) et gonfle inutilement la mémoire TeX sur un document long
% (~300 boîtes) : on ne charge que ce qui est réellement utilisé.
\usepackage[skins,breakable]{tcolorbox}
\usepackage{graphicx}
\usepackage{colortbl}
\usepackage{booktabs}
\usepackage{tabularx}
\usepackage{diagbox} % case d'en-tête diagonale des tableaux à double entrée
% Colonnes extensibles centrée / gauche / droite (C, L, R), souvent utilisées par les modèles
\newcolumntype{C}{>{\centering\arraybackslash}X}
\newcolumntype{L}{>{\raggedright\arraybackslash}X}
\newcolumntype{R}{>{\raggedleft\arraybackslash}X}
\usepackage{array}
\usepackage{multicol}
\usepackage{siunitx}
% parse-numbers=false : accepte aussi \SI{2,5 \times 10^{-3}}{...}
\sisetup{locale=FR,output-decimal-marker={,},group-minimum-digits=4,parse-numbers=false}
\DeclareSIUnit\ohm{\ensuremath{\symup{\Omega}}} % U+2126 absent de la police
\DeclareSIUnit\division{div} % réglages d'oscilloscope (ms/div, V/div)
\DeclareSIUnit\tr{tr} % tours (vitesse de rotation, ex. \SI{1500}{\tr\per\minute})
\DeclareSIUnit\molar{M} % concentration molaire (ex. \SI{3}{\molar}, \SI{1}{\micro\molar})
\DeclareSIUnit\an{an} % année (le modèle confond parfois avec \year, un compteur TeX) — ex. \SI{5}{\kilo\meter\per\an}
\DeclareSIUnit\FCFA{FCFA} % franc CFA (ex. \SI{500}{\FCFA\per\kilo\gram})
\DeclareSIUnit\degreeBrix{\textdegree Brix} % teneur en sucre d'un sirop/jus (réfractométrie)
\DeclareSIUnit\month{mois} % le modèle utilise parfois \month, un compteur TeX, comme unité de durée
\usepackage{qrcode}
% \degree : commande fréquemment utilisée par les modèles pour un angle ou une
% température en dehors de \SI{}{} (non fournie par les paquets chargés).
\providecommand{\degree}{^{\circ}}
% \qed (fin de démonstration), utilisé par les modèles sans amsthm
\providecommand{\qed}{\hfill$\square$}
% \barre : double barre des valeurs interdites dans un tableau de variations (tkz-tab)
\providecommand{\barre}{\ensuremath{\|}}
% environnement proof (amsthm non chargé) utilisé par les modèles
\ifdefined\proof\else\newenvironment{proof}[1][Démonstration]{\par\noindent\textit{#1.}\ }{\qed\par}\fi
\usepackage{lastpage}
\usepackage{hyperref}
\hypersetup{hidelinks}

% ============================================================
% Palette « Pop » OnBuch+ — mêmes hex que la classe P (plus_ui.dart)
% ============================================================
\definecolor{poporange}{HTML}{F59321}
\definecolor{popdark}{HTML}{E07A0C}
\definecolor{popcream}{HTML}{FFF6E8}
\definecolor{popink}{HTML}{1C1714}
\definecolor{poppurple}{HTML}{7A5AE0}
\definecolor{popgreen}{HTML}{2E9E6A}
\definecolor{poppink}{HTML}{E0457B}
\definecolor{popblue}{HTML}{2F7DE1}
\definecolor{popgold}{HTML}{C98A12}
\definecolor{popmuted}{HTML}{8A7F76}
\definecolor{popline}{HTML}{EADFD2}
\definecolor{popgray}{HTML}{8A7F76}
\colorlet{popgrey}{popgray}
% Alias tolérants (le modèle nomme parfois les couleurs différemment)
\colorlet{popwhite}{white}
\colorlet{popblack}{popink}
\colorlet{popred}{poppink}
\colorlet{popyellow}{popgold}
% Teintes claires (fonds de tableaux/encadrés) — le modèle nomme parfois
% les couleurs avec un suffixe L (ex. popblueL) sans les avoir définies.
\colorlet{poporangeL}{poporange!20}
\colorlet{popdarkL}{popdark!20}
\colorlet{poppurpleL}{poppurple!20}
\colorlet{popgreenL}{popgreen!20}
\colorlet{poppinkL}{poppink!20}
\colorlet{popblueL}{popblue!20}
\colorlet{popgoldL}{popgold!20}
\colorlet{popredL}{popred!20}
\colorlet{popgrayL}{popgray!20}
% Teintes claires pour fonds de tableaux / zones
\colorlet{popblueL}{popblue!10!white}
\colorlet{poporangeL}{poporange!14!white}
\colorlet{popgreenL}{popgreen!12!white}
\colorlet{poppurpleL}{poppurple!12!white}
\colorlet{poppinkL}{poppink!12!white}
\colorlet{popgoldL}{popgold!14!white}

\pagecolor{popcream}

% ============================================================
% Typographie
% ============================================================
\defaultfontfeatures{Ligatures=TeX}
\setmainfont{PlusJakartaSans-Medium.ttf}[Path=../../../fonts/,BoldFont=PlusJakartaSans-Bold.ttf]
% Maths en sans-serif (Fira Math) avec chiffres et lettres droites de Plus
% Jakarta, pour que les formules se fondent dans le texte.
\usepackage[math-style=ISO,bold-style=ISO]{unicode-math}
\setmathfont{FiraMath-Regular.otf}[Path=../../../fonts/]
\setmathfont{PlusJakartaSans-Medium.ttf}[Path=../../../fonts/,range={up/num,up/{Latin,latin},\mathup}]
\usepackage{icomma} % virgule décimale sans espace en mode math
% Caractères mathématiques Unicode absents des polices de texte (Plus Jakarta,
% Archivo Black) : rendus par leur équivalent mathématique, en texte comme en
% formule — sinon ils s'affichent en carré vide (ex. « Arithmétique dans ℤ »).
\usepackage{newunicodechar}
\newunicodechar{ℝ}{\ensuremath{\mathbb{R}}}
\newunicodechar{ℤ}{\ensuremath{\mathbb{Z}}}
\newunicodechar{ℂ}{\ensuremath{\mathbb{C}}}
\newunicodechar{ℕ}{\ensuremath{\mathbb{N}}}
\newunicodechar{ℚ}{\ensuremath{\mathbb{Q}}}
\newunicodechar{𝔻}{\ensuremath{\mathbb{D}}}
\newunicodechar{α}{\ensuremath{\alpha}}
\newunicodechar{β}{\ensuremath{\beta}}
\newunicodechar{γ}{\ensuremath{\gamma}}
\newunicodechar{δ}{\ensuremath{\delta}}
\newunicodechar{ε}{\ensuremath{\varepsilon}}
\newunicodechar{θ}{\ensuremath{\theta}}
\newunicodechar{λ}{\ensuremath{\lambda}}
\newunicodechar{μ}{\ensuremath{\mu}}
\newunicodechar{σ}{\ensuremath{\sigma}}
\newunicodechar{τ}{\ensuremath{\tau}}
\newunicodechar{φ}{\ensuremath{\varphi}}
\newunicodechar{ψ}{\ensuremath{\psi}}
\newunicodechar{ω}{\ensuremath{\omega}}
\newunicodechar{Σ}{\ensuremath{\Sigma}}
% majuscules produites par \MakeUppercase dans les titres (α-aminés -> Α)
\newunicodechar{Α}{\ensuremath{\alpha}}
\newunicodechar{Β}{\ensuremath{\beta}}
\newunicodechar{Γ}{\ensuremath{\Gamma}}
\newunicodechar{Δ}{\ensuremath{\Delta}}
\newunicodechar{Ε}{\ensuremath{\varepsilon}}
\newunicodechar{Θ}{\ensuremath{\Theta}}
\newunicodechar{Λ}{\ensuremath{\Lambda}}
\newunicodechar{Μ}{\ensuremath{\mu}}
\newunicodechar{Τ}{\ensuremath{\tau}}
\newunicodechar{Φ}{\ensuremath{\Phi}}
\newunicodechar{Ψ}{\ensuremath{\Psi}}
\newunicodechar{Ω}{\ensuremath{\Omega}}
\newunicodechar{↦}{\ensuremath{\mapsto}}
\newunicodechar{∈}{\ensuremath{\in}}
\newunicodechar{∉}{\ensuremath{\notin}}
\newunicodechar{∧}{\ensuremath{\wedge}}
\newunicodechar{⇔}{\ensuremath{\Leftrightarrow}}
\newunicodechar{⟺}{\ensuremath{\Longleftrightarrow}}
\newunicodechar{⇒}{\ensuremath{\Rightarrow}}
\newunicodechar{⟹}{\ensuremath{\Longrightarrow}}
\newunicodechar{∘}{\ensuremath{\circ}}
\newunicodechar{∪}{\ensuremath{\cup}}
\newunicodechar{∩}{\ensuremath{\cap}}
\newunicodechar{≡}{\ensuremath{\equiv}}
\newunicodechar{⊂}{\ensuremath{\subset}}
\newunicodechar{⊥}{\ensuremath{\perp}}
\newunicodechar{∃}{\ensuremath{\exists}}
\newunicodechar{∀}{\ensuremath{\forall}}
\newunicodechar{∛}{\ensuremath{\sqrt[3]{\,}}}
\newunicodechar{∜}{\ensuremath{\sqrt[4]{\,}}}
\newunicodechar{⃗}{\ensuremath{{}^{\to}}}
\newunicodechar{ⁿ}{\ifmmode^{n}\else\textsuperscript{\ensuremath{n}}\fi}
\newunicodechar{ˣ}{\ifmmode^{x}\else\textsuperscript{\ensuremath{x}}\fi}
\newunicodechar{ᵇ}{\ifmmode^{b}\else\textsuperscript{\ensuremath{b}}\fi}
\newunicodechar{ʳ}{\ifmmode^{r}\else\textsuperscript{\ensuremath{r}}\fi}
\newunicodechar{ᵘ}{\ifmmode^{u}\else\textsuperscript{\ensuremath{u}}\fi}
\newunicodechar{ʸ}{\ifmmode^{y}\else\textsuperscript{\ensuremath{y}}\fi}
\newunicodechar{ᵃ}{\ifmmode^{a}\else\textsuperscript{\ensuremath{a}}\fi}
\newunicodechar{ᵉ}{\ifmmode^{e}\else\textsuperscript{\ensuremath{e}}\fi}
\newunicodechar{ᵏ}{\ifmmode^{k}\else\textsuperscript{\ensuremath{k}}\fi}
\newunicodechar{ᵖ}{\ifmmode^{p}\else\textsuperscript{\ensuremath{p}}\fi}
\newunicodechar{ᴺ}{\ifmmode^{N}\else\textsuperscript{\ensuremath{N}}\fi}
\newunicodechar{ᶜ}{\ifmmode^{c}\else\textsuperscript{\ensuremath{c}}\fi}
\newunicodechar{ʰ}{\ifmmode^{h}\else\textsuperscript{\ensuremath{h}}\fi}
\newunicodechar{ᵠ}{\ifmmode^{q}\else\textsuperscript{\ensuremath{q}}\fi}
\newunicodechar{ₙ}{\ifmmode_{n}\else\textsubscript{\ensuremath{n}}\fi}
\newunicodechar{ₐ}{\ifmmode_{a}\else\textsubscript{\ensuremath{a}}\fi}
\newunicodechar{ᵢ}{\ifmmode_{i}\else\textsubscript{\ensuremath{i}}\fi}
\newunicodechar{ᵣ}{\ifmmode_{r}\else\textsubscript{\ensuremath{r}}\fi}
\newunicodechar{ₖ}{\ifmmode_{k}\else\textsubscript{\ensuremath{k}}\fi}
\newunicodechar{ᵦ}{\ifmmode_{\beta}\else\textsubscript{\ensuremath{\beta}}\fi}
\newunicodechar{ₓ}{\ifmmode_{x}\else\textsubscript{\ensuremath{x}}\fi}
\newfontfamily\popdisplay{ArchivoBlack-Regular.ttf}[Path=../../../fonts/]
\newfontfamily\poplogo{SpaceGrotesk-Bold.ttf}[Path=../../../fonts/]
\newfontfamily\popsemi{PlusJakartaSans-SemiBold.ttf}[Path=../../../fonts/]
\newfontfamily\popxbold{PlusJakartaSans-ExtraBold.ttf}[Path=../../../fonts/]
\newfontfamily\popxboldls{PlusJakartaSans-ExtraBold.ttf}[Path=../../../fonts/,LetterSpace=3.0]

\setlist[enumerate]{leftmargin=1.7em,itemsep=5pt,topsep=4pt}
\setlist[itemize]{leftmargin=1.7em,itemsep=4pt,topsep=4pt,label={\color{poporange}\textbullet}}
\setlength{\parindent}{0pt}
\setlength{\parskip}{5pt}
\renewcommand{\arraystretch}{1.25}

% pgfplots : style Pop par défaut
\pgfplotsset{
  every axis/.append style={axis line style={popink,line width=1pt},tick style={popink},
    grid style={popline,line width=0.5pt},label style={font=\small},tick label style={font=\footnotesize}},
  popcurve/.style={poporange,line width=1.8pt,no markers,smooth},
}
% Même style disponible dans un \draw TikZ simple (hors pgfplots)
\tikzset{popcurve/.style={poporange,line width=1.8pt}}
\tikzset{popgrid/.style={popline,very thin}}
\colorlet{popcurve}{poporange} % le nom du style est parfois utilisé comme couleur

% ============================================================
% Pill / squiggle
% ============================================================
\newcommand{\poppill}[3]{%
  \tikz[baseline=-0.55ex]\node[draw=popink,line width=1.2pt,rounded corners=2.6pt,%
    fill=#2,inner xsep=6.5pt,inner ysep=3pt]{%
    \popxboldls\fontsize{7}{7}\selectfont\color{#3}\MakeUppercase{#1}};%
}
\newcommand{\popsquiggle}[1]{%
  \tikz[baseline=-0.4ex]{
    \draw[#1,line width=2.6pt,line cap=round]
      plot [smooth] coordinates {(0,0.09) (0.34,0.01) (0.68,0.21) (1.02,0.01) (1.36,0.10)};
  }%
}
% Mot-clé surligné (marqueur orange sous le texte)
\newcommand{\cle}[1]{{\popxbold\color{popdark}#1}}

% ============================================================
% En-tête / pied
% ============================================================
\pagestyle{fancy}
\fancyhf{}
\renewcommand{\headrulewidth}{0pt}
\renewcommand{\footrulewidth}{0pt}
\lhead{\raisebox{-0.15ex}{\poplogo\fontsize{13}{13}\selectfont\color{popink}OnBuch\color{poporange}+}}
\rhead{\poppill{\DOCMATIERE\ \textbullet\ \DOCNIVEAU\ \textbullet\ Leçon \DOCLECON}{white}{popink}}
\lfoot{\popxbold\fontsize{7.6}{7.6}\selectfont\color{popmuted}\MakeUppercase{Cours signé OnBuch+}\ \textbullet\ {\popsemi\DOCTITRE}}
\rfoot{\tikz[baseline=-0.6ex]\node[rounded corners=8.5pt,fill=popink,minimum height=17pt,minimum width=17pt,inner xsep=5pt,inner sep=0]{\popxbold\fontsize{8}{8}\selectfont\color{popcream}\thepage\,/\,\pageref*{LastPage}};}

% ============================================================
% Titres
% ============================================================
\newcommand{\chaptitre}[2]{%
  \begin{tcolorbox}[enhanced,colback=poporange,colframe=popink,boxrule=2.6pt,arc=11pt,
      drop shadow=popink,left=15pt,right=15pt,top=12pt,bottom=11pt,before skip=3pt,after skip=18pt]
    \poppill{#2}{popink}{popcream}\par\vspace{9pt}
    {\raggedright\hyphenpenalty=10000\exhyphenpenalty=10000\popdisplay\fontsize{22}{24}\selectfont\color{popcream}\MakeUppercase{#1}\par}
  \end{tcolorbox}
}
% Section numérotée : pastille orange avec le numéro + titre Archivo
\newcounter{popsec}
\newcommand{\coursec}[1]{%
  \stepcounter{popsec}\setcounter{popsub}{0}%
  \par\vspace{16pt}\noindent
  \begin{minipage}{\linewidth}
  \tikz[baseline=-0.7ex]\node[draw=popink,line width=1.6pt,fill=poporange,rounded corners=4pt,minimum size=22pt,inner sep=0]{\popdisplay\fontsize{12}{12}\selectfont\color{popcream}\thepopsec};%
  \hspace{8pt}{\popdisplay\fontsize{15}{17}\selectfont\color{popink}\MakeUppercase{#1}}\par
  \vspace{4pt}\noindent{\color{poporange}\rule{36pt}{3.6pt}}
  \end{minipage}\par\nopagebreak\vspace{8pt}
}
\newcounter{popsub}
\newcommand{\courssub}[1]{%
  \stepcounter{popsub}%
  \par\vspace{9pt}\noindent{\popxbold\fontsize{11.5}{13}\selectfont\color{popink}{\color{poporange}\thepopsec.\thepopsub}\hspace{6pt}#1}\par\nopagebreak\vspace{4pt}
}

% ============================================================
% Boîtes pédagogiques — même recette : fond blanc, bordure épaisse colorée,
% ombre dure encre, badge plein. Titre optionnel [#1].
% ============================================================
\tcbset{popbase/.style={breakable,enhanced,colback=white,boxrule=2.3pt,arc=9pt,
  shadow={3pt}{-3pt}{0pt}{popink},left=14pt,right=14pt,top=11pt,bottom=11pt,before skip=11pt,after skip=15pt,
  fontupper=\popsemi\fontsize{10.6}{14.5}\selectfont\color{popink},
  fontlower=\popsemi\fontsize{10.6}{14.5}\selectfont\color{popink}}}
% Coche dessinée (le glyphe ✓ n'existe pas dans les polices du document)
\newcommand{\popcheck}[1]{\tikz[baseline=-0.5ex]\draw[#1,line width=1.6pt,line cap=round,line join=round](-0.13,0.0)--(-0.04,-0.09)--(0.13,0.1);}
\providecommand{\coche}{\popcheck{popgreen}}
\providecommand{\resultat}[1]{\par\smallskip\noindent\fcolorbox{popgreen}{popgreenL}{\parbox{\dimexpr\linewidth-2\fboxsep-2\fboxrule\relax}{\textbf{Résultat.} #1}}\par\smallskip}
\newcommand{\popboxhead}[3]{\poppill{#1}{#2}{white}\ifx\relax#3\relax\else\hspace{6pt}{\popxbold\fontsize{10.6}{12}\selectfont\color{popink}#3}\fi\par\vspace{7pt}}

\newtcolorbox{aretenir}[1][]{popbase,colframe=popgreen,before upper={\popboxhead{À retenir}{popgreen}{#1}}}
% Texte en espagnol (document de lecture, dialogue, poème, extrait) : cadre
% d'étude. Un titre optionnel (source, auteur) s'affiche dans l'en-tête.
\newtcolorbox{textebox}[1][]{popbase,colframe=popblue,before upper={\popboxhead{Texto}{popblue}{#1}\begin{otherlanguage}{spanish}},after upper={\end{otherlanguage}}}
% Marques ✓ / ✗ (forme correcte / fautive) souvent utilisées par les modèles
\providecommand{\cross}{\textcolor{poppink}{\ensuremath{\times}}}
\providecommand{\xmark}{\cross}\providecommand{\crossmark}{\cross}\providecommand{\cmark}{\ensuremath{\checkmark}}
% Ligne de réponse / justification à compléter (inventée par les modèles)
\providecommand{\justification}[1]{\par\noindent\textit{Justification~:} \dotfill\par}
% \textipa (package tipa non chargé) : la police ne couvre pas l'API -> simple passage
\providecommand{\textipa}[1]{#1}
% Soulignement (ulem non chargé) : \uline, \ul
\providecommand{\uline}[1]{\underline{#1}}\providecommand{\ul}[1]{\underline{#1}}
% Ordinaux français écrits en macro par les modèles : 1\ière, 3\ième
\newcommand{\ière}{\textsuperscript{re}}\newcommand{\ième}{\textsuperscript{e}}
% Mot ou phrase en espagnol à l'intérieur d'un paragraphe français
\newcommand{\esp}[1]{\ifmmode\text{\textit{#1}}\else\begingroup\selectlanguage{spanish}\itshape #1\endgroup\fi}
\newtcolorbox{exemplebox}[1][]{popbase,colframe=poporange,before upper={\popboxhead{Exemple}{poporange}{#1}}}
\newtcolorbox{definition}[1][]{popbase,colframe=popblue,before upper={\popboxhead{Définition}{popblue}{#1}}}
\newtcolorbox{propriete}[1][]{popbase,colframe=poppurple,before upper={\popboxhead{Propriété}{poppurple}{#1}}}
\newtcolorbox{methode}[1][]{popbase,colframe=popgold,before upper={\popboxhead{Méthode}{popgold}{#1}}}
\newtcolorbox{attention}[1][]{popbase,colframe=poppink,before upper={\popboxhead{Attention}{poppink}{#1}}}
\newtcolorbox{experience}[1][]{popbase,colframe=popblue,colback=popblueL,before upper={\popboxhead{Expérience}{popblue}{#1}}}
% « Le savais-tu ? » : carte violette pleine (contexte, culture, Cameroun)
\newtcolorbox{savaistu}[1][]{popbase,colback=poppurple,colframe=popink,
  fontupper=\popsemi\fontsize{10.4}{14.2}\selectfont\color{white},
  before upper={\poppill{Le savais-tu\,?}{white}{poppurple}\ifx\relax#1\relax\else\hspace{6pt}{\popxbold\color{white}#1}\fi\par\vspace{7pt}}}
% Exercice résolu : énoncé en haut, \tcblower puis la solution sur fond crème
\newcounter{popexo}\newcounter{popexr}
\newtcolorbox{exoresolu}[1][]{popbase,colframe=poporange,colbacklower=popcream,
  segmentation style={popink,line width=1.2pt,dashed},
  before upper={\stepcounter{popexr}\popboxhead{Exercice résolu \thepopexr}{poporange}{#1}},
  before lower={\poppill{Solution}{popink}{popcream}\par\vspace{6pt}}}
% Exercice (non résolu en place) — la correction est dans la partie Corrigés
\newtcolorbox{exercice}[1][]{popbase,colframe=popink,
  before upper={\stepcounter{popexo}\popboxhead{Exercice \thepopexo}{popink}{#1}}}
% Corrigé
\newtcolorbox{corrige}[1][]{popbase,colframe=popgreen,colback=popgreenL,
  before upper={\popboxhead{Corrigé}{popgreen}{#1}}}
% Objectifs / prérequis / bilan
\newtcolorbox{objectifs}{popbase,colframe=popink,colback=poporangeL,
  before upper={\popboxhead{Objectifs de la leçon}{popink}{}}}
\newtcolorbox{prerequis}{popbase,colframe=popink,colback=popblueL,
  before upper={\popboxhead{Prérequis}{popblue}{}}}
\newtcolorbox{bilan}{popbase,colframe=popink,colback=white,boxrule=2.8pt,
  before upper={\popboxhead{Fiche bilan}{poporange}{L'essentiel à connaître pour l'examen}}}
% Figure encadrée avec légende
\newtcolorbox{popfigure}[1][]{enhanced,breakable=false,colback=white,colframe=popline,boxrule=1.4pt,arc=8pt,
  left=10pt,right=10pt,top=10pt,bottom=8pt,before skip=10pt,after skip=12pt,halign=center,
  lower separated=false,halign lower=center,
  fontlower=\popsemi\fontsize{9}{11.5}\selectfont\color{popmuted},
  #1}
\newcommand{\legende}[1]{\par\vspace{6pt}{\popsemi\fontsize{9}{11.5}\selectfont\color{popmuted}#1}}

% ============================================================
% Signature OnBuch+
% ============================================================
\newcommand{\poplogoinline}[1]{{\poplogo\fontsize{#1}{#1}\selectfont\color{popink}OnBuch\color{poporange}+}}
\newcommand{\signatureonbuch}{%
  \begin{tcolorbox}[enhanced,colback=popink,colframe=popink,boxrule=2.2pt,arc=10pt,
      drop shadow=poporange,left=14pt,right=14pt,top=10pt,bottom=10pt,before skip=0pt,after skip=0pt]
    \tikz[baseline=-0.7ex]\node[circle,fill=popgreen,draw=popcream,line width=1.3pt,minimum size=22pt,inner sep=0]{\popcheck{white}};%
    \hspace{9pt}\begin{minipage}[c]{0.62\linewidth}
      {\popxbold\fontsize{12}{13}\selectfont\color{popcream}Signé\ \textbullet\ {\poplogo OnBuch\color{poporange}+}}\par\vspace{2pt}
      {\raggedright\popsemi\fontsize{8}{10.5}\selectfont\color{popline}\MakeUppercase{Équipe pédagogique OnBuch+}\\\MakeUppercase{Conforme au programme officiel MINESEC}\par}
    \end{minipage}\hfill
    {\poplogo\fontsize{22}{22}\selectfont\color{popcream}OnBuch\color{poporange}+}
  \end{tcolorbox}%
}

% ============================================================
% Page de garde
% #1 titre · #2 matière · #3 niveau · #4 module · #5 n° leçon · #6 durée ·
% #7 description / objectifs (texte ou itemize)
% ============================================================
\newcommand{\pagegarde}[7]{%
  \thispagestyle{empty}
  \newgeometry{a4paper,margin=1.7cm,top=1.6cm,bottom=1.4cm}%
  \begin{tikzpicture}[remember picture,overlay]
    \fill[poporange!18] ([xshift=-3.2cm,yshift=-3.6cm]current page.north east) circle (4.6cm);
    \fill[poppurple!14] ([xshift=2.4cm,yshift=7.6cm]current page.south west) circle (3.8cm);
  \end{tikzpicture}%
  {\poplogo\fontsize{30}{30}\selectfont\color{popink}OnBuch\color{poporange}+}\hfill
  \raisebox{0.6ex}{\poppill{Cours signé \textbullet\ Premium}{popink}{popcream}}\par
  \vspace{1pt}\popsquiggle{poppurple}\par
  \vspace{8pt}
  \begin{tcolorbox}[enhanced,colback=poporange,colframe=popink,boxrule=2.8pt,arc=13pt,
      drop shadow=popink,left=18pt,right=18pt,top=12pt,bottom=13pt,after skip=10pt]
    \poppill{#2 \textbullet\ #3}{popink}{popcream}\hspace{5pt}\poppill{#4}{white}{popink}\par\vspace{8pt}
    {\popdisplay\fontsize{14}{14}\selectfont\color{popink}LEÇON #5}\par\vspace{4pt}
    {\raggedright\hyphenpenalty=10000\exhyphenpenalty=10000\popdisplay\fontsize{25}{27}\selectfont\color{popcream}\MakeUppercase{#1}\par}
    \vspace{8pt}
    {\popxbold\fontsize{9.5}{9.5}\selectfont\color{popink}\MakeUppercase{Durée conseillée : #6}}
  \end{tcolorbox}
  \begin{tcolorbox}[enhanced,colback=white,colframe=popink,boxrule=2.2pt,arc=10pt,
      drop shadow=popink,left=16pt,right=16pt,top=10pt,bottom=10pt,after skip=8pt]
    \poppill{Dans cette leçon}{poppurple}{white}\par\vspace{5pt}
    \popsemi\fontsize{9.3}{12.6}\selectfont\color{popink}#7
  \end{tcolorbox}
  \noindent\begin{minipage}[t]{0.487\linewidth}
    \begin{tcolorbox}[enhanced,colback=popgreen,colframe=popink,boxrule=2.2pt,arc=10pt,
        drop shadow=popink,left=11pt,right=11pt,top=10pt,bottom=10pt,height=3.1cm,valign=center]
      \tcbox[colback=white,colframe=popink,boxrule=1.3pt,arc=5pt,boxsep=0pt,nobeforeafter,
        left=3pt,right=3pt,top=3pt,bottom=3pt,valign=center]{\qrcode[height=1.9cm]{https://play.google.com/store/apps/details?id=com.luvvix.onbuch&hl=fr}}%
      \hspace{8pt}\begin{minipage}[c]{0.5\linewidth}
        {\popdisplay\fontsize{11}{12.5}\selectfont\color{white}\MakeUppercase{Télécharge OnBuch}}\par\vspace{4pt}
        {\raggedright\popsemi\fontsize{8.4}{11}\selectfont\color{white}Gratuit \textbullet\ Google Play\\Tuteur IA, annales, concours, résultats\par}
      \end{minipage}
    \end{tcolorbox}
  \end{minipage}\hfill
  \begin{minipage}[t]{0.487\linewidth}
    \begin{tcolorbox}[enhanced,colback=poppurple,colframe=popink,boxrule=2.2pt,arc=10pt,
        drop shadow=popink,left=11pt,right=11pt,top=10pt,bottom=10pt,height=3.1cm,valign=center]
      \tikz[baseline=-0.7ex]\node[circle,fill=white,draw=popink,line width=1.3pt,minimum size=1.6cm,inner sep=0]{\popdisplay\fontsize{24}{24}\selectfont\color{poppurple}+};%
      \hspace{8pt}\begin{minipage}[c]{0.6\linewidth}
        {\popdisplay\fontsize{11}{12.5}\selectfont\color{white}\MakeUppercase{Passe à OnBuch+}}\par\vspace{4pt}
        {\raggedright\popsemi\fontsize{8.4}{11}\selectfont\color{white}Tous les cours signés, Tuteur IA illimité, fiches PDF — onglet OnBuch+ de l'app\par}
      \end{minipage}
    \end{tcolorbox}
  \end{minipage}
  \vspace{8pt}
  \popcontenu
  \vfill
  \signatureonbuch
  \restoregeometry
  \newpage
}

% Rangée « Ce cours contient » (page de garde)
\newcommand{\poptile}[3]{% #1 couleur #2 gros texte #3 légende
  \begin{tcolorbox}[enhanced,colback=#1,colframe=popink,boxrule=2pt,arc=9pt,shadow={2.5pt}{-2.5pt}{0pt}{popink},
      left=6pt,right=6pt,top=9pt,bottom=9pt,halign=center,valign=center,height=1.75cm,width=\linewidth,nobeforeafter]
    {\popdisplay\fontsize{12.5}{13}\selectfont\color{white}\MakeUppercase{#2}}\par\vspace{3pt}
    {\centering\popsemi\fontsize{7.8}{9.5}\selectfont\color{white}#3\par}
  \end{tcolorbox}}
\newcommand{\popcontenu}{%
  \par\noindent{\popxboldls\fontsize{7.5}{7.5}\selectfont\color{popmuted}CE COURS SIGNÉ CONTIENT}\par\vspace{7pt}
  \noindent\begin{minipage}[t]{0.235\linewidth}\poptile{popblue}{Cours}{complet, illustré, pas à pas}\end{minipage}\hfill
  \begin{minipage}[t]{0.235\linewidth}\poptile{poporange}{Méthodes}{et exercices résolus}\end{minipage}\hfill
  \begin{minipage}[t]{0.235\linewidth}\poptile{poppink}{Exercices}{type examen + corrigés}\end{minipage}\hfill
  \begin{minipage}[t]{0.235\linewidth}\poptile{popgreen}{Fiche bilan}{l'essentiel à retenir}\end{minipage}\par}

% Sommaire
\newtcolorbox{sommaire}{enhanced,colback=white,colframe=popink,boxrule=2.2pt,arc=10pt,
  drop shadow=popink,left=18pt,right=18pt,top=13pt,bottom=13pt,before skip=6pt,after skip=20pt,
  before upper={\poppill{Sommaire}{poppurple}{white}\par\vspace{9pt}\popsemi\fontsize{10.6}{14.5}\selectfont\color{popink}}}

% Signature de fin de cours — poussée tout en bas de la dernière page
\newcommand{\findecours}{%
  \par\vfill\nopagebreak
  \begin{center}\popsquiggle{poporange}\end{center}\vspace{4pt}
  \signatureonbuch
}
\AtBeginDocument{\def\llbracket{\mathopen{[\mkern-3.5mu[}}\def\rrbracket{\mathclose{]\mkern-3.5mu]}}} % intervalles d'entiers (glyphe ⟦ absent de la police)
% Glyphes absents de Fira Math : redéfinis à partir de glyphes disponibles
\AtBeginDocument{%
  \def\setminus{\mathbin{\backslash}}\let\smallsetminus\setminus
  \def\vdots{\mathord{\vbox{\baselineskip3.2pt\lineskiplimit0pt\kern4pt\hbox{.}\hbox{.}\hbox{.}}}}%
  \def\ddots{\mathinner{\mkern1mu\raise6.5pt\hbox{.}\mkern2mu\raise3.6pt\hbox{.}\mkern2mu\raise0.7pt\hbox{.}\mkern1mu}}%
  \def\triangle{\mathord{\scalebox{0.9}{$\symup{\Delta}$}}}%
  \def\nabla{\mathord{\raisebox{\depth}{\scalebox{1}[-1]{$\symup{\Delta}$}}}}%
  \def\checkmark{\popcheck{popdark}}%
  \def\star{\mathbin{\ast}}\def\bigstar{\mathord{\ast}}%
  \def\diamond{\mathbin{\scalebox{0.6}{$\Diamond$}}}%
  \def\bigcup{\mathop{\mathchoice{\raisebox{-0.3em}{\scalebox{1.7}{$\cup$}}}{\scalebox{1.25}{$\cup$}}{\cup}{\cup}}\displaylimits}%
  \def\bigcap{\mathop{\mathchoice{\raisebox{-0.3em}{\scalebox{1.7}{$\cap$}}}{\scalebox{1.25}{$\cap$}}{\cap}{\cap}}\displaylimits}%
  \def\lesseqqgtr{\mathrel{\lessgtr}}\def\bigoplus{\mathop{\oplus}\displaylimits}%
}
\providecommand{\ssi}{si et seulement si} % abréviation fréquente
\AtBeginDocument{\providecommand{\wideparen}{\overparen}} % arc de cercle

%% FIGFIX-BEGIN v3 — correctifs globaux des figures (docs/onbuch-generation/tools/figfix)
\usepackage{adjustbox}\usepackage{etoolbox}
% toute figure TikZ/pgfplots plus large que la ligne est réduite à la largeur disponible
\BeforeBeginEnvironment{tikzpicture}{\begin{adjustbox}{max width=\linewidth}}
\AfterEndEnvironment{tikzpicture}{\end{adjustbox}}
% légendes pgfplots lisibles par-dessus les courbes
\pgfplotsset{every axis/.append style={legend style={fill=white,fill opacity=0.92,text opacity=1,draw=black!15,inner sep=3pt}}}
% étiquettes d'arêtes (syntaxe edge["..."]) avec fond blanc
\tikzset{every edge quotes/.append style={fill=white,inner sep=1.2pt}}
% bulles de cartes mentales : texte à la ligne dans la bulle, bulle un peu plus grande
\tikzset{every concept/.append style={text width=2.0cm,align=center,minimum size=2.3cm,inner sep=2pt}}
%% FIGFIX-END
\begin{document}

\pagegarde{El deporte, el ocio y el bienestar}{Espagnol}{1ère A}{Módulo 2 : Environnement, santé et bien-être}{EA07}{2 h}{Cette leçon mixte du module « Environnement, santé et bien-être » permet d'acquérir le lexique du sport et des loisirs, de lire et commenter un texte sur le bien-être, et de maîtriser l'expression du goût (gustar) et de la fréquence (soler + infinitif) pour parler de ses habitudes sportives.
\vspace{4pt}
\begin{multicols}{2}\small
\begin{itemize}
\item À la fin de cette leçon, je sais lire et comprendre un texte sur le sport, les loisirs et le bien-être
\item À la fin de cette leçon, je sais utiliser le lexique du sport et du ocio (deporte, entrenamiento, partido, aficionado, campeón, equipo...)
\item À la fin de cette leçon, je sais exprimer mes goûts avec me gusta / me encanta + nom ou infinitif
\item À la fin de cette leçon, je sais exprimer la fréquence et l'habitude avec soler + infinitif et les adverbes de fréquence
\item À la fin de cette leçon, je sais conjuguer soler et practicar au présent de l'indicatif
\item À la fin de cette leçon, je sais traduire des phrases du français vers l'espagnol et inversement sur le thème du sport
\end{itemize}\end{multicols}}

\begin{sommaire}
\begin{multicols}{2}
\begin{enumerate}
\item Texte support : El deporte, una fuente de bienestar
\item Lexique thématique : el deporte y el ocio
\item Grammaire 1 : l'expression du goût — gustar
\item Grammaire 2 : l'habitude et la fréquence — soler + infinitif
\item Traduction : thème et version sur le sport
\item Commentaire de texte et production guidée
\item Activité d'intégration
\item Méthodes et exercices résolus
\item Exercices
\item Corrigés des exercices
\item Fiche bilan
\end{enumerate}
\end{multicols}
\end{sommaire}

\begin{objectifs}
\begin{itemize}
\item À la fin de cette leçon, je sais lire et comprendre un texte sur le sport, les loisirs et le bien-être
\item À la fin de cette leçon, je sais utiliser le lexique du sport et du ocio (deporte, entrenamiento, partido, aficionado, campeón, equipo...)
\item À la fin de cette leçon, je sais exprimer mes goûts avec me gusta / me encanta + nom ou infinitif
\item À la fin de cette leçon, je sais exprimer la fréquence et l'habitude avec soler + infinitif et les adverbes de fréquence
\item À la fin de cette leçon, je sais conjuguer soler et practicar au présent de l'indicatif
\item À la fin de cette leçon, je sais traduire des phrases du français vers l'espagnol et inversement sur le thème du sport
\item À la fin de cette leçon, je sais parler de mes loisirs et comparer le sport au Cameroun et dans les pays hispanophones
\item À la fin de cette leçon, je sais rédiger un court texte pour encourager la pratique sportive
\end{itemize}
\end{objectifs}

\begin{prerequis}
\begin{itemize}
\item Présent de l'indicatif des verbes réguliers en -ar, -er, -ir (Seconde)
\item Les pronoms personnels sujets et les pronoms toniques (a mí, a ti...)
\item Le lexique de la santé et du corps humain (leçons EA05-EA06)
\item Les verbes à changement orthographique o > ue (poder, dormir) vus en Seconde
\item Les adverbes de quantité et de temps simples (muy, mucho, siempre)
\end{itemize}
\end{prerequis}

\newpage

\coursec{Texte support : El deporte, una fuente de bienestar}

C'est samedi matin à Yaoundé : au stade Ahmadou Ahidjo, des jeunes s'entraînent à l'athlétisme pendant que d'autres jouent au football dans les quartiers. Le soir, toute la famille se réunit devant le match des Lions Indomptables. Et toi, quel sport pratiques-tu ? Aimes-tu le sport ou préfères-tu d'autres loisirs ? En Espagne et en Amérique latine aussi, le sport et le \esp{ocio} (les loisirs) occupent une place centrale dans la vie quotidienne.

\textbf{Problématique :} comment parler de nos goûts sportifs et évaluer l'importance du sport pour la santé physique et mentale ? Pour répondre à cette question, nous allons lire le récit de Samuel, un jeune Camerounais sportif, et le comparer à celui de son correspondant espagnol, Pablo.

\textbf{Objectifs de la leçon :} à la fin de cette séquence, tu seras capable de lire et comprendre un texte sur le sport, de parler de tes loisirs et de tes goûts avec \esp{me gusta}, d'exprimer tes habitudes avec \esp{soler} + infinitif, et d'encourager un camarade à pratiquer un sport.

\courssub{Lecture globale du texte}

\begin{methode}[Comment lire un texte en espagnol ?]
\begin{enumerate}
\item \textbf{Lire le titre et anticiper :} \esp{El deporte, una fuente de bienestar}. Que va-t-on trouver dans ce texte ? Le titre annonce les bienfaits du sport.
\item \textbf{Première lecture globale, en silence, sans dictionnaire :} chercher seulement l'idée générale (de quoi parle le texte ? qui parle ?).
\item \textbf{Repérer les mots transparents :} ce sont les mots qui ressemblent au français : \esp{deporte}, \esp{atletismo}, \esp{mental}, \esp{física}, \esp{problemas}.
\item \textbf{Deuxième lecture avec le glossaire :} cette fois, on cherche les détails (qui, quoi, quand, combien de fois, pourquoi).
\item \textbf{Répondre aux questions en citant le texte :} chaque réponse doit être justifiée par une phrase ou un groupe de mots du texte.
\end{enumerate}
\end{methode}

\textbf{Consigne de classe :} première lecture silencieuse (3 minutes), puis lecture à voix haute par l'enseignant. Pendant la lecture à voix haute, souligne dans ton cahier les mots que tu connais déjà.

\begin{textebox}[El deporte, una fuente de bienestar]
Me llamo Samuel y vivo en Yaoundé. Suelo practicar deporte cuatro veces por semana : los martes y los jueves tengo entrenamiento de atletismo, y el sábado juego un partido de fútbol con mi equipo del barrio.

Me gusta mucho el deporte porque es bueno para la salud física y mental. Cuando corro, me siento libre y olvido los problemas. Después del esfuerzo, duermo mejor y estoy más tranquilo en clase.

Mi amigo español Pablo es un gran aficionado al fútbol : nunca se pierde un partido del Real Madrid. Él suele montar en bicicleta los domingos al aire libre, con su padre. Pablo quiere ser campeón de ciclismo un día.

Los dos pensamos que el deporte es una fuente de bienestar : nos hace más fuertes, más alegres y nos permite conocer a nuevos amigos. El deporte también une a las personas : cuando juegan los Leones Indomables, ¡todo el país está contento! ¡Anímate a practicar deporte!
\end{textebox}

\begin{definition}[Glossaire]
\begin{itemize}
\item \esp{fuente} : source (féminin : \esp{la fuente})
\item \esp{salud} : santé (féminin : \esp{la salud})
\item \esp{esfuerzo} : effort (masculin : \esp{el esfuerzo})
\item \esp{animarse} : se motiver, s'encourager (verbe pronominal ; à l'impératif : \esp{¡anímate!} = motive-toi !)
\item \esp{al aire libre} : en plein air
\item \esp{barrio} : quartier (masculin : \esp{el barrio})
\item \esp{corro} : je cours (verbe \esp{correr}, 1\textsuperscript{re} personne du singulier du présent)
\item \esp{une} : unit, rassemble (verbe \esp{unir}, 3\textsuperscript{e} personne du singulier)
\item \esp{perderse} : rater, manquer (un spectacle, un match) ; aussi : se perdre, s'égarer
\end{itemize}
\end{definition}

\textbf{Traduction du texte :} « Je m'appelle Samuel et je vis à Yaoundé. J'ai l'habitude de faire du sport quatre fois par semaine : le mardi et le jeudi, j'ai un entraînement d'athlétisme, et le samedi je joue un match de football avec mon équipe du quartier. J'aime beaucoup le sport parce qu'il est bon pour la santé physique et mentale. Quand je cours, je me sens libre et j'oublie les problèmes. Après l'effort, je dors mieux et je suis plus calme en classe. Mon ami espagnol Pablo est un grand passionné de football : il ne rate jamais un match du Real Madrid. Lui a l'habitude de faire du vélo le dimanche en plein air, avec son père. Pablo veut être champion de cyclisme un jour. Nous pensons tous les deux que le sport est une source de bien-être : il nous rend plus forts, plus joyeux et il nous permet de rencontrer de nouveaux amis. Le sport rassemble aussi les gens : quand les Lions Indomptables jouent, tout le pays est content ! Motive-toi à faire du sport ! »

\courssub{Repérage du lexique officiel dans le texte}

Le texte contient les neuf mots du lexique officiel de la leçon. Retrouve-les et mémorise-les grâce au tableau suivant :

\begin{center}
\begin{tabularx}{\linewidth}{|X|X|X|}
\hline
\rowcolor{popblueL} Español & Français & Phrase du texte \\
\hline
el deporte & le sport & \esp{Suelo practicar deporte...} \\
\hline
el fútbol & le football & \esp{...juego un partido de fútbol...} \\
\hline
el atletismo & l'athlétisme & \esp{...entrenamiento de atletismo...} \\
\hline
el entrenamiento & l'entraînement & \esp{...tengo entrenamiento...} \\
\hline
el partido & le match & \esp{...nunca se pierde un partido...} \\
\hline
el ocio & les loisirs & (titre de la leçon : \esp{el deporte, el ocio y el bienestar}) \\
\hline
el aficionado & le passionné, le supporter & \esp{...un gran aficionado al fútbol...} \\
\hline
el campeón & le champion & \esp{...quiere ser campeón de ciclismo...} \\
\hline
el equipo & l'équipe & \esp{...con mi equipo del barrio...} \\
\hline
\end{tabularx}
\end{center}

\begin{attention}[Faux amis et pièges de vocabulaire]
\begin{itemize}
\item \esp{el partido} signifie « le match » dans le contexte sportif (il peut aussi signifier « le parti politique » dans un contexte politique). \esp{jugar un partido} = jouer un match.
\item Ne pas confondre \esp{perder} et \esp{perderse} : \esp{perder un partido} = perdre un match ; \esp{perderse un partido} = rater, manquer un match ; \esp{perderse} (seul) = se perdre, s'égarer. Dans le texte : \esp{Pablo nunca se pierde un partido} (Pablo ne rate jamais un match) ; mais \esp{Mi equipo perdió el partido} (mon équipe a perdu le match) et \esp{Me perdí en el mercado} (je me suis perdu au marché).
\item Attention au verbe \esp{jugar} : il se conjugue avec le changement \textbf{u > ue} : \esp{yo juego}, \esp{tú juegas}, \esp{él juega}, mais \esp{nosotros jugamos}, \esp{vosotros jugáis}, \esp{ellos juegan}. Dans le texte : \esp{el sábado juego un partido} (le samedi je joue un match). On dit \esp{jugar al fútbol} (avec la préposition \esp{a} + \esp{el} = \esp{al}) ou \esp{jugar un partido de fútbol}.
\item \esp{el aficionado} est un passionné qui regarde ; celui qui pratique est \esp{el deportista} (le sportif). Féminins : \esp{la aficionada}, \esp{la deportista}.
\item \esp{el campeón} a pour féminin \esp{la campeona} : \esp{Pablo quiere ser campeón} / \esp{María quiere ser campeona}.
\end{itemize}
\end{attention}

\courssub{Les idées principales du texte}

Après la deuxième lecture, on peut organiser les idées du texte autour de trois grands bienfaits du sport. Observe la carte mentale suivante :

\begin{popfigure}
\begin{tikzpicture}[mindmap, grow cyclic, every node/.style={concept, text=white, align=center}, root concept/.append style={fill=popblue, font=\bfseries}, level 1/.append style={level distance=3.2cm, sibling angle=120}]
\node[root concept]{El deporte}
  child[concept color=popgreen]{ node[align=center]{salud\\física} }
  child[concept color=poporange]{ node[align=center]{salud\\mental} }
  child[concept color=poppurple]{ node[align=center]{vida\\social} };
\end{tikzpicture}
\legende{Carte mentale : les trois bienfaits du sport dans le texte}
\end{popfigure}

\begin{aretenir}[Les trois idées principales]
\begin{enumerate}
\item \textbf{Les bienfaits physiques} (\esp{salud física}) : \esp{es bueno para la salud física} ; \esp{nos hace más fuertes} ; \esp{duermo mejor}. Le sport rend plus fort et améliore le sommeil.
\item \textbf{Les bienfaits mentaux} (\esp{salud mental}) : \esp{olvido los problemas} ; \esp{estoy más tranquilo en clase} ; \esp{nos hace más alegres}. Le sport libère l'esprit et apaise.
\item \textbf{La dimension sociale} (\esp{vida social}) : \esp{nos permite conocer a nuevos amigos} ; \esp{el deporte une a las personas}. Le sport crée des amitiés et rassemble tout un pays derrière son équipe.
\end{enumerate}
\end{aretenir}

\begin{attention}[Ser ou estar dans le texte ?]
Le texte illustre bien la différence entre \esp{ser} et \esp{estar} :
\begin{itemize}
\item \esp{El deporte \textbf{es} bueno para la salud} : caractéristique générale, permanente $\rightarrow$ \esp{ser}.
\item \esp{\textbf{Estoy} más tranquilo en clase} : état temporaire, résultat d'un changement $\rightarrow$ \esp{estar}.
\item \esp{Todo el país \textbf{está} contento} : état passager (le jour du match) $\rightarrow$ \esp{estar}.
\end{itemize}
Retiens : \esp{ser} = ce qu'on est de façon durable ; \esp{estar} = comment on se trouve à un moment donné.
\end{attention}

\courssub{Questions de compréhension}

\begin{exercice}[Compréhension de l'écrit]
Réponds en espagnol, par des phrases complètes, en citant le texte.
\begin{enumerate}
\item ¿Qué deportes practica Samuel ?
\item ¿Cuántas veces por semana entrena ?
\item ¿Qué beneficios del deporte menciona el texto ?
\item ¿Qué deporte prefiere Pablo ?
\item ¿Por qué el deporte es importante para la salud mental ?
\end{enumerate}
\end{exercice}

\begin{corrige}[Questions de compréhension]
\begin{enumerate}
\item \esp{Samuel practica el atletismo y el fútbol.} (Samuel pratique l'athlétisme et le football.) Citation : \esp{« tengo entrenamiento de atletismo, y el sábado juego un partido de fútbol »}.
\item \esp{Entrena cuatro veces por semana.} (Il s'entraîne quatre fois par semaine.) Citation : \esp{« Suelo practicar deporte cuatro veces por semana »}.
\item \esp{El texto menciona beneficios físicos (más fuertes, duerme mejor), mentales (olvida los problemas, está más tranquilo) y sociales (conocer nuevos amigos).} (Le texte mentionne des bienfaits physiques, mentaux et sociaux.)
\item \esp{Pablo es aficionado al fútbol y también practica el ciclismo : suele montar en bicicleta los domingos.} (Pablo est passionné de football et pratique aussi le cyclisme.) Citation : \esp{« es un gran aficionado al fútbol »}.
\item \esp{El deporte es importante para la salud mental porque libera el estrés : « Cuando corro, me siento libre y olvido los problemas ».} Le sport permet d'oublier les problèmes et d'être plus calme.
\end{enumerate}
\end{corrige}

\begin{exoresolu}[Vrai ou faux ? Justifie par une citation du texte]
\textbf{Énoncé.} Dis si les phrases suivantes sont vraies (\esp{verdadero}) ou fausses (\esp{falso}) et justifie ta réponse par une citation du texte.
\begin{enumerate}
\item Samuel vive en Douala.
\item Pablo nunca ve los partidos del Real Madrid.
\item El deporte permite conocer a nuevos amigos.
\end{enumerate}
\tcblower
\textbf{Solution détaillée.}
\begin{enumerate}
\item \textbf{Falso.} Le texte dit : \esp{« Me llamo Samuel y vivo en Yaoundé »}. Samuel habite à Yaoundé, pas à Douala.
\item \textbf{Falso.} Le texte dit : \esp{« nunca se pierde un partido del Real Madrid »}. Pablo ne rate jamais un match du Real Madrid : il les regarde tous.
\item \textbf{Verdadero.} Le texte dit : \esp{« nos permite conocer a nuevos amigos »}. C'est la dimension sociale du sport.
\end{enumerate}
\textbf{Méthode :} pour justifier, recopie toujours entre guillemets le groupe de mots exact du texte qui prouve ta réponse. Une réponse sans citation n'est pas complète.
\end{exoresolu}

\begin{experience}[À toi de parler !]
Par binômes, posez-vous les questions suivantes et répondez avec des phrases complètes :
\begin{itemize}
\item \esp{¿Qué deporte practicas ?} (Quel sport pratiques-tu ?)
\item \esp{¿Cuántas veces por semana entrenas ?} (Combien de fois par semaine t'entraînes-tu ?)
\item \esp{¿Eres aficionado al fútbol ? ¿Cuál es tu equipo favorito ?} (Es-tu passionné de football ? Quelle est ton équipe préférée ?)
\end{itemize}
Modèle de réponse : \esp{Practico el baloncesto dos veces por semana. Soy aficionado al fútbol y mi equipo favorito son los Leones Indomables.} (Je pratique le basketball deux fois par semaine. Je suis passionné de football et mon équipe préférée, ce sont les Lions Indomptables.)
Remarque l'accord : \esp{mi equipo favorito son los Leones Indomables} est correct car le verbe \esp{ser} s'accorde ici avec le nom pluriel qui le suit ; on peut aussi dire simplement \esp{mi equipo favorito es el Canon de Yaoundé}.
\end{experience}

\begin{savaistu}[Le sport au Cameroun et en Espagne]
Le Cameroun a remporté la médaille d'or olympique de football aux Jeux de Sydney en 2000 : une victoire historique des Lions Indomptables qui a fait vibrer tout le pays, de Yaoundé à Douala. De son côté, l'Espagne a gagné la Coupe du monde en 2010 en Afrique du Sud, puis le Championnat d'Europe en 2012 et en 2024. En espagnol, on dit \esp{ganar un partido} (gagner un match) et \esp{ganar una medalla de oro} (gagner une médaille d'or). En espagnol, les Lions Indomptables s'appellent \esp{los Leones Indomables}. Et n'oublie pas la Guinée équatoriale, pays hispanophone voisin du Cameroun, où le football est aussi le sport roi : son équipe nationale féminine, \esp{el Nzalang Femenino}, est bien connue en Afrique !
\end{savaistu}

\begin{aretenir}[Ce qu'il faut retenir de cette lecture]
\begin{itemize}
\item Le texte de Samuel montre que le sport est \cle{una fuente de bienestar} : il apporte des bienfaits \cle{físicos}, \cle{mentales} et \cle{sociales}.
\item Neuf mots de lexique officiel à connaître par cœur : \esp{deporte, fútbol, atletismo, entrenamiento, partido, ocio, aficionado, campeón, equipo}.
\item Pour lire un texte : anticiper avec le titre, lire d'abord sans dictionnaire, repérer les mots transparents, puis répondre en citant le texte.
\item Deux expressions clés déjà rencontrées et que nous étudierons en grammaire : \esp{me gusta el deporte} (j'aime le sport) et \esp{suelo practicar} (j'ai l'habitude de pratiquer).
\item Distinction utile : \esp{perder un partido} (perdre un match) $\neq$ \esp{perderse un partido} (rater un match).
\end{itemize}
\end{aretenir}

\coursec{Lexique thématique : el deporte y el ocio}

Dans la section précédente, nous avons lu le texte \emph{El deporte, una fuente de bienestar} et nous avons découvert que le sport est essentiel pour la santé physique et mentale. Ce texte contenait de nombreux mots nouveaux : \esp{entrenamiento}, \esp{equipo}, \esp{aficionados}, \esp{estar en forma}... Pour bien comprendre un texte sur le sport et pouvoir parler de nos loisirs, il faut maintenant \textbf{organiser et mémoriser ce vocabulaire}. C'est l'objectif de cette section : construire, champ par champ, un lexique solide que tu réutiliseras dans les sections de grammaire, de traduction et d'expression orale.

\courssub{Observation : le vocabulaire dans son contexte}

Avant d'apprendre des listes de mots, observons-les dans une situation réelle. Lis ce court reportage original sur un match au Cameroun.

\begin{textebox}[Un sábado de fútbol en Yaoundé]
El sábado por la tarde, el estadio de Yaoundé está lleno de aficionados. Hoy juega el equipo del barrio contra el campeón del campeonato escolar. Los jugadores llegan temprano para el entrenamiento: corren, estiran y practican tiros al arco. El entrenador da las últimas instrucciones antes del partido.

A las cuatro, el árbitro pita el comienzo. El público grita y anima a su equipo. En el segundo tiempo, el delantero del barrio marca un gol y el estadio explota de alegría. Al final, el partido termina dos a uno: ¡el equipo local gana!

Después del partido, muchos jóvenes hablan de su tiempo libre. Unos prefieren jugar al fútbol, otros practican atletismo o hacen natación. Todos están de acuerdo: el deporte es bueno para la salud física y mental. Estar en forma y sentirse fuerte: eso es el verdadero bienestar.
\end{textebox}

\textbf{Glossaire :} \esp{lleno} = plein ; \esp{estiran} = ils s'étirent ; \esp{tiros al arco} = tirs au but ; \esp{pita} = siffle ; \esp{anima} = encourage ; \esp{delantero} = attaquant ; \esp{explotar de alegría} = exploser de joie ; \esp{estar de acuerdo} = être d'accord.

\textbf{Questions de compréhension :}
\begin{enumerate}
\item ¿Quién juega hoy en el estadio de Yaoundé?
\item ¿Qué hacen los jugadores antes del partido?
\item ¿Quién marca un gol en el segundo tiempo?
\item ¿Qué deportes practican los jóvenes en su tiempo libre?
\end{enumerate}

Tu remarques que les mots se regroupent naturellement : les \cle{sports}, la \cle{pratique sportive}, les \cle{lieux et personnes}, les \cle{loisirs} et le \cle{bien-être}. Étudions chaque champ.

\courssub{Champ 1 — Les sports}

\begin{center}
\begin{tabularx}{\linewidth}{|X|X|X|}\hline
\rowcolor{popblueL} \textbf{Español} & \textbf{Français} & \textbf{Remarque} \\ \hline
el deporte & le sport & masculin : \esp{el deporte es bueno} \\ \hline
el fútbol & le football & accent sur la \emph{u} \\ \hline
el atletismo & l'athlétisme & \esp{practicar atletismo} \\ \hline
el baloncesto & le basketball & aussi \esp{el básquetbol} \\ \hline
el voleibol & le volleyball & masculin \\ \hline
la natación & la natation & féminin : \esp{la natación} \\ \hline
el ciclismo & le cyclisme & \esp{montar en bicicleta} = faire du vélo \\ \hline
el boxeo & la boxe & attention : masculin en espagnol \\ \hline
\end{tabularx}
\end{center}

\begin{attention}[Faux amis et pièges]
\esp{El fútbol} signifie « le football », mais \esp{el fútbol americano} signifie « le football américain » : ne traduis jamais \esp{fútbol} par « football américain ». De même, \esp{un campeón} est un \emph{champion}, pas un champignon : le champignon se dit \esp{una seta}. Enfin, \esp{el boxeo} est masculin alors que « la boxe » est féminin en français.
\end{attention}

\courssub{Champ 2 — La pratique sportive}

\begin{center}
\begin{tabularx}{\linewidth}{|X|X|X|}\hline
\rowcolor{popgreenL} \textbf{Español} & \textbf{Français} & \textbf{Remarque} \\ \hline
el entrenamiento & l'entraînement & nom masculin \\ \hline
entrenar(se) & (s')entraîner & \esp{me entreno} = je m'entraîne \\ \hline
el partido & le match & \esp{un partido de fútbol} \\ \hline
el equipo & l'équipe & masculin : \esp{el equipo} \\ \hline
el/la entrenador(a) & l'entraîneur / l'entraîneuse & double forme \\ \hline
el campeonato & le championnat & \esp{el campeonato escolar} \\ \hline
el/la campeón(-ona) & le/la champion(ne) & \esp{la campeona} \\ \hline
ganar & gagner & \esp{ganar un partido} \\ \hline
perder & perdre & irrégulier : \esp{pierdo} (e $>$ ie) \\ \hline
marcar un gol & marquer un but & \esp{gol} sans accent \\ \hline
\end{tabularx}
\end{center}

\begin{exemplebox}[Exemples traduits]
\esp{Mi equipo se entrena todos los días.} « Mon équipe s'entraîne tous les jours. »

\esp{El entrenador habla con los jugadores antes del partido.} « L'entraîneur parle avec les joueurs avant le match. »

\esp{Samuel Eto'o marcó muchos goles y ganó muchos campeonatos.} « Samuel Eto'o a marqué beaucoup de buts et a gagné beaucoup de championnats. »

\esp{Ayer perdimos el partido, pero no perdemos la esperanza.} « Hier nous avons perdu le match, mais nous ne perdons pas espoir. »
\end{exemplebox}

\courssub{Champ 3 — Les amateurs et les lieux}

\begin{center}
\begin{tabularx}{\linewidth}{|X|X|X|}\hline
\rowcolor{poppurpleL} \textbf{Español} & \textbf{Français} & \textbf{Remarque} \\ \hline
el/la aficionado(a) & le/la supporter, amateur & \esp{soy aficionado al fútbol} \\ \hline
el estadio & le stade & \esp{el estadio está lleno} \\ \hline
el gimnasio & le gymnase, la salle de sport & prononcé « yim-na-sio » \\ \hline
la cancha / el campo & le terrain & \esp{el campo de fútbol} \\ \hline
el público & le public & masculin \\ \hline
el árbitro & l'arbitre & \esp{el árbitro pita} \\ \hline
\end{tabularx}
\end{center}

\begin{attention}[Être supporter de : aficionado A]
Pour dire « je suis supporter de... », l'espagnol utilise la préposition \esp{a} : \esp{Soy aficionado al fútbol.} « Je suis amateur de football. » Ne dis pas \esp{aficionado de fútbol} : c'est un calque du français.
\end{attention}

\begin{savaistu}[Un hincha y la afición]
En espagnol, un supporter passionné se dit aussi \esp{un hincha} (mot très utilisé en Amérique latine). \esp{La afición} désigne l'ensemble des supporters d'une équipe : \esp{La afición del equipo canta durante todo el partido} « Les supporters de l'équipe chantent pendant tout le match ». Au Cameroun, la afición des Lions Indomptables est célèbre dans toute l'Afrique !
\end{savaistu}

\courssub{Champ 4 — El ocio (les loisirs)}

\begin{definition}[El ocio]
\esp{El ocio} désigne l'ensemble des activités pratiquées pendant le temps libre (\esp{el tiempo libre}), par opposition au travail et aux études. Exemple : \esp{En mi tiempo libre, escucho música y salgo con amigos.} « Pendant mon temps libre, j'écoute de la musique et je sors avec des amis. »
\end{definition}

\begin{center}
\begin{tabularx}{\linewidth}{|X|X|X|}\hline
\rowcolor{poporangeL} \textbf{Español} & \textbf{Français} & \textbf{Remarque} \\ \hline
el ocio & les loisirs & masculin \\ \hline
el tiempo libre & le temps libre & \esp{en mi tiempo libre} \\ \hline
ver la televisión & regarder la télévision & \esp{veo} (ver irrégulier) \\ \hline
escuchar música & écouter de la musique & \esp{escuchar} = écouter \\ \hline
salir con amigos & sortir avec des amis & \esp{salgo} (salir irrégulier) \\ \hline
leer & lire & \esp{leer un libro} \\ \hline
bailar & danser & verbe régulier en -ar \\ \hline
montar en bicicleta & faire du vélo & aussi \esp{andar en bicicleta} \\ \hline
\end{tabularx}
\end{center}

\begin{exemplebox}[Exemples traduits]
\esp{Los fines de semana, salgo con mis amigos al mercado.} « Le week-end, je sors avec mes amis au marché. »

\esp{Mi hermana prefiere leer y escuchar música.} « Ma sœur préfère lire et écouter de la musique. »

\esp{¿Qué haces en tu tiempo libre?} « Que fais-tu pendant ton temps libre ? »
\end{exemplebox}

\courssub{Champ 5 — Le bien-être}

\begin{center}
\begin{tabularx}{\linewidth}{|X|X|X|}\hline
\rowcolor{popgoldL} \textbf{Español} & \textbf{Français} & \textbf{Remarque} \\ \hline
la salud física y mental & la santé physique et mentale & \esp{salud} féminin \\ \hline
el bienestar & le bien-être & un seul mot \\ \hline
sentirse bien/libre/fuerte & se sentir bien/libre/fort & \esp{me siento bien} \\ \hline
estar en forma & être en forme & \esp{estar}, pas \esp{ser} \\ \hline
relajarse & se détendre & verbe pronominal \\ \hline
\end{tabularx}
\end{center}

\begin{attention}[Sentirse : verbe pronominal et irrégulier]
\esp{Sentirse} est un verbe pronominal (avec \esp{se}) et irrégulier (e $>$ ie) : \esp{me siento, te sientes, se siente, nos sentimos, os sentís, se sienten}. Exemple : \esp{Después del deporte, me siento libre y fuerte.} « Après le sport, je me sens libre et fort. » N'oublie jamais le pronom réfléchi : on ne dit pas \esp{siento bien} pour « je me sens bien ». Remarque aussi l'accord de l'adjectif : \esp{Ella se siente fuerte.} « Elle se sent forte. »
\end{attention}

\courssub{Les verbes clés : practicar, jugar, hacer}

Trois verbes dominent le discours sur le sport. Leur choix n'est pas libre : il dépend du type d'activité.

\begin{propriete}[Practicar / jugar a / hacer]
\begin{itemize}
\item \esp{Jugar a + article + sport} : pour les jeux et les sports avec ballon ou adversaire : \esp{jugar al fútbol, al baloncesto, al voleibol, al ajedrez} (aux échecs). Attention : \esp{jugar} est irrégulier (u $>$ ue) : \esp{juego, juegas, juega, jugamos, jugáis, juegan}.
\item \esp{Practicar + sport} : pour une discipline pratiquée régulièrement : \esp{practicar atletismo, practicar natación}.
\item \esp{Hacer + activité} : pour les activités physiques sans ballon : \esp{hacer deporte, hacer ejercicio, hacer ciclismo, hacer gimnasia}.
\end{itemize}
\end{propriete}

\begin{attention}[La contraction obligatoire : a + el = al]
On dit \esp{jugar AL fútbol} : la préposition \esp{a} suivie de l'article \esp{el} se contracte obligatoirement en \esp{al}. Erreur fréquente des francophones : écrire \esp{jugar el fútbol} (calque de « jouer le football ») ou \esp{jugar a el fútbol}. Les deux sont fautives. De même : \esp{juego al baloncesto}, \esp{jugamos al voleibol}.
\end{attention}

\begin{exemplebox}[Les trois constructions en action]
\esp{Juego al fútbol con mi equipo.} « Je joue au football avec mon équipe. »

\esp{Practico atletismo.} « Je pratique l'athlétisme. »

\esp{Hago ejercicio por la mañana.} « Je fais de l'exercice le matin. »
\end{exemplebox}

\begin{popfigure}
\begin{center}
\begin{tikzpicture}[scale=1]
\node[draw=poppurple, fill=poppurpleL, rounded corners, thick, align=center] (dep) at (0,1.5) {\textbf{EL DEPORTE}};
\node[draw=popblue, fill=popblueL, rounded corners, thick, align=center, minimum width=3.2cm, minimum height=1.1cm] (jugar) at (0,3.4) {\textbf{JUGAR A}\\ \esp{al fútbol, al voleibol}\\ jeux et ballons};
\node[draw=popgreen, fill=popgreenL, rounded corners, thick, align=center, minimum width=3.2cm, minimum height=1.1cm] (practicar) at (-3.6,0) {\textbf{PRACTICAR}\\ \esp{atletismo, natación}\\ disciplines};
\node[draw=poporange, fill=poporangeL, rounded corners, thick, align=center, minimum width=3.2cm, minimum height=1.1cm] (hacer) at (3.6,0) {\textbf{HACER}\\ \esp{deporte, ejercicio}\\ activités physiques};
\draw[-{Stealth}, thick, popmuted] (dep) -- (jugar);
\draw[-{Stealth}, thick, popmuted] (dep) -- (practicar);
\draw[-{Stealth}, thick, popmuted] (dep) -- (hacer);
\end{tikzpicture}
\end{center}
\legende{Le triangle du choix : jugar a, practicar ou hacer selon le type d'activité.}
\end{popfigure}

\begin{exoresolu}[Complète avec la bonne construction]
Énoncé. Complète chaque phrase avec \esp{juego al}, \esp{practico} ou \esp{hago}, puis traduis.

1. \esp{... fútbol con mis amigos los sábados.}

2. \esp{... natación en la piscina de Douala.}

3. \esp{... ejercicio por la mañana para estar en forma.}
\tcblower
Solution détaillée.

1. \esp{Juego al fútbol con mis amigos los sábados.} « Je joue au football avec mes amis le samedi. » Le football est un sport avec ballon : on utilise \esp{jugar a}, avec la contraction \esp{al} et la forme irrégulière \esp{juego} (u $>$ ue).

2. \esp{Practico natación en la piscina de Douala.} « Je pratique la natation à la piscine de Douala. » La natation est une discipline : on utilise \esp{practicar} sans article.

3. \esp{Hago ejercicio por la mañana para estar en forma.} « Je fais de l'exercice le matin pour être en forme. » Activité physique sans ballon : \esp{hacer}, forme irrégulière \esp{hago} à la première personne.
\end{exoresolu}

\begin{aretenir}[L'essentiel du lexique]
\begin{itemize}
\item Les sports : \esp{el fútbol, el atletismo, el baloncesto, la natación, el ciclismo, el boxeo}.
\item La pratique : \esp{el entrenamiento, el partido, el equipo, el campeonato, ganar, perder, marcar un gol}.
\item Les lieux et les gens : \esp{el estadio, el gimnasio, el campo, el árbitro, el aficionado}.
\item Les loisirs : \esp{el ocio, el tiempo libre, escuchar música, salir con amigos, bailar}.
\item Le bien-être : \esp{la salud, el bienestar, estar en forma, sentirse fuerte, relajarse}.
\item Trois verbes à distinguer : \esp{jugar AL fútbol} (ballon), \esp{practicar atletismo} (discipline), \esp{hacer ejercicio} (activité).
\end{itemize}
\end{aretenir}

\begin{exercice}[À toi de jouer]
1. Traduis en espagnol : a) Je joue au volleyball avec mon équipe. b) Ma sœur pratique la natation. c) Nous faisons du sport pour être en forme. d) Les supporters remplissent le stade.

2. Classe ces mots dans les cinq champs lexicaux : \esp{el árbitro, bailar, el campeonato, la salud, el baloncesto, relajarse, el gimnasio, marcar un gol, el tiempo libre, el boxeo}.

3. Rédige trois phrases sur tes propres loisirs en utilisant \esp{jugar a}, \esp{practicar} ou \esp{hacer} et deux mots du champ du bien-être.
\end{exercice}

\begin{corrige}[Exercice : À toi de jouer]
1. a) \esp{Juego al voleibol con mi equipo.} b) \esp{Mi hermana practica natación.} c) \esp{Hacemos deporte para estar en forma.} d) \esp{Los aficionados llenan el estadio.}

2. Sports : \esp{el baloncesto, el boxeo}. Pratique : \esp{el campeonato, marcar un gol}. Lieux et gens : \esp{el árbitro, el gimnasio}. Loisirs : \esp{bailar, el tiempo libre}. Bien-être : \esp{la salud, relajarse}.

3. Exemple de production : \esp{Juego al fútbol los sábados y después me siento libre. Practico atletismo para estar en forma. Hago ejercicio por la mañana porque el deporte es bueno para la salud.}
\end{corrige}

Ce lexique est désormais ton outil de travail. Dans la prochaine section, nous apprendrons à exprimer nos goûts avec le verbe \esp{gustar} : tu pourras alors dire quels sports et quels loisirs tu aimes, en utilisant tout le vocabulaire de cette section.

\coursec{Grammaire 1 : l'expression du goût — gustar}

Dans la section précédente, nous avons découvert le vocabulaire du sport et des loisirs : \esp{el fútbol}, \esp{el atletismo}, \esp{el entrenamiento}, \esp{el ocio}, \esp{el aficionado}\ldots{} Mais savoir nommer les activités ne suffit pas : il faut maintenant pouvoir dire ce que l'on \cle{aime} et ce que l'on \cle{n'aime pas}. C'est exactement ce que permet le verbe \esp{gustar}, le verbe le plus important — et le plus piégeux — de cette leçon. Observons-le d'abord en action.

\courssub{Observation : un verbe qui fonctionne « à l'envers »}

Reprenons deux phrases tirées de notre texte support :

\begin{exemplebox}[Deux phrases du texte]
\esp{Me gusta mucho el deporte.} « J'aime beaucoup le sport. »

\esp{Me gusta practicar deporte.} « J'aime pratiquer un sport. »
\end{exemplebox}

Soulignons \esp{me gusta} et regardons ce qui suit : dans la première phrase, on trouve un \cle{nom} (\esp{el deporte}) ; dans la seconde, un \cle{infinitif} (\esp{practicar}). Posons-nous maintenant une question grammaticale : quel est le sujet du verbe \esp{gusta} ?

En français, dans « j'aime le sport », le sujet est \emph{je}. En espagnol, c'est l'inverse : le sujet de \esp{gustar} est \emph{la chose aimée}, et la personne qui éprouve le goût est exprimée par un \cle{pronom complément}. Littéralement, \esp{me gusta el deporte} signifie « le sport me plaît » : \esp{el deporte} est le sujet, \esp{me} est le complément.

\begin{popfigure}
\begin{tikzpicture}[node distance=1.2cm, >=Stealth]
\node[draw=popblue, fill=popblueL, rounded corners, align=center, font=\bfseries] (sujet) {EL DEPORTE\\ \footnotesize la chose aimée = SUJET};
\node[draw=poporange, fill=poporangeL, rounded corners, align=center, font=\bfseries, right=2.2cm of sujet] (verbe) {ME GUSTA\\ \footnotesize le verbe s'accorde\\ \footnotesize avec la chose};
\node[draw=popgreen, fill=popgreenL, rounded corners, align=center, font=\bfseries, right=2.2cm of verbe] (pers) {A MÍ\\ \footnotesize la personne\\ \footnotesize = complément};
\draw[->, very thick, popblue] (sujet) -- node[above, font=\footnotesize]{plaît} (verbe);
\draw[->, very thick, popgreen] (verbe) -- node[above, font=\footnotesize]{à qui ?} (pers);
\node[below=0.9cm of verbe, align=center, font=\footnotesize\itshape, text=popmuted] {En français : « j'aime le sport » — le sujet est « je ».\\ En espagnol : « el deporte me gusta » — le sujet est « el deporte ».};
\end{tikzpicture}
\legende{Le mécanisme inversé de \esp{gustar} : la chose aimée commande le verbe.}
\end{popfigure}

\begin{aretenir}[L'idée-clé]
\esp{Gustar} ne signifie pas « aimer » mais « \cle{plaire à} ». Toute la construction en découle : la chose aimée est le sujet du verbe, et la personne est introduite par un pronom complément d'objet indirect (\esp{me, te, le, nos, os, les}).
\end{aretenir}

\courssub{Les formes : gusta ou gustan}

Bonne nouvelle : au présent, on n'utilise pratiquement que \cle{deux formes} de \esp{gustar} :

\begin{propriete}[La règle d'accord de gustar]
\begin{itemize}
\item \esp{Gustar} + \cle{infinitif} $\rightarrow$ toujours \esp{gusta} (3\up{e} personne du singulier) : \esp{Me gusta correr.} « J'aime courir. » ; \esp{Me gusta jugar al fútbol.} « J'aime jouer au football. »
\item \esp{Gustar} + \cle{nom singulier} $\rightarrow$ \esp{gusta} : \esp{Me gusta el atletismo.} « J'aime l'athlétisme. »
\item \esp{Gustar} + \cle{nom pluriel} $\rightarrow$ \esp{gustan} : \esp{Me gustan los deportes.} « J'aime les sports. » ; \esp{Me gustan los partidos.} « J'aime les matchs. »
\end{itemize}
Même avec plusieurs infinitifs, le verbe reste au singulier : \esp{Me gusta nadar y correr.} « J'aime nager et courir. »
\end{propriete}

Le pronom complément, lui, indique \emph{qui} éprouve le goût. Il est souvent renforcé (ou précisé) par une \cle{forme tonique} introduite par la préposition \esp{a} :

\begin{popfigure}
\begin{tikzpicture}
\node[draw=popdark, fill=popcream, rounded corners, align=center, font=\bfseries] (titre) at (0,0) {GUSTAR au présent};
\node[draw=popblue, fill=popblueL, rounded corners, align=center, below left=1.1cm and -2.5cm of titre] (p1) {\esp{a mí} \textbf{me}\\ \esp{gusta(n)}};
\node[draw=popblue, fill=popblueL, rounded corners, align=center, below=1.1cm of titre] (p2) {\esp{a ti} \textbf{te}\\ \esp{gusta(n)}};
\node[draw=popblue, fill=popblueL, rounded corners, align=center, below right=1.1cm and -2.5cm of titre] (p3) {\esp{a él / ella / usted} \textbf{le}\\ \esp{gusta(n)}};
\node[draw=popgreen, fill=popgreenL, rounded corners, align=center, below=2.6cm of p1] (p4) {\esp{a nosotros(as)} \textbf{nos}\\ \esp{gusta(n)}};
\node[draw=popgreen, fill=popgreenL, rounded corners, align=center, below=2.6cm of p2] (p5) {\esp{a vosotros(as)} \textbf{os}\\ \esp{gusta(n)}};
\node[draw=popgreen, fill=popgreenL, rounded corners, align=center, below=2.6cm of p3] (p6) {\esp{a ellos / ellas / ustedes} \textbf{les}\\ \esp{gusta(n)}};
\draw[->, thick, popmuted] (titre) -- (p1);
\draw[->, thick, popmuted] (titre) -- (p2);
\draw[->, thick, popmuted] (titre) -- (p3);
\draw[->, thick, popmuted] (p2) -- (p4);
\draw[->, thick, popmuted] (p2) -- (p5);
\draw[->, thick, popmuted] (p2) -- (p6);
\end{tikzpicture}
\legende{Les six combinaisons : forme tonique + pronom complément + \esp{gusta} ou \esp{gustan}.}
\end{popfigure}

\begin{exemplebox}[Les pronoms toniques en contexte]
\esp{A mí me gusta el fútbol, pero a mi hermano le gusta el baloncesto.} « Moi, j'aime le football, mais mon frère, lui, aime le basket-ball. »

\esp{A nosotros nos gustan los partidos del sábado.} « Nous aimons les matchs du samedi. »

\esp{¿A ti te gusta el atletismo?} « Et toi, tu aimes l'athlétisme ? »
\end{exemplebox}

La forme tonique (\esp{a mí}, \esp{a ti}\ldots) est \cle{facultative} pour \esp{me, te, nos, os}, mais elle devient \cle{indispensable} pour lever l'ambiguïté de \esp{le} et \esp{les} : \esp{le gusta} peut vouloir dire « il aime », « elle aime » ou « vous aimez » ; on précise donc \esp{a él le gusta}, \esp{a ella le gusta}, \esp{a usted le gusta}. Elle sert aussi à insister ou à opposer : \esp{A mí me gusta, pero a ti no.} « Moi j'aime, mais pas toi. »

\begin{attention}[Ne confonds pas \esp{mí} et \esp{mi}]
\esp{A mí} (avec accent) signifie « à moi » ; \esp{mi} (sans accent) signifie « mon, ma ». Compare : \esp{A mí me gusta mi equipo.} « Moi, j'aime mon équipe. » L'accent distingue le pronom tonique du possessif : c'est une faute classique de copie.
\end{attention}

\courssub{Comparaison français / espagnol}

\begin{center}
\begin{tabularx}{\linewidth}{|X|X|X|}
\hline
\rowcolor{popblueL} Français & Español & Remarque \\
\hline
J'aime le sport. & Me gusta el deporte. & Sujet français : \emph{je} ; sujet espagnol : \esp{el deporte}. \\
\hline
Tu aimes les matchs. & Te gustan los partidos. & Nom pluriel $\rightarrow$ \esp{gustan}. \\
\hline
Il adore l'entraînement. & Le encanta el entrenamiento. & \esp{encantar} = même construction. \\
\hline
Nous aimons courir. & Nos gusta correr. & Infinitif $\rightarrow$ toujours \esp{gusta}. \\
\hline
Elle n'aime pas perdre. & No le gusta perder. & La négation \esp{no} précède le pronom. \\
\hline
Vous aimez le football ? & ¿Os gusta el fútbol? & \esp{os} = « vous » en Espagne. \\
\hline
\end{tabularx}
\end{center}

\courssub{Négation, gradation et verbes de la même famille}

\textbf{La négation} : on place simplement \esp{no} devant le pronom complément : \esp{No me gusta el boxeo.} « Je n'aime pas la boxe. » Pour renforcer : \esp{No me gusta nada el boxeo.} « Je n'aime pas du tout la boxe. »

\textbf{La gradation} permet de nuancer l'intensité du goût :

\begin{itemize}
\item \esp{me gusta} : j'aime ;
\item \esp{me gusta mucho} : j'aime beaucoup ;
\item \esp{me encanta} : j'adore (litt. « cela me ravit ») ;
\item \esp{no me gusta (nada)} : je n'aime pas (du tout).
\end{itemize}

\begin{exemplebox}[Exemples traduits]
\esp{A mi hermano le encanta el fútbol.} « Mon frère adore le football. »

\esp{¿Te gusta el deporte?} « Aimes-tu le sport ? »

\esp{No nos gusta nada perder.} « Nous n'aimons pas du tout perdre. »

\esp{A mis amigos les gustan mucho los partidos de la selección de Camerún.} « Mes amis aiment beaucoup les matchs de la sélection du Cameroun. »
\end{exemplebox}

\textbf{Les verbes de la même famille} (utiles pour le Bac) se construisent exactement comme \esp{gustar} :

\begin{center}
\begin{tabularx}{\linewidth}{|X|X|X|}
\hline
\rowcolor{popgreenL} Verbe & Sens & Exemple \\
\hline
\esp{encantar} & adorer, ravir & \esp{Me encanta nadar.} « J'adore nager. » \\
\hline
\esp{interesar} & intéresser & \esp{Le interesa el atletismo.} « L'athlétisme l'intéresse. » \\
\hline
\esp{aburrir} & ennuyer & \esp{Nos aburren los entrenamientos largos.} « Les longs entraînements nous ennuient. » \\
\hline
\esp{doler} & faire mal & \esp{Me duelen las piernas.} « J'ai mal aux jambes. » \\
\hline
\end{tabularx}
\end{center}

\begin{aretenir}[Et aux autres temps ?]
La construction reste identique à tous les temps ; seule la terminaison du verbe change. À l'imparfait : \esp{De niño, me gustaba el atletismo.} « Enfant, j'aimais l'athlétisme. » Au futur : \esp{Te gustará el partido.} « Le match te plaira. » Retiens aussi le conditionnel de politesse et de souhait : \esp{Me gustaría practicar más deporte.} « J'aimerais pratiquer davantage le sport. » — une formule très appréciée à l'expression écrite du Bac.
\end{aretenir}

\courssub{Méthode de traduction et exercice résolu}

\begin{methode}[Traduire « j'aime faire du sport »]
\begin{enumerate}
\item Identifier \cle{la chose aimée} : ici « faire du sport », c'est-à-dire un infinitif $\rightarrow$ \esp{hacer deporte}.
\item Choisir la forme du verbe : devant un infinitif, toujours \esp{gusta}.
\item Placer le pronom complément correspondant à la personne : « j' » $\rightarrow$ \esp{me}.
\item Assembler dans l'ordre espagnol : pronom + verbe + chose aimée $\rightarrow$ \esp{Me gusta hacer deporte.}
\item (Optionnel) Ajouter la forme tonique pour insister : \esp{A mí me gusta hacer deporte.}
\end{enumerate}
\end{methode}

\begin{exoresolu}[Traduction guidée]
Énoncé : traduis en espagnol les phrases suivantes. a) « Mon équipe adore les entraînements. » b) « Est-ce que vous aimez le basket-ball ? » (vosotros) c) « Je n'aime pas du tout perdre un match. »
\tcblower
Solution détaillée :

a) La chose aimée est \esp{los entrenamientos} (nom pluriel) $\rightarrow$ verbe au pluriel ; « adore » $\rightarrow$ \esp{encantan} ; la personne est « mon équipe » (3\up{e} pers. sing.) $\rightarrow$ \esp{le}, précisée par \esp{a mi equipo}. Réponse : \esp{A mi equipo le encantan los entrenamientos.}

b) La chose aimée est \esp{el baloncesto} (nom singulier) $\rightarrow$ \esp{gusta} ; « vous » (vosotros) $\rightarrow$ \esp{os}. Réponse : \esp{¿Os gusta el baloncesto?}

c) La chose aimée est un infinitif, \esp{perder un partido} $\rightarrow$ \esp{gusta} ; négation renforcée : \esp{no me gusta nada}. Réponse : \esp{No me gusta nada perder un partido.}
\end{exoresolu}

\courssub{Les pièges des francophones}

\begin{attention}[Erreur n°1 : « yo gusto » n'existe pas !]
Pour dire « j'aime », ne traduis jamais mot à mot par \esp{yo gusto}. Le sujet du verbe est la chose aimée, pas la personne. On dit \esp{me gusta el deporte} (litt. « le sport me plaît »). C'est l'erreur la plus fréquente des candidats francophones au Bac — les correcteurs la repèrent immédiatement.
\end{attention}

\begin{attention}[L'article défini est obligatoire]
Après \esp{gustar}, le nom est toujours précédé de l'article défini : \esp{me gusta EL fútbol}, \esp{me gustan LOS deportes}. Dire \esp{me gusta fútbol} est une faute. Retiens : en espagnol, on aime \emph{le} football, \emph{la} natation, \emph{les} matchs.
\end{attention}

\begin{attention}[Attention à l'accord gusta / gustan]
Le verbe s'accorde avec ce qui suit : \esp{me gusta el partido} (singulier) mais \esp{me gustan los partidos} (pluriel). Et n'oublie pas la contraction : « jouer au foot » se dit \esp{jugar AL fútbol} (\esp{a} + \esp{el} = \esp{al}) : \esp{Me gusta jugar al fútbol.}
\end{attention}

\begin{exercice}[À toi de jouer]
1. Traduis en espagnol : a) « J'aime beaucoup l'athlétisme. » b) « Nous adorons les matchs du samedi. » c) « Elle n'aime pas la boxe. » d) « Aimes-tu courir le matin ? »

2. Corrige les fautes : a) \esp{Yo gusto el deporte.} b) \esp{Me gusta fútbol.} c) \esp{Me gusta los entrenamientos.}
\end{exercice}

\begin{corrige}[Exercice « À toi de jouer »]
1. a) \esp{Me gusta mucho el atletismo.} b) \esp{Nos encantan los partidos del sábado.} c) \esp{No le gusta el boxeo.} (ou, pour lever l'ambiguïté : \esp{A ella no le gusta el boxeo.}) d) \esp{¿Te gusta correr por la mañana?}

2. a) \esp{Me gusta el deporte.} (jamais \esp{yo gusto}) b) \esp{Me gusta el fútbol.} (article obligatoire) c) \esp{Me gustan los entrenamientos.} (nom pluriel $\rightarrow$ \esp{gustan})
\end{corrige}

Tu sais maintenant exprimer tes goûts sportifs. Dans la section suivante, nous verrons comment dire à quelle \cle{fréquence} tu pratiques ces activités, grâce au verbe \esp{soler} suivi de l'infinitif.

\coursec{Grammaire 2 : l'habitude et la fréquence — soler + infinitif}

Dans la section précédente, nous avons appris à exprimer nos goûts avec \esp{gustar} : \esp{Me gusta el fútbol}, \esp{Nos gusta el deporte}. Mais dire ce que l'on \emph{aime} ne suffit pas : il faut aussi pouvoir dire ce que l'on \emph{fait régulièrement}. Aimez-vous le sport... et le pratiquez-vous vraiment ? Combien de fois par semaine ? C'est exactement ce que permet le verbe \esp{soler}, que nous allons découvrir maintenant grâce aux phrases de notre texte support.

\courssub{Observation : partons des exemples du texte}

Reprenons deux phrases du texte \esp{El deporte, una fuente de bienestar} :

\begin{exemplebox}[Phrases observées]
\esp{Suelo practicar deporte cuatro veces por semana.} « D'habitude, je pratique un sport quatre fois par semaine. »\par
\esp{Él suele montar en bicicleta los domingos.} « Il a l'habitude de faire du vélo le dimanche. »
\end{exemplebox}

Observons attentivement ces deux phrases et posons-nous trois questions :

\begin{enumerate}
\item Quel verbe est conjugué ? \esp{suelo} et \esp{suele} sont des formes du verbe \esp{soler}.
\item Sous quelle forme se trouve le verbe d'action ? \esp{practicar} et \esp{montar} sont à l'\textbf{infinitif}.
\item Que devient le \textbf{o} de \esp{soler} ? Il devient \textbf{ue} : \esp{suelo}, \esp{suele}... mais attention, pas partout !
\end{enumerate}

Nous pouvons déjà deviner la structure : \textbf{soler conjugué + infinitif} exprime une action habituelle, régulière. En français, on dirait « d'habitude + verbe conjugué » ou « avoir l'habitude de + infinitif ».

\courssub{La règle : soler + infinitif}

\begin{propriete}[Soler + infinitif : l'habitude]
Le verbe \esp{soler} suivi d'un \textbf{infinitif} exprime une \textbf{habitude}, une action répétée régulièrement. Il se traduit en français par « avoir l'habitude de », « d'habitude », « généralement ».\par
\esp{Suelo entrenar los martes.} = D'habitude, je m'entraîne le mardi.\par\medskip
\textbf{Changement orthographique o $>$ ue :} \esp{soler} est un verbe à \cle{radical changeant} : le \textbf{o} du radical devient \textbf{ue} dans la chaussure (trois personnes du singulier + 3e personne du pluriel). La 1re et la 2e personne du pluriel \textbf{ne changent pas} : \esp{solemos}, \esp{soléis}.
\end{propriete}

Voici la conjugaison complète au présent de l'indicatif :

\begin{center}
\begin{tabular}{|l|l|l|}
\hline
\rowcolor{popblueL} \textbf{Personne} & \textbf{Forme} & \textbf{Traduction} \\
\hline
yo & \esp{suelo} & j'ai l'habitude de \\
\hline
tú & \esp{sueles} & tu as l'habitude de \\
\hline
él / ella / usted & \esp{suele} & il / elle a l'habitude de \\
\hline
nosotros / nosotras & \esp{solemos} & nous avons l'habitude de \\
\hline
vosotros / vosotras & \esp{soléis} & vous avez l'habitude de \\
\hline
ellos / ellas / ustedes & \esp{suelen} & ils / elles ont l'habitude de \\
\hline
\end{tabular}
\end{center}

\begin{popfigure}
\begin{tikzpicture}
\node[draw=popblue, fill=popblueL, rounded corners, font=\bfseries] (titre) at (0,0) {SOLER + infinitif = l'habitude};
\node[draw=poporange, fill=poporangeL, rounded corners, align=center] (s1) at (-4.5,-1.6) {yo\\ \textbf{suelo}\\ j'ai l'habitude de};
\node[draw=poporange, fill=poporangeL, rounded corners, align=center] (s2) at (0,-1.6) {tú\\ \textbf{sueles}\\ tu as l'habitude de};
\node[draw=poporange, fill=poporangeL, rounded corners, align=center] (s3) at (4.5,-1.6) {él / ella\\ \textbf{suele}\\ il / elle a l'habitude de};
\node[draw=popgreen, fill=popgreenL, rounded corners, align=center] (p1) at (-4.5,-3.6) {nosotros\\ \textbf{solemos}\\ (o inchangé !)};
\node[draw=popgreen, fill=popgreenL, rounded corners, align=center] (p2) at (0,-3.6) {vosotros\\ \textbf{soléis}\\ (o inchangé !)};
\node[draw=poporange, fill=poporangeL, rounded corners, align=center] (p3) at (4.5,-3.6) {ellos / ustedes\\ \textbf{suelen}\\ ils ont l'habitude de};
\draw[-{Stealth}, popblue, thick] (titre) -- (s1);
\draw[-{Stealth}, popblue, thick] (titre) -- (s2);
\draw[-{Stealth}, popblue, thick] (titre) -- (s3);
\draw[-{Stealth}, popblue, thick] (titre) -- (p1);
\draw[-{Stealth}, popblue, thick] (titre) -- (p2);
\draw[-{Stealth}, popblue, thick] (titre) -- (p3);
\end{tikzpicture}
\legende{Conjugaison de \esp{soler} au présent : en orange, les formes avec changement o $>$ ue ; en vert, les deux formes sans changement.}
\end{popfigure}

\begin{attention}[Le changement o $>$ ue ne s'applique pas partout]
Ne dites jamais \esp{nosotros suelemos} ! Comme pour \esp{poder} (\esp{puedo} mais \esp{podemos}) ou \esp{querer} (\esp{quiero} mais \esp{queremos}), le changement de voyelle ne touche \textbf{pas} les personnes \esp{nosotros} et \esp{vosotros} : \esp{solemos}, \esp{soléis}. Retenez l'image de la « chaussure » : le changement se produit à l'intérieur de la chaussure (singulier + 3e personne du pluriel), jamais à l'extérieur.
\end{attention}

\courssub{Emploi et équivalents français}

\esp{Soler} s'emploie pour parler de ce qui se répète dans le temps : un entraînement hebdomadaire, une habitude familiale, une routine. Comparons les deux langues :

\begin{center}
\begin{tabularx}{\linewidth}{|X|X|X|}
\hline
\rowcolor{popblueL} \textbf{Français} & \textbf{Espagnol} & \textbf{Remarque} \\
\hline
D'habitude, je fais du sport. & \esp{Suelo hacer deporte.} & L'espagnol conjugue \esp{soler} ; le verbe d'action reste à l'infinitif. \\
\hline
Nous avons l'habitude de jouer le samedi. & \esp{Solemos jugar los sábados.} & « avoir l'habitude de » = \esp{soler} + infinitif. \\
\hline
Généralement, elle court le matin. & \esp{Ella suele correr por la mañana.} & « généralement » se rend souvent par \esp{soler}. \\
\hline
D'habitude, regardes-tu les matchs ? & \esp{¿Sueles ver los partidos?} & À la forme interrogative, pas de mot supplémentaire : l'intonation suffit. \\
\hline
\end{tabularx}
\end{center}

\begin{aretenir}[La différence essentielle]
En français, on dit « d'habitude + verbe \textbf{conjugué} » : « d'habitude, je \emph{fais} du sport ». En espagnol, on conjugue \esp{soler} et le verbe d'action reste à l'\textbf{infinitif} : \esp{Suelo hacer deporte}. Ne conjuguez jamais les deux verbes : \esp{suelo hago} est une faute grave !
\end{aretenir}

\begin{exemplebox}[Exemples traduits]
\esp{Solemos jugar un partido el sábado.} « Nous avons l'habitude de jouer un match le samedi. »\par
\esp{¿Sueles hacer deporte?} « As-tu l'habitude de faire du sport ? »\par
\esp{Pablo suele ver los partidos del Real Madrid.} « Pablo regarde habituellement les matchs du Real Madrid. »\par
\esp{Mis amigos suelen ir al estadio los fines de semana.} « Mes amis vont d'habitude au stade les week-ends. »\par
\esp{En Yaoundé, los jóvenes suelen jugar al fútbol por la tarde.} « À Yaoundé, les jeunes jouent généralement au football l'après-midi. »\par
\esp{No suelo ver la televisión; prefiero entrenar.} « Je n'ai pas l'habitude de regarder la télévision ; je préfère m'entraîner. »
\end{exemplebox}

Remarquez dans le dernier exemple que la négation se place simplement devant \esp{soler} : \esp{no suelo} + infinitif = « je n'ai pas l'habitude de ».

\courssub{Les adverbes et expressions de fréquence}

Pour préciser le rythme d'une habitude, l'espagnol utilise des adverbes de fréquence et des expressions de temps. Les voici, classées de la plus forte à la plus faible fréquence :

\begin{center}
\begin{tabularx}{\linewidth}{|X|X|X|}
\hline
\rowcolor{popblueL} \textbf{Español} & \textbf{Français} & \textbf{Exemple} \\
\hline
\esp{siempre} & toujours & \esp{Siempre entreno por la mañana.} \\
\hline
\esp{normalmente / generalmente} & normalement, généralement & \esp{Normalmente jugamos los sábados.} \\
\hline
\esp{a menudo} & souvent & \esp{Voy a menudo al gimnasio.} \\
\hline
\esp{a veces} & parfois & \esp{A veces monto en bicicleta.} \\
\hline
\esp{casi nunca} & presque jamais & \esp{Casi nunca veo los partidos.} \\
\hline
\esp{nunca} & jamais & \esp{Nunca falto al entrenamiento.} \\
\hline
\end{tabularx}
\end{center}

Et les expressions de temps qui accompagnent naturellement \esp{soler} :

\begin{itemize}
\item \esp{una vez / dos veces / cuatro veces por semana} : une fois / deux fois / quatre fois par semaine ;
\item \esp{todos los días} : tous les jours ;
\item \esp{los lunes, los martes...} : le lundi, le mardi (tous les lundis, tous les mardis) ;
\item \esp{los fines de semana} : les week-ends ;
\item \esp{por la mañana / por la tarde / por la noche} : le matin / l'après-midi / le soir.
\end{itemize}

\begin{popfigure}
\begin{tikzpicture}
\draw[-{Stealth}, very thick, popdark] (0,0) -- (12.5,0);
\foreach \x/\label/\col in {
  0/{siempre\\100 \%}/popgreen,
  2.4/{normalmente}/popgreen!70!popblue,
  4.8/{a menudo}/popblue,
  7.2/{a veces}/poporange,
  9.6/{casi nunca}/poporange!70!popink,
  12/{nunca\\0 \%}/popink} {
  \fill[\col] (\x,0) circle (0.12);
  \node[above, align=center, font=\small\bfseries, text=\col] at (\x,0.15) {\label};
}
\node[below, font=\small\itshape, text=popmuted] at (0,-0.25) {toujours};
\node[below, font=\small\itshape, text=popmuted] at (2.4,-0.25) {normalement};
\node[below, font=\small\itshape, text=popmuted] at (4.8,-0.25) {souvent};
\node[below, font=\small\itshape, text=popmuted] at (7.2,-0.25) {parfois};
\node[below, font=\small\itshape, text=popmuted] at (9.6,-0.25) {presque jamais};
\node[below, font=\small\itshape, text=popmuted] at (12,-0.25) {jamais};
\end{tikzpicture}
\legende{Frise de fréquence : de \esp{siempre} (100 \% du temps) à \esp{nunca} (0 \%).}
\end{popfigure}

\begin{attention}[« vez » devient « veces » au pluriel]
Le mot \esp{vez} (fois) fait son pluriel en \esp{veces} : \esp{una vez}, \esp{dos veces}, \esp{tres veces por semana}. N'écrivez jamais \esp{dos vez} ! Règle générale : les noms terminés en \textbf{-z} changent le \textbf{z} en \textbf{c} devant \textbf{-es} : \esp{luz} $>$ \esp{luces}, \esp{pez} $>$ \esp{peces}, \esp{vez} $>$ \esp{veces}.
\end{attention}

\begin{attention}[Les jours de la semaine : article et minuscule]
Pour exprimer une habitude, les jours de la semaine prennent l'article défini pluriel : \esp{los lunes} = « le lundi » (tous les lundis), \esp{los domingos} = « le dimanche ». De plus, en espagnol, les jours de la semaine \textbf{ne prennent pas de majuscule} : on écrit \esp{lunes}, \esp{martes}, \esp{domingo}, contrairement au français « Lundi, Mardi ». Et pour dire « samedi prochain », on utilise le singulier : \esp{el sábado}.
\end{attention}

\courssub{Rappel : practicar, verbe régulier en -ar}

Le verbe \esp{practicar} (pratiquer, faire [un sport]) est un verbe \textbf{régulier} du premier groupe en \textbf{-ar}. Il se conjugue comme \esp{hablar} :

\begin{center}
\begin{tabular}{|l|l|l|}
\hline
\rowcolor{popgreenL} \textbf{Personne} & \textbf{Forme} & \textbf{Traduction} \\
\hline
yo & \esp{practico} & je pratique \\
\hline
tú & \esp{practicas} & tu pratiques \\
\hline
él / ella / usted & \esp{practica} & il / elle pratique \\
\hline
nosotros / nosotras & \esp{practicamos} & nous pratiquons \\
\hline
vosotros / vosotras & \esp{practicáis} & vous pratiquez \\
\hline
ellos / ellas / ustedes & \esp{practican} & ils / elles pratiquent \\
\hline
\end{tabular}
\end{center}

\begin{exemplebox}[Practicar en contexte]
\esp{Practico atletismo en el estadio de Douala.} « Je pratique l'athlétisme au stade de Douala. »\par
\esp{¿Qué deporte practicas?} « Quel sport pratiques-tu ? »\par
\esp{Mi hermana practica natación tres veces por semana.} « Ma sœur fait de la natation trois fois par semaine. »
\end{exemplebox}

Attention à l'orthographe : en français on écrit « pratiquer » avec \emph{qu}, en espagnol c'est \esp{practicar} avec un seul \textbf{c}. Et ne confondez pas avec le faux ami : \esp{practicar un deporte} = « pratiquer un sport », jamais « pratiquer » au sens de « fréquenter ».

\courssub{Pour aller plus loin : soler à l'imparfait}

\begin{propriete}[Solía + infinitif : l'habitude au passé]
À l'\textbf{imparfait}, \esp{soler} exprime une habitude \textbf{passée} : \esp{solía practicar} = « j'avais l'habitude de pratiquer », « je pratiquais (autrefois) ».\par
Conjugaison à l'imparfait (régulière, sans changement de voyelle) : \esp{solía, solías, solía, solíamos, solíais, solían}.\par
\esp{De niño, solía jugar al fútbol todos los días.} « Enfant, je jouais au football tous les jours. »\par
\esp{Antes solíamos ver los partidos en casa de mi abuelo.} « Avant, nous avions l'habitude de regarder les matchs chez mon grand-père. »
\end{propriete}

Cette forme sera très utile au baccalauréat pour raconter des souvenirs ou comparer passé et présent : \esp{Antes solía ver los partidos, pero ahora prefiero jugar.} « Avant, je regardais les matchs, mais maintenant je préfère jouer. »

\begin{savaistu}[Un verbe « défectif » venu du latin]
\esp{Soler} vient du latin \emph{solere}, « avoir coutume, avoir l'habitude ». C'est ce que les grammairiens appellent un verbe \cle{défectif} : on ne l'emploie presque jamais qu'au présent et à l'imparfait de l'indicatif. Vous n'entendrez pratiquement jamais \esp{soler} au futur ou au subjonctif : pour dire « j'aurai l'habitude de », l'espagnol préfère d'autres tournures comme \esp{acostumbrar} (s'habituer à) ou simplement un adverbe de fréquence. Retenez donc : \esp{soler} = présent et imparfait, point final !
\end{savaistu}

\courssub{Exercice résolu}

\begin{exoresolu}[Conjugue et traduis]
Conjuguez \esp{soler} à la personne demandée, puis traduisez la phrase complète en français.\par
1. (yo) \esp{\ldots{} practicar deporte cuatro veces por semana.}\par
2. (nosotros) \esp{\ldots{} jugar un partido el sábado.}\par
3. (tú) \esp{¿\ldots{} montar en bicicleta los domingos?}\par
4. (ellos) \esp{\ldots{} ver los partidos del campeonato.}\par
5. (ella, imparfait) \esp{De niña, \ldots{} nadar todos los días.}
\tcblower
\textbf{Solutions détaillées :}\par
1. \esp{Suelo practicar deporte cuatro veces por semana.} « D'habitude, je pratique un sport quatre fois par semaine. » (1re personne du singulier : changement o $>$ ue.)\par
2. \esp{Solemos jugar un partido el sábado.} « Nous avons l'habitude de jouer un match le samedi. » (1re personne du pluriel : \textbf{pas} de changement de voyelle !)\par
3. \esp{¿Sueles montar en bicicleta los domingos?} « As-tu l'habitude de faire du vélo le dimanche ? » (2e personne du singulier : \esp{sueles}.)\par
4. \esp{Suelen ver los partidos del campeonato.} « Ils regardent habituellement les matchs du championnat. » (3e personne du pluriel : \esp{suelen}.)\par
5. \esp{De niña, solía nadar todos los días.} « Enfant, elle avait l'habitude de nager tous les jours. » (Imparfait : habitude passée, conjugaison régulière sans changement de voyelle.)
\end{exoresolu}

\begin{exercice}[À toi de jouer]
Traduisez en espagnol :\par
1. J'ai l'habitude de m'entraîner le mardi.\par
2. Nous jouons souvent au football le week-end.\par
3. As-tu l'habitude de regarder les matchs à la télévision ?\par
4. Mon équipe s'entraîne trois fois par semaine.\par
5. Avant, ils allaient au stade tous les dimanches.
\end{exercice}

\begin{corrige}[Exercice « À toi de jouer »]
1. \esp{Suelo entrenar los martes.} (ou \esp{Suelo entrenarme los martes} avec le verbe pronominal.)\par
2. \esp{Solemos jugar al fútbol los fines de semana.} ou \esp{A menudo jugamos al fútbol los fines de semana.}\par
3. \esp{¿Sueles ver los partidos en la televisión?}\par
4. \esp{Mi equipo entrena tres veces por semana.} ou \esp{Mi equipo suele entrenar tres veces por semana.} (Attention au pluriel : \esp{veces} !)\par
5. \esp{Antes solían ir al estadio todos los domingos.} (Imparfait de \esp{soler} pour une habitude passée.)
\end{corrige}

\begin{aretenir}[L'essentiel de la section]
\begin{itemize}
\item \esp{Soler} + infinitif = « avoir l'habitude de », « d'habitude » : \esp{Suelo practicar deporte.}
\item Changement o $>$ ue dans la chaussure : \esp{suelo, sueles, suele, suelen}, mais \esp{solemos, soléis}.
\item Le verbe d'action reste toujours à l'\textbf{infinitif} : jamais \esp{suelo hago}.
\item Fréquence : \esp{siempre, normalmente, a menudo, a veces, casi nunca, nunca} ; \esp{dos veces por semana} ; \esp{los lunes} (article + minuscule).
\item Au passé : \esp{solía} + infinitif = « j'avais l'habitude de ».
\end{itemize}
\end{aretenir}

Maintenant que nous savons exprimer nos goûts (\esp{me gusta}) et nos habitudes (\esp{suelo} + infinitif), nous sommes prêts à passer à la section suivante : la traduction, où ces deux outils nous serviront à traduire des phrases complètes sur le sport, en thème comme en version.

\coursec{Traduction : thème et version sur le sport}

Dans la section précédente, nous avons appris à exprimer l'habitude avec \esp{soler + infinitif} : \esp{Suelo entrenar los jueves} « J'ai l'habitude de m'entraîner le jeudi ». Nous savons aussi, depuis la grammaire 1, dire ce qui nous plaît avec \esp{gustar}. Le moment est venu de mettre ces deux outils à l'épreuve dans l'exercice roi de l'apprentissage des langues : la \cle{traduction}. Nous allons d'abord traduire de l'espagnol vers le français (\cle{version}), puis du français vers l'espagnol (\cle{thème}), en justifiant chaque choix comme de vrais traducteurs professionnels. Cette double compétence — version et thème — est au cœur des épreuves d'espagnol de la classe de Première et du Baccalauréat.

\courssub{La méthode de traduction : observer avant de traduire}

Traduire n'est pas remplacer mot à mot. C'est d'abord \emph{comprendre la structure} de la phrase, puis la reconstruire dans l'autre langue. Avant tout exercice, observons deux phrases modèles tirées de notre texte support :

\esp{Me gustan los deportes de equipo.} « J'aime les sports d'équipe. »

\esp{Suele correr por la mañana.} « D'habitude, elle court le matin. »

Dans la première phrase, le verbe \esp{gustan} est au pluriel parce que ce qui plaît, \esp{los deportes}, est pluriel ; le pronom \esp{me} se place avant le verbe. Dans la seconde, \esp{suele} est conjugué (3\up{e} personne) et l'infinitif \esp{correr} reste inchangé : c'est exactement la construction française « avoir l'habitude de + infinitif ».

\begin{methode}[Traduire une phrase avec « aimer »]
\begin{enumerate}
\item \textbf{Repérer ce qui est aimé} : est-ce un infinitif (\emph{courir}), un nom singulier (\emph{le football}) ou un nom pluriel (\emph{les matchs}) ?
\item \textbf{Choisir la forme du verbe} : infinitif ou nom singulier $\rightarrow$ \esp{gusta} ; nom pluriel $\rightarrow$ \esp{gustan}.
\item \textbf{Placer le pronom complément} (\esp{me, te, le, nos, os, les}) \emph{avant} le verbe.
\item \textbf{Mettre l'article défini} devant le nom aimé : \esp{el fútbol}, \esp{los deportes}. En espagnol, on aime toujours « \emph{le} quelque chose ».
\item \textbf{Vérifier} : relire la phrase entière et contrôler l'accord.
\end{enumerate}
\end{methode}

\begin{exemplebox}[Deux modèles à imiter]
\esp{Me gustan los deportes de equipo.} « J'aime les sports d'équipe. » (nom pluriel $\rightarrow$ \esp{gustan} + article \esp{los}).

\esp{Suele correr por la mañana.} « D'habitude, elle court le matin. » (\esp{soler} conjugué + infinitif ; \esp{por la mañana} = le matin, sans article en français).
\end{exemplebox}

\begin{popfigure}
\begin{tikzpicture}[
  node distance=0.55cm,
  etape/.style={rectangle, rounded corners, draw=popblue, fill=popblueL, text width=5.6cm, align=center, font=\small, inner sep=5pt},
  fleche/.style={-{Stealth[length=2.5mm]}, thick, popdark}
]
\node[etape] (e1) {\textbf{1. Lire} la phrase entière et comprendre le sens};
\node[etape, below=of e1] (e2) {\textbf{2. Identifier la structure} : goût ? habitude ? action simple ?};
\node[etape, below=of e2] (e3) {\textbf{3. Choisir l'outil} : \esp{gustar} / \esp{soler + inf.} / verbe simple};
\node[etape, below=of e3] (e4) {\textbf{4. Conjuguer} : \esp{gusta/gustan}, \esp{suelo/suele...}};
\node[etape, below=of e4, draw=poporange, fill=poporangeL] (e5) {\textbf{5. Vérifier} : article après \esp{gustar}, accord, \esp{al}, \esp{veces}};
\draw[fleche] (e1) -- (e2);
\draw[fleche] (e2) -- (e3);
\draw[fleche] (e3) -- (e4);
\draw[fleche] (e4) -- (e5);
\end{tikzpicture}
\legende{Organigramme de la méthode de traduction : cinq étapes, de la lecture à la vérification.}
\end{popfigure}

\courssub{Version : traduction guidée d'un extrait du texte support}

Voici un extrait du texte support de la leçon. Lisez-le d'abord en silence, puis à voix haute, avant de le traduire. Soulignez au passage tous les verbes conjugués : c'est le premier réflexe du traducteur.

\begin{textebox}[El deporte, una fuente de bienestar (extrait)]
En mi barrio de Yaoundé, a muchos jóvenes les gusta el fútbol. Los fines de semana, suelen jugar al fútbol en el campo del barrio. A mi hermano le gustan los deportes de equipo, pero a mí me gusta más el atletismo. Suelo entrenar tres veces por semana porque el deporte es bueno para la salud. Los médicos dicen que practicar un deporte también es bueno para la mente: quien hace deporte suele estar más contento y duerme mejor.
\end{textebox}

\textbf{Glossaire :} \esp{el barrio} : le quartier ; \esp{el campo} : le terrain ; \esp{la mente} : l'esprit ; \esp{quien} : celui qui ; \esp{duerme} : (il) dort, verbe \esp{dormir}.

\begin{exoresolu}[Version guidée : cinq phrases]
Traduisez en français : (1) \esp{En mi barrio de Yaoundé, a muchos jóvenes les gusta el fútbol.} (2) \esp{Los fines de semana, suelen jugar al fútbol en el campo del barrio.} (3) \esp{A mi hermano le gustan los deportes de equipo, pero a mí me gusta más el atletismo.} (4) \esp{Suelo entrenar tres veces por semana porque el deporte es bueno para la salud.} (5) \esp{Quien hace deporte suele estar más contento y duerme mejor.}
\tcblower
\textbf{(1)} « Dans mon quartier de Yaoundé, beaucoup de jeunes aiment le football. » Justification : \esp{les gusta} = « cela plaît à eux » ; ce qui plaît est \esp{el fútbol} (singulier), donc \esp{gusta} ; en français on inverse la construction : « les jeunes aiment ».

\textbf{(2)} « Les week-ends, ils ont l'habitude de jouer au football sur le terrain du quartier. » Justification : \esp{suelen + infinitif} = habitude à la 3\up{e} personne du pluriel ; \esp{jugar al fútbol} : \esp{al} = \esp{a + el}, contraction obligatoire, car on joue \emph{à} un sport ; \esp{los fines de semana} avec article exprime l'habitude (« les week-ends »).

\textbf{(3)} « Mon frère aime les sports d'équipe, mais moi je préfère l'athlétisme. » Justification : \esp{le gustan} au pluriel car \esp{los deportes} est pluriel ; \esp{a mi hermano} et \esp{a mí} précisent à qui cela plaît ; \esp{me gusta más} = « je préfère ».

\textbf{(4)} « J'ai l'habitude de m'entraîner trois fois par semaine parce que le sport est bon pour la santé. » Justification : \esp{suelo entrenar} = habitude à la 1\up{re} personne ; \esp{veces} = « fois » ; \esp{por semana} = « par semaine » ; \esp{es bueno} avec \esp{ser} car c'est une qualité permanente.

\textbf{(5)} « Celui qui fait du sport est généralement plus content et dort mieux. » Justification : \esp{suele estar} = habitude + \esp{estar} (état temporaire : être content) ; \esp{hace deporte} = « fait du sport ».
\end{exoresolu}

\courssub{Thème — série 1 : phrases simples}

Passons au \cle{thème}, c'est-à-dire à la traduction du français vers l'espagnol. Appliquez l'organigramme : lire, identifier, choisir, conjuguer, vérifier.

\begin{exoresolu}[Thème, série 1 : trois phrases simples]
Traduisez en espagnol : (1) « J'aime le football. » (2) « Nous avons l'habitude de nous entraîner le jeudi. » (3) « Mon équipe gagne souvent les matchs. »
\tcblower
\textbf{(1)} \esp{Me gusta el fútbol.} Structure de goût : ce qui est aimé est un nom singulier $\rightarrow$ \esp{gusta} ; article obligatoire \esp{el} ; pronom \esp{me} avant le verbe. Piège évité : ne jamais dire \esp{yo gusto}.

\textbf{(2)} \esp{Solemos entrenarnos los jueves.} Structure d'habitude : \esp{soler} conjugué à la 1\up{re} personne du pluriel (\esp{solemos}) + infinitif ; « nous entraîner » est pronominal $\rightarrow$ \esp{entrenarnos} ; « le jeudi » avec l'habitude se dit \esp{los jueves} (article + pluriel).

\textbf{(3)} \esp{Mi equipo gana a menudo los partidos.} Phrase simple : sujet \esp{mi equipo} (singulier) $\rightarrow$ \esp{gana} ; « souvent » = \esp{a menudo} ; « les matchs » = \esp{los partidos}.
\end{exoresolu}

\courssub{Thème — série 2 : phrases complexes}

Les phrases complexes combinent plusieurs structures : goût + opposition, habitude + cause. Il faut découper la phrase en morceaux et traiter chaque morceau avec la méthode.

\begin{exoresolu}[Thème, série 2 : deux phrases complexes]
Traduisez : (1) « Mon frère adore l'athlétisme mais il n'aime pas du tout perdre. » (2) « D'habitude, je fais du sport trois fois par semaine parce que c'est bon pour la santé. »
\tcblower
\textbf{(1)} \esp{A mi hermano le encanta el atletismo, pero no le gusta nada perder.} Découpage : « adorer » = \esp{encantar}, qui fonctionne exactement comme \esp{gustar} ; ce qui est adoré, \esp{el atletismo}, est singulier $\rightarrow$ \esp{le encanta} ; \esp{a mi hermano} précise la personne. « Il n'aime pas du tout perdre » : ce qui n'est pas aimé est un \emph{infinitif} (\esp{perder}) $\rightarrow$ \esp{gusta} ; « pas du tout » = \esp{no... nada} ; la négation \esp{no} se place avant le pronom : \esp{no le gusta nada}.

\textbf{(2)} \esp{Normalmente, hago deporte tres veces por semana porque es bueno para la salud.} Découpage : « d'habitude » = \esp{normalmente} (autre solution possible : \esp{suelo hacer deporte} avec \esp{soler}) ; « je fais du sport » = \esp{hago deporte}, verbe \esp{hacer}, 1\up{re} personne irrégulière \esp{hago} ; « trois fois » = \esp{tres veces} ; « par semaine » = \esp{por semana}, sans article ; « c'est bon pour la santé » = \esp{es bueno para la salud}, avec \esp{ser} (qualité permanente) et la préposition \esp{para}.
\end{exoresolu}

\begin{attention}[« Par semaine » et « le week-end »]
« Par semaine » se traduit \esp{por semana}, \textbf{sans article} : ne dites jamais \esp{por la semana}. De même : \esp{por día} « par jour », \esp{por mes} « par mois ». En revanche, « le week-end » se dit \esp{el fin de semana}, et pour exprimer l'habitude on utilise le pluriel avec article : \esp{los fines de semana} « les week-ends, le week-end (en général) ».
\end{attention}

\begin{attention}[« C'est bon pour la santé » : ser, jamais estar]
\esp{Es bueno para la salud.} « Bon » exprime ici une qualité permanente, une vérité générale : on utilise donc \esp{ser}. La tournure \esp{está bueno} existe, mais elle signifie « il est appétissant » (en parlant d'un plat) ou, familièrement, « elle est belle » : jamais « c'est bon pour la santé ». Retenez aussi la préposition : « bon \emph{pour} » = \esp{bueno para}.
\end{attention}

\begin{attention}[Les verbes de la famille de \esp{gustar}]
\esp{encantar} « adorer », \esp{interesar} « intéresser » et \esp{preferir} « préférer » se construisent comme \esp{gustar} : pronom complément + verbe accordé avec la chose aimée. \esp{Me encanta el fútbol.} « J'adore le football. » \esp{¿Te interesa el atletismo?} « L'athlétisme t'intéresse-t-il ? » Attention : \esp{preferir} est un verbe à diphtongaison (\esp{prefiero, prefieres, prefiere...}) mais il garde la même construction : \esp{Mi hermana prefiere el atletismo.} « Ma sœur préfère l'athlétisme. »
\end{attention}

\courssub{Corrections collectives et bilan des pièges}

Lors de la correction collective, chaque élève justifie sa traduction à voix haute : « J'ai mis \esp{gustan} parce que... ». Cette verbalisation est la clé de la mémorisation. Voici le tableau récapitulatif des pièges rencontrés dans cette séance.

\begin{center}
\begin{tabularx}{\linewidth}{|X|X|X|}
\hline
\rowcolor{popblueL} \textbf{Piège} & \textbf{Erreur du francophone} & \textbf{Forme correcte} \\
\hline
Article après \esp{gustar} & \esp{Me gusta fútbol.} & \esp{Me gusta el fútbol.} \\
\hline
Contraction \esp{a + el} & \esp{Juego a el fútbol.} & \esp{Juego al fútbol.} \\
\hline
Accord \esp{gusta/gustan} & \esp{Me gusta los deportes.} & \esp{Me gustan los deportes.} \\
\hline
Sujet français « je » & \esp{Yo gusto el fútbol.} & \esp{Me gusta el fútbol.} \\
\hline
« Fois » & \esp{tres tiempos} & \esp{tres veces} \\
\hline
« Par semaine » & \esp{por la semana} & \esp{por semana} \\
\hline
Jours et habitude & \esp{el jueves} (un seul jeudi) ou \esp{en jueves} & \esp{los jueves} « le jeudi » \\
\hline
« Bon pour la santé » & \esp{está bueno para la salud} & \esp{es bueno para la salud} \\
\hline
\end{tabularx}
\end{center}

\begin{aretenir}[Les cinq commandements du traducteur]
\begin{enumerate}
\item Après \esp{gustar}, le nom aimé prend toujours l'article : \esp{Me gusta el atletismo.}
\item \esp{gusta} avec un infinitif ou un singulier, \esp{gustan} avec un pluriel.
\item On joue \emph{à} un sport : \esp{jugar al fútbol} (contraction \esp{al} obligatoire).
\item L'habitude : \esp{soler} conjugué + infinitif ; jours et week-ends au pluriel avec article : \esp{los lunes}, \esp{los fines de semana}.
\item « Fois » = \esp{veces} ; « par semaine » = \esp{por semana} ; « c'est bon pour » = \esp{es bueno para}.
\end{enumerate}
\end{aretenir}

\begin{exercice}[À vous de traduire]
Traduisez en espagnol, puis justifiez chaque choix à l'écrit :
\begin{enumerate}
\item « Nous aimons les matchs de football le samedi. »
\item « D'habitude, elle court le matin parce que c'est bon pour la santé. »
\item « Mon équipe a l'habitude de s'entraîner deux fois par semaine. »
\item « J'aime beaucoup l'athlétisme, mais je n'aime pas perdre. »
\end{enumerate}
\end{exercice}

\begin{corrige}[Exercice : À vous de traduire]
\textbf{1.} \esp{Nos gustan los partidos de fútbol los sábados.} Nom pluriel $\rightarrow$ \esp{gustan} ; article \esp{los} ; habitude le samedi $\rightarrow$ \esp{los sábados}.

\textbf{2.} \esp{Suele correr por la mañana porque es bueno para la salud.} Habitude $\rightarrow$ \esp{suele} + infinitif ; qualité permanente $\rightarrow$ \esp{es bueno} ; préposition \esp{para}.

\textbf{3.} \esp{Mi equipo suele entrenarse dos veces por semana.} \esp{soler} à la 3\up{e} personne du singulier ; verbe pronominal \esp{entrenarse} ; \esp{veces} ; \esp{por semana} sans article.

\textbf{4.} \esp{Me gusta mucho el atletismo, pero no me gusta perder.} Nom singulier et infinitif $\rightarrow$ \esp{gusta} dans les deux cas ; négation \esp{no} avant le pronom \esp{me}.
\end{corrige}

Vous savez désormais traduire, dans les deux sens, les trois structures essentielles de cette leçon : le goût avec \esp{gustar}, l'habitude avec \esp{soler}, et le jeu avec \esp{jugar a}. La prochaine section vous permettra de réutiliser ces acquis dans un commentaire de texte et une production guidée sur le sport et le bien-être.

\coursec{Commentaire de texte et production guidée}

Après avoir traduit des phrases sur le sport en thème et en version, nous savons maintenant manier le vocabulaire et la grammaire de cette leçon. Il est temps de les réunir dans l'exercice roi de la classe de Première et du Baccalauréat : le \cle{commentaire de texte}, suivi d'une \cle{production écrite guidée}. L'objectif est double : montrer que vous avez compris le texte support \esp{El deporte, una fuente de bienestar}, et apprendre à exprimer votre propre avis en espagnol correct.

\courssub{Qu'est-ce qu'un commentaire de texte ?}

\begin{definition}[Le commentaire de texte]
Le \cle{commentaire de texte} (\esp{comentario de texto}) est un exercice écrit dans lequel l'élève présente un texte, en dégage l'idée principale, analyse les idées secondaires et le vocabulaire, puis donne son avis personnel argumenté. En espagnol, on dit : \esp{comentar un texto}. Au Bac, il est rédigé entièrement en espagnol.
\end{definition}

Un bon commentaire n'est ni un résumé mot à mot, ni une simple opinion. C'est un texte organisé en quatre étapes, comme le montre l'organigramme suivant.

\begin{popfigure}
\begin{tikzpicture}[
  node distance=0.55cm,
  etape/.style={rectangle, rounded corners=4pt, draw=popblue, fill=popblueL, text width=6.2cm, align=center, font=\small, inner sep=6pt},
  fleche/.style={-{Stealth[length=3mm]}, thick, poporange}
]
\node[etape, align=center] (e1){\textbf{1. Presentación}\\ \esp{El texto es un artículo... trata de...}};
\node[etape, below=of e1, align=center] (e2){\textbf{2. Idea principal}\\ \esp{La idea principal es que...}};
\node[etape, below=of e2, align=center] (e3){\textbf{3. Análisis}\\ ideas secundarias + léxico (citas)};
\node[etape, below=of e3, align=center] (e4){\textbf{4. Opinión personal}\\ \esp{En mi opinión... Pienso que...}};
\draw[fleche] (e1) -- (e2);
\draw[fleche] (e2) -- (e3);
\draw[fleche] (e3) -- (e4);
\end{tikzpicture}
\legende{Les quatre étapes du commentaire de texte : présentation, idée principale, analyse, avis personnel.}
\end{popfigure}

\courssub{Étape 1 : présenter le texte}

La présentation répond à trois questions : quel est le \cle{type de texte} (article de presse, dialogue, lettre, poème), quel est le \cle{thème}, quelle est la \cle{source} ou la situation de communication ? Deux ou trois phrases suffisent : inutile de résumer tout le texte à ce stade.

\begin{exemplebox}[Présenter le texte : modèles de phrases]
\esp{El texto es un artículo de prensa.} — Le texte est un article de presse.\par
\esp{El texto trata de la importancia del deporte.} — Le texte traite de l'importance du sport.\par
\esp{El tema del texto es el deporte y el bienestar.} — Le thème du texte est le sport et le bien-être.\par
\esp{Según el autor, el deporte es una fuente de bienestar.} — Selon l'auteur, le sport est une source de bien-être.
\end{exemplebox}

\begin{attention}[Faux ami : tratar de]
\esp{Tratar de} + nom signifie « traiter de, parler de » : \esp{El texto trata del deporte.} — Le texte traite du sport. Mais \esp{tratar de} + infinitif signifie « essayer de » : \esp{Trato de practicar deporte cada día.} — J'essaie de pratiquer un sport chaque jour. Ne confondez pas les deux constructions dans votre commentaire. Attention aussi à la contraction obligatoire : \esp{de} + \esp{el} = \esp{del} (\esp{trata del deporte}).
\end{attention}

\courssub{Étape 2 : dégager l'idée principale}

L'\cle{idée principale} (\esp{idea principal}) est le message essentiel du texte, résumé en une ou deux phrases. Pour la trouver, demandez-vous : que veut démontrer l'auteur ? Évitez de recopier une phrase entière du texte : reformulez avec vos propres mots.

\begin{exemplebox}[Formuler l'idée principale]
\esp{La idea principal del texto es que el deporte es bueno para la salud física y mental.} — L'idée principale du texte est que le sport est bon pour la santé physique et mentale.\par
\esp{El autor quiere animar a los jóvenes a practicar deporte.} — L'auteur veut encourager les jeunes à pratiquer un sport.
\end{exemplebox}

\courssub{Étape 3 : analyser les idées secondaires et le lexique}

On relève ensuite les \cle{idées secondaires} (par exemple : les bienfaits physiques, les bienfaits mentaux, l'exemple du football au Cameroun) et on les justifie en \cle{citant le texte en espagnol} entre guillemets. On peut aussi commenter le champ lexical du sport : \esp{entrenamiento}, \esp{partido}, \esp{equipo}, \esp{aficionado}, \esp{campeón}. La citation est la preuve que vous avez réellement lu le texte : elle est indispensable.

\begin{exemplebox}[Analyser et citer]
\esp{El autor dice que el deporte «fortalece el corazón»: es un beneficio físico.} — L'auteur dit que le sport « fortifie le cœur » : c'est un bienfait physique.\par
\esp{También explica que el deporte «reduce el estrés»: es un beneficio mental.} — Il explique aussi que le sport « réduit le stress » : c'est un bienfait mental.\par
\esp{Además, el texto utiliza un léxico del deporte: «equipo», «partido», «entrenamiento».} — De plus, le texte utilise un lexique du sport : « équipe », « match », « entraînement ».
\end{exemplebox}

\courssub{Étape 4 : donner son avis personnel argumenté}

C'est l'étape la plus personnelle : vous dites si vous êtes d'accord ou non, et \cle{pourquoi}. Vous pouvez comparer la situation du Cameroun et celle des pays hispanophones. Un avis sans argument ne vaut rien : chaque \esp{pienso que} doit être suivi d'un \esp{porque}.

\begin{exemplebox}[Exprimer son avis]
\esp{En mi opinión, el deporte es necesario para los jóvenes.} — À mon avis, le sport est nécessaire pour les jeunes.\par
\esp{Pienso que el fútbol es muy importante en Camerún.} — Je pense que le football est très important au Cameroun.\par
\esp{Estoy de acuerdo porque el deporte une a las personas.} — Je suis d'accord parce que le sport unit les personnes.\par
\esp{Sin embargo, en mi barrio no hay suficientes campos de deporte.} — Cependant, dans mon quartier, il n'y a pas assez de terrains de sport.
\end{exemplebox}

\begin{center}
\begin{tabularx}{\linewidth}{|X|X|X|}
\hline
\rowcolor{popblueL} Français & Español & Étape du commentaire \\
\hline
Le texte traite de... & El texto trata de... & Présentation \\
\hline
Selon l'auteur... & Según el autor... & Présentation / analyse \\
\hline
L'idée principale est que... & La idea principal es que... & Idée principale \\
\hline
De plus / aussi & Además / También & Analyse \\
\hline
Cependant & Sin embargo & Nuance \\
\hline
C'est pourquoi & Por eso & Conclusion logique \\
\hline
À mon avis / je pense que & En mi opinión / Pienso que & Avis personnel \\
\hline
Je suis d'accord parce que... & Estoy de acuerdo porque... & Avis argumenté \\
\hline
\end{tabularx}
\end{center}

\begin{methode}[Réussir le commentaire de texte au Bac]
\begin{enumerate}
\item \textbf{Lire} le texte deux fois et souligner les mots du champ lexical du sport.
\item \textbf{Présenter} le texte en deux phrases : type, thème, source (\esp{El texto es un artículo... trata de...}).
\item \textbf{Dégager} l'idée principale en une phrase claire, sans recopier le texte.
\item \textbf{Analyser} deux ou trois idées secondaires en citant le texte en espagnol entre guillemets : citer le texte rapporte des points.
\item \textbf{Utiliser} les connecteurs : \esp{además}, \esp{también}, \esp{por eso}, \esp{sin embargo}.
\item \textbf{Terminer} toujours par un avis personnel argumenté avec \esp{En mi opinión...} ou \esp{Pienso que... porque...}.
\item \textbf{Relire} pour vérifier les accords, les accents et la conjugaison (\esp{me gusta}, \esp{suelo} + infinitif).
\end{enumerate}
\end{methode}

\courssub{Commentaire modèle rédigé}

Voici un commentaire complet (environ douze lignes) sur le texte support de la leçon. Observez la progression en quatre étapes et les connecteurs qui assurent la cohérence du texte.

\begin{textebox}[Comentario modelo : El deporte, una fuente de bienestar]
El texto es un artículo sobre el deporte y el bienestar. Trata de la importancia del deporte para la salud de los jóvenes, en Camerún y en los países hispanohablantes.\par
La idea principal es que el deporte es una fuente de bienestar físico y mental. Según el autor, el deporte «fortalece el corazón y los músculos»: es un beneficio físico. Además, el texto explica que el deporte «reduce el estrés y mejora el ánimo»: es un beneficio mental. También utiliza un léxico del deporte muy rico: «equipo», «partido», «entrenamiento», «aficionado», «campeón».\par
En mi opinión, el deporte es necesario para los jóvenes. Pienso que en Camerún el fútbol une a los barrios, como en España, donde los aficionados siguen a sus equipos cada semana. Sin embargo, muchos jóvenes no tienen campos para entrenar. Por eso, estoy de acuerdo con el autor: debemos practicar deporte con regularidad, porque es bueno para el cuerpo y para la mente.
\end{textebox}

\textbf{Glossaire :} \esp{fuente} = source ; \esp{fortalece} = fortifie ; \esp{músculos} = muscles ; \esp{reduce el estrés} = réduit le stress ; \esp{mejora el ánimo} = améliore le moral ; \esp{barrios} = quartiers ; \esp{siguen} = suivent, supportent ; \esp{campos} = terrains ; \esp{mente} = esprit.

\textbf{Questions de vérification :}
\begin{enumerate}
\item Quel est le type de texte commenté et quel est son thème ?
\item Quels sont les deux bienfaits du sport cités dans le commentaire ?
\item Quels connecteurs l'élève utilise-t-il pour organiser son analyse ?
\item Quelle comparaison fait-il entre le Cameroun et l'Espagne ?
\end{enumerate}

\begin{exoresolu}[Repérer les quatre étapes dans le commentaire modèle]
Énoncé : recopiez le tableau et indiquez à quelle étape (1 à 4) appartient chaque phrase extraite du commentaire modèle : a) \esp{En mi opinión, el deporte es necesario para los jóvenes.} b) \esp{El texto es un artículo sobre el deporte y el bienestar.} c) \esp{La idea principal es que el deporte es una fuente de bienestar.} d) \esp{El autor dice que el deporte «fortalece el corazón».}
\tcblower
Solution : a) étape 4, avis personnel (marqueur \esp{En mi opinión}) ; b) étape 1, présentation (type et thème du texte) ; c) étape 2, idée principale (marqueur \esp{La idea principal es que...}) ; d) étape 3, analyse, car la phrase cite le texte entre guillemets pour justifier une idée secondaire (le bienfait physique).
\end{exoresolu}

\begin{savaistu}[Le commentaire au Baccalauréat camerounais]
Au Bac, le commentaire de texte est noté sur trois critères : la \cle{compréhension} du texte, la \cle{justesse de la langue} (grammaire, orthographe, accents) et la \cle{pertinence de l'avis personnel}. Bonne nouvelle : citer le texte en espagnol pour justifier votre analyse rapporte des points. Un avis personnel sans argument, en revanche, ne suffit jamais : pensez toujours à ajouter \esp{porque...}.
\end{savaistu}

\courssub{Production guidée : Mi deporte favorito}

À votre tour de produire ! Rédigez un paragraphe de 6 à 8 lignes intitulé \esp{Mi deporte favorito} en suivant la trame imposée. Cette trame réutilise toute la grammaire de la leçon : \esp{gustar}, \esp{encantar} et \esp{soler} + infinitif.

\begin{methode}[Trame de la production « Mi deporte favorito »]
\begin{enumerate}
\item \textbf{Le sport pratiqué} : \esp{Mi deporte favorito es el fútbol / el atletismo / el baloncesto...}
\item \textbf{La fréquence} avec \esp{soler} + infinitif : \esp{Suelo entrenar tres veces por semana con mi equipo.}
\item \textbf{Le goût} avec \esp{gustar} / \esp{encantar} : \esp{Me gusta mucho jugar con mis amigos. Me encanta marcar goles.}
\item \textbf{Les bienfaits pour la santé} : \esp{El deporte es bueno para mi salud porque fortalece mi cuerpo y reduce el estrés.}
\item \textbf{Conclusion} : \esp{Por eso, pienso que todos los jóvenes deben practicar deporte.}
\end{enumerate}
\end{methode}

\begin{exemplebox}[Modèle de production rédigée]
\esp{Mi deporte favorito es el fútbol. Suelo jugar con mis amigos del barrio cuatro veces por semana, después de las clases. Me gusta mucho correr detrás del balón y me encanta marcar goles. Además, el fútbol es bueno para mi salud: fortalece mi corazón y mis piernas, y reduce el estrés de los exámenes. También me ayuda a tener amigos. Por eso, pienso que todos los jóvenes de Camerún deben practicar un deporte.}\par
Traduction : Mon sport favori est le football. J'ai l'habitude de jouer avec mes amis du quartier quatre fois par semaine, après les cours. J'aime beaucoup courir après le ballon et j'adore marquer des buts. De plus, le football est bon pour ma santé : il fortifie mon cœur et mes jambes, et il réduit le stress des examens. Il m'aide aussi à avoir des amis. C'est pourquoi je pense que tous les jeunes du Cameroun doivent pratiquer un sport.
\end{exemplebox}

\begin{attention}[Pièges des francophones dans la production]
\begin{itemize}
\item \esp{Soler} se conjugue comme \esp{volver} (o $\to$ ue) : \esp{yo suelo}, et non « yo solo » (\esp{solo} = seul !).
\item Après \esp{me gusta}, mettez l'infinitif : \esp{me gusta jugar}, jamais « me gusta juego ».
\item \esp{Hacer deporte} et \esp{practicar deporte} sont corrects ; ne traduisez pas « faire du sport » mot à mot.
\item N'oubliez pas les accents : \esp{fútbol}, \esp{opinión}, \esp{además}, \esp{corazón}.
\item « Par semaine » se dit \esp{por semana} (et non « para semana ») : \esp{entreno tres veces por semana}.
\end{itemize}
\end{attention}

\begin{exercice}[Rédaction guidée : Mi deporte favorito]
Rédigez votre propre paragraphe de 6 à 8 lignes en suivant la trame imposée (sport + fréquence avec \esp{soler} + goût avec \esp{gustar}/\esp{encantar} + bienfaits + conclusion). Utilisez au moins cinq mots du lexique de la leçon : \esp{deporte}, \esp{entrenamiento}, \esp{partido}, \esp{equipo}, \esp{aficionado}, \esp{campeón}, \esp{ocio}.
\end{exercice}

\begin{corrige}[Exercice 1 : grille d'auto-évaluation]
Relisez votre paragraphe et cochez mentalement chaque critère :
\begin{itemize}
\item Ai-je utilisé \esp{soler} + infinitif correctement conjugué (\esp{suelo} + infinitif) ?
\item Ai-je utilisé \esp{gustar} ou \esp{encantar} avec \esp{me} et un infinitif ou un nom singulier ?
\item Ai-je employé au moins cinq mots du lexique du sport et des loisirs ?
\item Ai-je cité au moins deux bienfaits du sport pour la santé (physique et mentale) ?
\item Ai-je utilisé au moins deux connecteurs (\esp{además}, \esp{también}, \esp{por eso}, \esp{sin embargo}) ?
\item Ai-je vérifié les accents (\esp{fútbol}, \esp{opinión}, \esp{corazón}) et la ponctuation espagnole (¿ ? ¡ !) ?
\end{itemize}
Si vous avez coché les six critères, votre production est prête pour le Bac !
\end{corrige}

\courssub{Débat d'encouragement : ¿Por qué debemos practicar deporte ?}

\begin{experience}[Débat de classe : animar a practicar deporte]
La classe se divise en deux équipes. À tour de rôle, chaque élève donne un argument pour encourager la pratique sportive, en commençant obligatoirement par \esp{Pienso que... porque...} ou \esp{En mi opinión... porque...}. Un élève ne peut pas répéter l'argument d'un camarade.\par
Exemples d'arguments attendus :\par
\esp{Pienso que debemos practicar deporte porque es bueno para el corazón.} — Je pense que nous devons pratiquer un sport parce que c'est bon pour le cœur.\par
\esp{En mi opinión, el deporte es importante porque reduce el estrés de los exámenes.} — À mon avis, le sport est important parce qu'il réduit le stress des examens.\par
\esp{Pienso que el deporte une a los jóvenes del barrio.} — Je pense que le sport unit les jeunes du quartier.\par
Le professeur note au tableau les arguments dans deux colonnes : \esp{beneficios físicos} et \esp{beneficios mentales y sociales}. L'équipe qui trouve le plus d'arguments différents gagne le titre de \esp{campeón del debate}.
\end{experience}

\begin{aretenir}[L'essentiel de la section]
\begin{itemize}
\item Le commentaire de texte suit quatre étapes : présentation, idée principale, analyse avec citations, avis personnel argumenté.
\item Les phrases-outils : \esp{El texto trata de...}, \esp{Según el autor...}, \esp{La idea principal es que...}, \esp{En mi opinión...}, \esp{Estoy de acuerdo porque...}.
\item Les connecteurs \esp{además}, \esp{también}, \esp{por eso}, \esp{sin embargo} donnent de la cohérence au texte.
\item La production guidée réutilise la grammaire de la leçon : \esp{suelo} + infinitif pour la fréquence, \esp{me gusta} / \esp{me encanta} pour le goût.
\item Citer le texte en espagnol et terminer par un avis argumenté : deux habitudes qui rapportent des points au Bac.
\end{itemize}
\end{aretenir}

\newpage

\coursec{Activité d'intégration}

\courssub{Objectifs de l'activité}

À l'issue de cette activité d'intégration, l'élève sera capable de :

\begin{itemize}
\item mobiliser le \cle{lexique du sport et du loisir} (deporte, equipo, entrenamiento, partido, aficionado, campeón...) dans une production authentique ;
\item exprimer ses \cle{goûts} avec \esp{gustar} et ses nuances (\esp{me gusta}, \esp{me encanta}, \esp{no me gusta nada}) ;
\item exprimer l'\cle{habitude et la fréquence} avec \esp{soler + infinitif} et les compléments de fréquence (\esp{dos veces por semana}, \esp{todos los días}) ;
\item construire une \cle{argumentation simple} pour convaincre : donner au moins deux bienfaits du sport sur la santé physique et mentale ;
\item travailler en équipe, présenter oralement un message publicitaire et interagir en espagnol avec les autres groupes.
\end{itemize}

\courssub{Situation}

Tu es élève en classe de Première A au lycée bilingue de Yaoundé. La semaine prochaine, ton établissement organise la \cle{Semaine du sport} : tournois de football et de handball, courses d'athlétisme, démonstrations de danse. Le club d'espagnol a été chargé de la communication : il faut motiver les élèves, souvent sédentaires, à s'inscrire et à bouger. Ton groupe a décidé de créer une \cle{affiche} ou un \cle{message radiophonique} de deux minutes pour la radio du lycée, « Radio Juventud ».

Pour t'inspirer, voici le message que la radio a diffusé l'année dernière et qui a eu beaucoup de succès :

\begin{textebox}[Radio Juventud — mensaje del año pasado]
¡Hola, chicos y chicas! Soy Marta, de la clase de Première A. ¿Sabéis qué? A mí me encanta el balonmano. Suelo entrenar con mi equipo tres veces por semana, los martes, los jueves y los sábados por la tarde, en el campo del barrio.

¿Por qué me gusta tanto? Primero, porque el deporte es bueno para el cuerpo: fortalece el corazón y los músculos. Segundo, porque es bueno para la mente: después de un partido, me siento tranquila y duermo mejor. Además, con mi equipo tengo nuevos amigos y aprendo a trabajar en grupo.

En España, muchos jóvenes practican deporte en clubes escolares; aquí en Camerún, el fútbol es el rey, ¡pero todos los deportes son importantes!

Entonces, ¿qué esperas? ¡Anímate a practicar deporte! Tu cuerpo y tu mente te lo van a agradecer.
\end{textebox}

\textbf{Glossaire :} \esp{balonmano} : handball ; \esp{el campo del barrio} : le terrain du quartier ; \esp{fortalece} : renforce ; \esp{me siento} : je me sens ; \esp{duermo mejor} : je dors mieux ; \esp{el rey} : le roi ; \esp{¿qué esperas?} : qu'attends-tu ? ; \esp{te lo van a agradecer} : ils t'en seront reconnaissants.

\courssub{Consignes}

Par groupes de 4, réalisez les tâches suivantes dans l'ordre.

\begin{enumerate}
\item \textbf{Compréhension.} Lisez le message de Marta à voix haute dans votre groupe. Relevez : a) le sport présenté ; b) la fréquence de l'entraînement ; c) les deux arguments sur la santé ; d) la comparaison entre l'Espagne et le Cameroun.

\item \textbf{Choix du sport.} Mettez-vous d'accord sur un sport à promouvoir (fútbol, atletismo, baloncesto, voleibol, balonmano, natación, danse traditionnelle...). Chaque membre du groupe doit dire une phrase avec \esp{gustar} : \esp{A mí me gusta... porque...}.

\item \textbf{Préparation du texte.} Rédigez ensemble un texte de 10 à 15 lignes qui contient obligatoirement :
\begin{itemize}
\item une accroche pour attirer l'attention (\esp{¡Hola! ¿Sabéis qué?}) ;
\item la présentation du sport et de votre équipe ou club ;
\item au moins deux phrases avec \esp{gustar} ou \esp{encantar} correctement construites ;
\item au moins une phrase avec \esp{soler + infinitif} et une expression de fréquence (\esp{dos veces por semana}, \esp{todos los fines de semana}) ;
\item deux arguments sur les bienfaits pour la santé physique et la santé mentale ;
\item si possible, une comparaison entre le Cameroun et un pays hispanophone ;
\item un slogan final impératif : \esp{¡Anímate a practicar deporte!}.
\end{itemize}

\item \textbf{Réalisation.} Fabriquez l'affiche (titre, slogan, dessins, 5 mots de lexique minimum bien visibles) ou répartissez les rôles pour le message radiophonique (présentateur, deux témoignages, conclusion).

\item \textbf{Présentation.} Chaque groupe présente son affiche ou joue son message devant la classe en deux minutes maximum. Parlez fort, articulez, regardez le public.

\item \textbf{Interaction.} Après chaque présentation, les autres groupes posent deux questions en espagnol avec les structures de la leçon : \esp{¿Te gusta el fútbol?}, \esp{¿Sueles entrenar por la mañana?}, \esp{¿Cuántas veces por semana practicáis?}. Le groupe interrogé répond en phrases complètes.
\end{enumerate}

\begin{attention}[Critères d'évaluation de l'enseignant]
L'enseignant évalue chaque groupe avec une grille : lexique du sport (5 mots minimum employés correctement), emploi correct de \esp{gustar} (avec \esp{a mí me}, l'article et le bon accord), emploi correct de \esp{soler + infinitif} conjugué, qualité des deux arguments de santé, fluidité et prononciation. Attention aux erreurs typiques : ne dites jamais \esp{yo gusto}, dites \esp{me gusta} ; après \esp{soler}, mettez toujours un infinitif : \esp{suelo practicar}, jamais \esp{suelo practico}.
\end{attention}

\courssub{Ressources à mobiliser}

\begin{aretenir}[Boîte à outils de la campagne]
\begin{itemize}
\item \textbf{Lexique :} el deporte, el fútbol, el atletismo, el equipo, el entrenamiento, el partido, el aficionado, el campeón, la salud, el ocio, el bienestar.
\item \textbf{Goût :} \esp{A mí me gusta el atletismo.} « J'aime l'athlétisme. » ; \esp{Me encanta jugar al fútbol.} « J'adore jouer au football. » ; \esp{¿Te gusta el deporte?} « Aimes-tu le sport ? »
\item \textbf{Habitude :} \esp{Suelo entrenar tres veces por semana.} « J'ai l'habitude de m'entraîner trois fois par semaine. » Rappel : soler est irrégulier (o $\to$ ue) : suelo, sueles, suele, solemos, soléis, suelen.
\item \textbf{Fréquence :} todos los días, una vez / dos veces por semana, los fines de semana, por la mañana / por la tarde.
\item \textbf{Argumenter :} \esp{primero... segundo... además...} ; \esp{el deporte es bueno para la salud} ; \esp{fortalece el corazón} ; \esp{reduce el estrés} ; \esp{ayuda a dormir mejor}.
\item \textbf{Encourager :} \esp{¡Anímate!}, \esp{¡Ven con nosotros!}, \esp{¿Qué esperas?}
\end{itemize}
\end{aretenir}

\courssub{Proposition de production et corrigé commenté}

\begin{exoresolu}[Message radiophonique modèle — Grupo « Los Leones »]
Rédigez le message de votre groupe pour Radio Juventud en respectant toutes les consignes.
\tcblower
\textbf{Production modèle :}

\esp{¡Hola, amigos de Radio Juventud! Somos « Los Leones », un grupo de Première A, y hoy queremos hablaros del atletismo.}

\esp{A nosotros nos encanta correr. Yo suelo correr por la mañana, antes de las clases, tres veces por semana, y mi compañera Aïcha suele entrenar con el club del liceo los sábados.}

\esp{¿Por qué nos gusta el atletismo? Primero, porque correr fortalece el corazón y los músculos: es muy bueno para la salud física. Segundo, porque reduce el estrés de los exámenes: después de correr, estamos más tranquilos y trabajamos mejor. Es bueno para la salud mental.}

\esp{En Guinea Ecuatorial y en España, muchos jóvenes practican atletismo en la escuela; aquí en Camerún también tenemos grandes campeones. ¡Nosotros podemos ser los próximos!}

\esp{Entonces, ¿qué esperas? ¡Anímate a practicar deporte! ¡Tu cuerpo y tu mente te lo van a agradecer!}

\textbf{Commentaire des choix de langue :}
\begin{itemize}
\item \esp{nos encanta correr} : avec « nous », gustar s'accorde avec l'infinitif \esp{correr} (singulier) et le pronom est \esp{nos} ; le renforcement \esp{a nosotros} est correct.
\item \esp{suelo correr}, \esp{suele entrenar} : soler est bien conjugué (o $\to$ ue) et suivi d'un infinitif, avec des compléments de fréquence variés.
\item Les arguments sont organisés par \esp{primero} et \esp{segundo}, et distinguent clairement santé physique (\esp{fortalece el corazón}) et santé mentale (\esp{reduce el estrés}).
\item La comparaison culturelle (\esp{Guinea Ecuatorial}, \esp{España}, \esp{Camerún}) valorise le contexte hispanophone proche du Cameroun.
\item Le slogan final reprend l'impératif \esp{¡Anímate!} avec le pronom \esp{te} soudé, comme dans l'affiche officielle.
\end{itemize}
\end{exoresolu}

\courssub{Bilan de l'activité}

Cette campagne « ¡Anímate a practicar deporte ! » t'a permis de réutiliser dans une situation réelle tout le contenu de la leçon : le lexique du sport et du loisir, l'expression du goût avec \esp{gustar}, l'habitude avec \esp{soler + infinitif} et les compléments de fréquence. Tu as aussi appris à construire une argumentation simple en deux points (santé physique, santé mentale) et à interagir à l'oral en posant des questions avec \esp{¿Te gusta...?} et \esp{¿Sueles...?}. Enfin, en comparant le Cameroun, l'Espagne et la Guinée équatoriale, tu as découvert que le sport est une langue universelle du monde hispanophone. Retiens le slogan : \esp{El deporte es vida} — « Le sport, c'est la vie ».

\newpage

\coursec{Méthodes et exercices résolus}

\courssub{Méthode 1 : Lire et comprendre un texte sur le sport}

\begin{methode}[Réussir les questions de compréhension sur un texte sportif]
\begin{enumerate}
\item \textbf{Lire d'abord les questions}, puis le texte : tu sais ainsi quelles informations chercher (personnes, lieux, fréquence, bénéfices pour la santé).
\item \textbf{Repérer le lexique du champ sportif} : \esp{deporte, entrenamiento, partido, equipo, campeón, aficionado}. Souligne chaque mot-clé et note à côté son équivalent français.
\item \textbf{Localiser la réponse dans le texte} : chaque réponse doit s'appuyer sur un passage précis. Recopie la structure de la phrase du texte en changeant le point de vue (\esp{Los jóvenes practican...} $\rightarrow$ \esp{Practican...}).
\item \textbf{Répondre par une phrase complète en espagnol}, jamais par un mot isolé : \esp{¿Dónde entrena el equipo?} $\rightarrow$ \esp{El equipo entrena en el estadio de Yaoundé.}
\item \textbf{Pour les questions d'opinion} (\esp{¿Crees que...?}), utilise les structures du goût : \esp{Me gusta..., creo que el deporte es importante porque...}.
\item \textbf{Vérifier} : accord sujet-verbe, accents écrits, ponctuation espagnole (¿...? ¡...!).
\end{enumerate}
\end{methode}

\begin{exoresolu}[Compréhension de l'écrit — type Bac]
\textbf{Texte.} \esp{En el barrio de Mvog-Ada, en Yaoundé, un grupo de jóvenes se reúne cada tarde para jugar al fútbol. «El deporte nos mantiene sanos y nos ayuda a olvidar los problemas», explica Junior, el capitán del equipo. Además, el médico del barrio afirma que los jóvenes que hacen deporte duermen mejor y trabajan con más concentración en la escuela.}

\textbf{Énoncé.} Réponds en espagnol par des phrases complètes : 1) ¿Dónde se reúnen los jóvenes? 2) ¿Qué hacen cada tarde? 3) Según Junior, ¿qué beneficios tiene el deporte? 4) ¿Qué afirma el médico del barrio?
\tcblower
\textbf{Raisonnement.} Chaque question commence par un mot interrogatif qui indique l'information à chercher : \esp{dónde} (lieu), \esp{qué} (action), \esp{según Junior} (opinion du capitán), \esp{el médico} (parole du médecin). On reprend la structure de la phrase du texte en la reformulant.

\textbf{Réponses rédigées.}
\begin{enumerate}
\item \esp{Los jóvenes se reúnen en el barrio de Mvog-Ada, en Yaoundé.}
\item \esp{Cada tarde juegan al fútbol.} (Attention : \esp{jugar} est un verbe à changement radical o $\rightarrow$ ue : \esp{juegan}.)
\item \esp{Según Junior, el deporte los mantiene sanos y los ayuda a olvidar los problemas.}
\item \esp{El médico afirma que los jóvenes que hacen deporte duermen mejor y trabajan con más concentración en la escuela.}
\end{enumerate}

\textbf{Conclusion.} Les réponses reprennent le vocabulaire du texte (\esp{reunirse, jugar, mantenerse sano, dormir mejor}) sans le copier mot à mot : c'est exactement ce que le correcteur attend.
\end{exoresolu}

\courssub{Méthode 2 : Parler de ses loisirs et de ses goûts avec \esp{gustar}}

\begin{methode}[Construire une phrase avec \esp{gustar}]
\begin{enumerate}
\item \textbf{Retenir la logique inversée} : en espagnol, on ne dit pas « j'aime le sport » mais « le sport me plaît » : \esp{Me gusta el deporte.} Le sujet grammatical est la chose aimée.
\item \textbf{Choisir le pronom d'objet indirect} : \esp{me, te, le, nos, os, les}.
\item \textbf{Accorder le verbe avec ce qui plaît} : singulier $\rightarrow$ \esp{gusta} ; pluriel $\rightarrow$ \esp{gustan}. \esp{Me gusta el fútbol.} / \esp{Me gustan los deportes de equipo.}
\item \textbf{Devant un infinitif, toujours le singulier} : \esp{Me gusta jugar al fútbol.} « J'aime jouer au football. »
\item \textbf{Pour insister ou préciser la personne}, ajouter \esp{a + pronom/nom} : \esp{A mí me gusta el atletismo. A María le gusta nadar.}
\item \textbf{La négation} se place devant le pronom : \esp{No me gusta el boxeo.}
\end{enumerate}
\end{methode}

\begin{aretenir}[Les pronoms de \esp{gustar}]
\begin{center}
\begin{tabular}{|l|l|l|}
\hline
\rowcolor{popblueL} Personne & Pronom & Exemple \\
\hline
je & \esp{me} & \esp{Me gusta el fútbol.} \\
\hline
tu & \esp{te} & \esp{¿Te gusta nadar?} \\
\hline
il/elle & \esp{le} & \esp{A María le gusta correr.} \\
\hline
nous & \esp{nos} & \esp{Nos gustan los partidos.} \\
\hline
vous & \esp{os} & \esp{Os gusta el atletismo.} \\
\hline
ils/elles & \esp{les} & \esp{Les gustan los videojuegos.} \\
\hline
\end{tabular}
\end{center}
Rappel : le verbe s'accorde avec la chose aimée, jamais avec la personne.
\end{aretenir}

\begin{exoresolu}[Expression des goûts — transformation de phrases]
\textbf{Énoncé.} Transforme les phrases suivantes avec \esp{gustar} selon le modèle. Modèle : \esp{Yo amo el baloncesto.} $\rightarrow$ \esp{Me gusta el baloncesto.}
1) \esp{Yo amo los partidos de fútbol.} 2) \esp{Ella ama correr por la mañana.} 3) \esp{Nosotros amamos el atletismo.} 4) \esp{Ellos no aman el boxeo.} 5) Traduis : « Mes amis aiment beaucoup les jeux vidéo, mais moi je préfère le sport. »
\tcblower
\textbf{Raisonnement.} On inverse la construction : la chose aimée devient sujet. On choisit \esp{gusta} ou \esp{gustan} selon le nombre, et le pronom selon la personne.

\textbf{Solutions.}
\begin{enumerate}
\item \esp{Me gustan los partidos de fútbol.} (pluriel $\rightarrow$ \esp{gustan})
\item \esp{Le gusta correr por la mañana.} (infinitif $\rightarrow$ singulier \esp{gusta})
\item \esp{Nos gusta el atletismo.}
\item \esp{No les gusta el boxeo.} (la négation précède le pronom)
\item \esp{A mis amigos les gustan mucho los videojuegos, pero a mí me gusta más el deporte.} (ou \esp{...pero yo prefiero el deporte.}) Justification : \esp{les gustan} car \esp{los videojuegos} est pluriel ; \esp{a mí} avec accent pour insister sur l'opposition.
\end{enumerate}

\textbf{Conclusion.} La question à se poser est toujours : « qu'est-ce qui plaît ? » — c'est ce mot qui commande \esp{gusta} ou \esp{gustan}.
\end{exoresolu}

\courssub{Méthode 3 : Exprimer l'habitude et la fréquence avec \esp{soler} + infinitif}

\begin{methode}[Employer \esp{soler} + infinitif et les marqueurs de fréquence]
\begin{enumerate}
\item \textbf{Sens} : \esp{soler + infinitif} = « avoir l'habitude de, faire habituellement ». \esp{Suelo entrenar los sábados.} « J'ai l'habitude de m'entraîner le samedi. »
\item \textbf{Conjugaison au présent} (verbe à changement radical o $\rightarrow$ ue) : \esp{suelo, sueles, suele, solemos, soléis, suelen}.
\item \textbf{L'infinitif suit directement}, sans préposition : \esp{Solemos jugar al fútbol después de clase.} « Nous jouons habituellement au football après les cours. »
\item \textbf{Combiner avec les adverbes de fréquence} pour nuancer : \esp{siempre} (toujours), \esp{normalmente} (normalement), \esp{a menudo} (souvent), \esp{a veces} (parfois), \esp{casi nunca} (presque jamais), \esp{nunca} (jamais).
\item \textbf{Expressions de temps utiles} : \esp{cada día} (chaque jour), \esp{una vez a la semana} (une fois par semaine), \esp{dos veces al mes} (deux fois par mois).
\item \textbf{Au passé} (habitude révolue), utilise l'imparfait : \esp{De niño, solía nadar en el río.} « Enfant, j'avais l'habitude de nager dans la rivière. »
\end{enumerate}
\end{methode}

\begin{propriete}[Conjugaison de \esp{soler} au présent de l'indicatif]
\begin{center}
\begin{tabular}{|l|l|}
\hline
\rowcolor{popblueL} Personne & Forme \\
\hline
\esp{yo} & \esp{suelo} \\
\hline
\esp{tú} & \esp{sueles} \\
\hline
\esp{él/ella/usted} & \esp{suele} \\
\hline
\esp{nosotros/nosotras} & \esp{solemos} \\
\hline
\esp{vosotros/vosotras} & \esp{soléis} \\
\hline
\esp{ellos/ellas/ustedes} & \esp{suelen} \\
\hline
\end{tabular}
\end{center}
Règle : le changement o $\rightarrow$ ue a lieu quand l'accent tonique tombe sur la racine (\esp{suelo, sueles, suele, suelen}) ; il n'a pas lieu quand l'accent tombe sur la désinence (\esp{solemos, soléis}).
\end{propriete}

\begin{exoresolu}[Thème — l'habitude et la fréquence]
\textbf{Énoncé.} Traduis en espagnol. 1) J'ai l'habitude de faire du sport trois fois par semaine. 2) Nous jouons habituellement au football après les cours. 3) Ma sœur court parfois le matin avant l'école. 4) Les champions s'entraînent chaque jour. 5) Quand j'étais enfant, je regardais souvent les matchs à la télévision.
\tcblower
\textbf{Raisonnement.} « Avoir l'habitude de » au présent $\rightarrow$ \esp{soler + infinitif} au présent. « Faire du sport » se dit \esp{hacer deporte} (jamais \esp{hacer sport}). « S'entraîner » = \esp{entrenarse}. Pour la phrase 5, habitude passée $\rightarrow$ imparfait.

\textbf{Solutions.}
\begin{enumerate}
\item \esp{Suelo hacer deporte tres veces a la semana.} (\esp{soler} : o $\rightarrow$ ue, \esp{suelo} ; « trois fois par semaine » = \esp{tres veces a la semana}.)
\item \esp{Solemos jugar al fútbol después de las clases.} (\esp{jugar al fútbol} : à + el = \esp{al}.)
\item \esp{Mi hermana a veces corre por la mañana antes de la escuela.} (ou \esp{Mi hermana suele correr...})
\item \esp{Los campeones entrenan cada día.} (ou \esp{se entrenan} ; \esp{cada} est invariable.)
\item \esp{Cuando era niño, solía ver los partidos en la televisión.} (\esp{era} et \esp{solía} : imparfait de l'habitude ; \esp{ver los partidos} = « regarder les matchs ».)
\end{enumerate}

\textbf{Conclusion.} Deux réflexes : \esp{soler} se conjugue avec le changement o $\rightarrow$ ue, et l'habitude passée exige l'imparfait.
\end{exoresolu}

\courssub{Méthode 4 : Traduire une version sur le sport et le bien-être}

\begin{methode}[Réussir la version (espagnol $\rightarrow$ français)]
\begin{enumerate}
\item \textbf{Lire tout le texte une première fois} pour dégager le sens global (qui ? quoi ? où ? pourquoi ?).
\item \textbf{Identifier les temps} : présent d'habitude, imparfait, futur ? Le français doit rendre le même temps.
\item \textbf{Traduire le lexique sportif avec précision} : \esp{entrenamiento} = « entraînement » ; \esp{partido} = « match » ; \esp{aficionado} = « supporter, amateur » ; \esp{equipo} = « équipe » ; \esp{ocio} = « loisirs ».
\item \textbf{Se méfier des faux amis} : \esp{actualmente} = « actuellement » (et non « effectivement, en réalité ») ; \esp{asistir a un partido} = « assister à un match » (être présent, pas aider).
\item \textbf{Rendre les tournures espagnoles par des tournures françaises naturelles} : \esp{el deporte nos mantiene sanos} = « le sport nous maintient en bonne santé ».
\item \textbf{Relire la traduction française} : elle doit être correcte et fluide, sans calque de l'espagnol.
\end{enumerate}
\end{methode}

\begin{exoresolu}[Version — type Bac]
\textbf{Énoncé.} Traduis en français le texte suivant.

\esp{En muchos países hispanohablantes, el fútbol es más que un deporte: es una pasión. Los aficionados asisten a los partidos cada fin de semana y apoyan a su equipo con entusiasmo. Pero el deporte no es solo espectáculo: los médicos recuerdan que practicar una actividad física con regularidad mejora la salud del cuerpo y de la mente.}
\tcblower
\textbf{Raisonnement.} Texte au présent de vérité générale $\rightarrow$ présent en français. Points de vigilance : \esp{más que} = « plus que » ; \esp{asisten a} = « assistent à » (faux ami : pas « aident ») ; \esp{apoyan} = « soutiennent, encouragent » ; \esp{con regularidad} = « régulièrement » ; \esp{la salud del cuerpo y de la mente} = « la santé du corps et de l'esprit ».

\textbf{Traduction.} « Dans de nombreux pays hispanophones, le football est plus qu'un sport : c'est une passion. Les supporters assistent aux matchs chaque week-end et encouragent leur équipe avec enthousiasme. Mais le sport n'est pas seulement un spectacle : les médecins rappellent que pratiquer une activité physique régulièrement améliore la santé du corps et de l'esprit. »

\textbf{Conclusion.} La version exige à la fois la fidélité au sens et un français naturel : on traduit des idées, pas des mots isolés.
\end{exoresolu}

\courssub{Méthode 5 : Encourager la pratique sportive — production guidée}

\begin{methode}[Rédiger un court texte pour encourager le sport]
\begin{enumerate}
\item \textbf{Accroche} : une question ou un constat. \esp{¿Sabes que el deporte puede cambiar tu vida?}
\item \textbf{Arguments santé physique} : \esp{fortalece el corazón, mantiene el cuerpo sano, ayuda a dormir mejor}.
\item \textbf{Arguments santé mentale} : \esp{reduce el estrés, mejora la concentración, nos hace más felices}.
\item \textbf{Arguments sociaux} : \esp{permite hacer amigos, enseña el trabajo en equipo}.
\item \textbf{Conclusion impérative} : utilise l'impératif affirmatif. \esp{¡Haz deporte! ¡Practica con tus amigos! ¡Muévete!} (irréguliers : \esp{hacer} $\rightarrow$ \esp{haz} ; \esp{ir} $\rightarrow$ \esp{ve} ; \esp{ser} $\rightarrow$ \esp{sé} ; \esp{venir} $\rightarrow$ \esp{ven}.)
\item \textbf{Longueur Bac} : 60 à 100 mots, phrases simples et correctes plutôt que longues et fautives.
\end{enumerate}
\end{methode}

\begin{exoresolu}[Production écrite guidée — type Bac]
\textbf{Énoncé.} Tu es membre du club sportif de ton lycée à Douala. Rédige un court message (60-80 mots) pour encourager tes camarades à faire du sport. Utilise au moins deux verbes à l'impératif, une structure avec \esp{gustar} et une avec \esp{soler}.
\tcblower
\textbf{Raisonnement.} Plan : accroche (question) $\rightarrow$ deux arguments (santé physique, santé mentale) $\rightarrow$ goût (\esp{te gustará}) $\rightarrow$ habitude (\esp{soler}) $\rightarrow$ conclusion à l'impératif (\esp{ven, haz, sé}).

\textbf{Production modèle.}

\esp{¿Quieres estar sano y feliz? El deporte es la solución. Practicar una actividad física fortalece el cuerpo y reduce el estrés después de las clases. Además, te gustará hacer nuevos amigos en el equipo. En nuestro club, solemos entrenar los martes y los jueves por la tarde. ¡Ven al estadio, haz deporte con nosotros y sé campeón de tu propia vida!}

\textbf{Analyse grammaticale.} \esp{te gustará} : futur de \esp{gustar} avec pronom \esp{te} ; \esp{solemos entrenar} : \esp{soler} à la première personne du pluriel + infinitif ; \esp{ven} (impératif de \esp{venir}), \esp{haz} (\esp{hacer}), \esp{sé} (\esp{ser}) : trois impératifs irréguliers corrects ; \esp{por la tarde} = « l'après-midi ».

\textbf{Conclusion.} 70 mots, toutes les contraintes respectées, aucune faute : c'est le modèle à reproduire.
\end{exoresolu}

\begin{savaistu}[Le sport dans le monde hispanophone]
Le football est roi en Espagne et en Amérique latine : les grands clubs comme le Real Madrid ou le FC Barcelone comptent des millions de supporters (\esp{aficionados}) sur tous les continents, et de nombreux joueurs africains, dont des Camerounais, y ont fait carrière. Mais le monde hispanophone brille aussi dans d'autres disciplines : l'athlétisme (\esp{atletismo}) au Kenya hispanophone... non, plutôt en Guinée équatoriale, seul pays hispanophone d'Afrique, où le sport scolaire se développe ; le cyclisme en Colombie ; le baseball à Cuba et au Venezuela. Au Cameroun, le football reste le sport le plus pratiqué : chaque quartier de Yaoundé ou de Douala a son équipe et son terrain, et les matchs du dimanche rassemblent tout le voisinage.
\end{savaistu}

\begin{attention}[Les 5 erreurs qui coûtent des points au Bac]
\begin{enumerate}
\item \textbf{Oublier l'accent sur les mots interrogatifs} : \esp{¿Qué deporte practicas? ¿Dónde entrenas?} — \esp{que, donde} sans accent changent de sens (ce sont alors des relatifs).
\item \textbf{Accorder \esp{gustar} avec la mauvaise personne} : on écrit \esp{Me gustan los deportes} (pluriel) et non \esp{Me gusta los deportes}. C'est la chose aimée qui commande le verbe.
\item \textbf{Calquer le français} : « faire du sport » se dit \esp{hacer deporte} (sans article, sans « du ») ; « jouer au football » = \esp{jugar al fútbol} (avec \esp{al} obligatoire, jamais \esp{jugar el fútbol}).
\item \textbf{Conjuguer \esp{soler} sans le changement radical} : \esp{suelo, sueles, suele, suelen} (o $\rightarrow$ ue), mais \esp{solemos, soléis} (l'accent tonique tombe sur la désinence, donc pas de diphtongaison).
\item \textbf{Confondre les faux amis} : \esp{asistir a un partido} = « assister à un match » (et non « aider ») ; \esp{un deporte} reste masculin : \esp{el deporte, este deporte} ; \esp{actualmente} = « actuellement », pas « effectivement ».
\end{enumerate}
\end{attention}

\newpage

\coursec{Exercices}

\courssub{Je vérifie mes connaissances}

\begin{exercice}[QCM : le lexique du sport]
Entoure la bonne réponse.
\begin{enumerate}
\item Un \esp{partido} est : a) un entraînement \quad b) un match \quad c) un trophée
\item Un \esp{aficionado} est : a) un supporter \quad b) un arbitre \quad c) un champion
\item \esp{El entrenamiento} signifie : a) la compétition \quad b) l'entraînement \quad c) le repos
\item Un \esp{campeón} est celui qui : a) perd toujours \quad b) gagne une compétition \quad c) regarde le match
\item \esp{El ocio} signifie : a) le travail \quad b) les loisirs \quad c) la fatigue
\end{enumerate}
\end{exercice}

\begin{exercice}[Vrai ou faux ? Justifie ta réponse]
Réponds par \esp{Verdadero} ou \esp{Falso} et justifie en une phrase en français.
\begin{enumerate}
\item Le verbe \esp{gustar} se construit comme le verbe français « aimer » : \esp{Yo gusto el fútbol}.
\item \esp{Soler} est toujours suivi d'un infinitif.
\item \esp{Suelo entrenar los lunes} signifie « je m'entraînerai lundi ».
\item \esp{Me gustan los deportes} : on utilise \esp{gustan} parce que \esp{deportes} est pluriel.
\end{enumerate}
\end{exercice}

\begin{exercice}[Vocabulaire : trouve le mot]
Complète chaque phrase avec un mot du lexique de la leçon : \esp{equipo}, \esp{atletismo}, \esp{partido}, \esp{aficionados}, \esp{campeón}, \esp{ocio}.
\begin{enumerate}
\item Los \ldots\ldots\ldots\ldots animan a su \ldots\ldots\ldots\ldots en el estadio de Yaoundé.
\item El sábado hay un \ldots\ldots\ldots\ldots de fútbol entre los dos liceos.
\item Samuel Eto'o fue \ldots\ldots\ldots\ldots de África con los Leones Indomables.
\item El \ldots\ldots\ldots\ldots incluye la carrera, el salto y el lanzamiento.
\item El deporte es mi actividad de \ldots\ldots\ldots\ldots favorita.
\end{enumerate}
\end{exercice}

\begin{exercice}[Conjugaison : gustar et soler]
Complète avec la forme correcte du verbe entre parenthèses.
\begin{enumerate}
\item A mí me \ldots\ldots\ldots\ldots (gustar) mucho el baloncesto.
\item A mis amigos les \ldots\ldots\ldots\ldots (gustar) los partidos de fútbol.
\item ¿A ti te \ldots\ldots\ldots\ldots (gustar) nadar ?
\item Nosotros \ldots\ldots\ldots\ldots (soler) entrenar los martes.
\item Ella \ldots\ldots\ldots\ldots (soler) correr por la mañana.
\item ¿Vosotros \ldots\ldots\ldots\ldots (soler) ver los partidos en la televisión ?
\end{enumerate}
\end{exercice}

\courssub{Je m'entraîne}

\begin{exercice}[Texte à trous : gustar et soler]
Conjugue les dix verbes entre parenthèses au présent de l'indicatif.

\esp{En mi familia, todos amamos el deporte. A mí me (1. gustar) \ldots\ldots\ldots\ldots mucho el fútbol y a mi hermana le (2. gustar) \ldots\ldots\ldots\ldots el atletismo. Yo (3. soler) \ldots\ldots\ldots\ldots entrenar tres veces por semana con mi equipo. Los fines de semana, nosotros (4. soler) \ldots\ldots\ldots\ldots jugar un partido en el barrio. A mis padres les (5. gustar) \ldots\ldots\ldots\ldots caminar por la mañana : ellos (6. soler) \ldots\ldots\ldots\ldots levantarse muy temprano. ¿Y tú? ¿Te (7. gustar) \ldots\ldots\ldots\ldots los deportes de equipo? ¿(8. Soler) \ldots\ldots\ldots\ldots hacer ejercicio? A mis amigos y a mí nos (9. gustar) \ldots\ldots\ldots\ldots también los videojuegos, pero nosotros no (10. soler) \ldots\ldots\ldots\ldots jugar durante la semana.}
\end{exercice}

\begin{exercice}[Thème : traduis en espagnol]
Traduis les phrases suivantes en espagnol.
\begin{enumerate}
\item J'aime beaucoup le sport.
\item Mon équipe s'entraîne deux fois par semaine.
\item D'habitude, nous jouons un match le samedi.
\item Mon ami n'aime pas du tout perdre.
\item Le sport est bon pour la santé mentale.
\end{enumerate}
\end{exercice}

\begin{exercice}[Reformulation : l'habitude avec soler]
Transforme chaque phrase en utilisant \esp{soler} + infinitif, puis traduis la phrase obtenue en français.

Exemple : \esp{Normalmente corro por la mañana.} $\rightarrow$ \esp{Suelo correr por la mañana.}
\begin{enumerate}
\item Normalmente jugamos al fútbol los sábados.
\item De costumbre, mi hermana nada en la piscina los domingos.
\item Generalmente, los aficionados ven el partido en la televisión.
\item Habitualmente hago ejercicio después de las clases.
\item Los lunes, mi equipo entrena en el estadio, como siempre.
\end{enumerate}
\end{exercice}

\begin{exercice}[Appariements : mots et définitions]
Relie chaque mot espagnol à sa définition française.

\begin{center}
\begin{tabularx}{\linewidth}{|X|X|}
\hline
\rowcolor{popblueL} \textbf{Español} & \textbf{Français} \\
\hline
1. el entrenamiento & a) personne qui adore un sport ou une équipe \\
\hline
2. el aficionado & b) période de préparation physique avant la compétition \\
\hline
3. el campeón & c) rencontre sportive entre deux équipes \\
\hline
4. el partido & d) celui qui remporte la compétition \\
\hline
5. el ocio & e) temps libre consacré au repos et aux loisirs \\
\hline
6. el equipo & f) groupe de joueurs qui jouent ensemble \\
\hline
\end{tabularx}
\end{center}
\end{exercice}

\courssub{Je me prépare au Bac}

\begin{textebox}[El fútbol femenino crece en España]
Hace veinte años, pocas niñas españolas jugaban al fútbol. El fútbol era un deporte « para chicos » y las niñas practicaban otros deportes como el baloncesto o la gimnasia. Hoy, la situación es muy diferente : miles de niñas y adolescentes entrenan cada semana en clubes de todo el país.

Este cambio tiene varias razones. Primero, las jugadoras españolas han ganado títulos importantes y ahora son modelos para las jóvenes. Segundo, la televisión emite cada vez más partidos femeninos, y los aficionados suelen verlos con entusiasmo. Tercero, muchas escuelas han creado equipos de fútbol para las alumnas.

A Lucía, una jugadora de quince años de Madrid, le gusta mucho su deporte : « Suelo entrenar cuatro veces por semana. Me gustan los entrenamientos y me encantan los partidos del sábado. El fútbol me da amigas, salud y alegría ».

Sin embargo, quedan problemas : las jugadoras profesionales ganan menos dinero que los hombres y algunos clubes no tienen buenos campos para las chicas. Pero las jóvenes como Lucía no se rinden : sueñan con ser campeonas del mundo.
\end{textebox}

\textbf{Glossaire} : \esp{pocas} « peu de » ; \esp{miles de} « des milliers de » ; \esp{emitir} « diffuser » ; \esp{quedar} « rester » ; \esp{el campo} « le terrain » ; \esp{rendirse} « abandonner, se rendre » ; \esp{soñar con} « rêver de ».

\begin{exercice}[Compréhension de l'écrit : le football féminin en Espagne]
Lis le texte ci-dessus, puis réponds aux questions.

\textbf{A. Comprensión} — Réponds par \esp{Verdadero} ou \esp{Falso} et justifie par une citation du texte.
\begin{enumerate}
\item Il y a vingt ans, beaucoup de filles espagnoles jouaient au football.
\item La télévision diffuse de plus en plus de matchs féminins.
\item Lucía s'entraîne seulement le samedi.
\item Tous les problèmes du football féminin sont résolus.
\end{enumerate}

\textbf{B. Preguntas abiertas} — Réponds en espagnol avec des phrases complètes.
\begin{enumerate}
\item ¿Cuáles son las tres razones del cambio en el fútbol femenino español ?
\item ¿Qué le gusta a Lucía de su deporte ?
\item ¿Qué problemas tienen todavía las jugadoras profesionales ?
\end{enumerate}

\textbf{C. Lengua}
\begin{enumerate}
\item Trouve dans le texte deux verbes à l'imparfait et explique cet emploi.
\item Relève deux exemples du verbe \esp{gustar} et deux exemples du verbe \esp{soler}, puis traduis les phrases en français.
\item Trouve dans le texte l'équivalent espagnol de : « les supporters », « l'entraînement », « championnes ».
\end{enumerate}
\end{exercice}

\begin{textebox}[Las costumbres de Diego]
Diego es un adolescente mexicano de dieciséis años. Le gustan mucho los deportes, especialmente el fútbol. Suele entrenar con su equipo tres veces por semana, después de las clases. Normalmente, los sábados juega un partido con sus amigos en el parque. A Diego no le gusta nada perder, pero siempre respeta al otro equipo.
\end{textebox}

\begin{exercice}[Traduction : thème et version]
\textbf{A. Versión} — Traduis en français le texte \esp{Las costumbres de Diego} ci-dessus (5 phrases).

\textbf{B. Thème} — Traduis en espagnol les phrases suivantes.
\begin{enumerate}
\item J'aime beaucoup le sport.
\item Mon équipe s'entraîne deux fois par semaine.
\item D'habitude, nous jouons un match le samedi.
\item Mon ami n'aime pas du tout perdre.
\item Le sport est bon pour la santé mentale.
\end{enumerate}
\end{exercice}

\begin{exercice}[Expression écrite : tu deporte favorito]
\textbf{Sujet :} \esp{Habla de tu deporte favorito y explica por qué el deporte es importante para los jóvenes.}

Rédige un texte en espagnol de \textbf{80 à 100 mots}. Tu dois :
\begin{itemize}
\item présenter ton sport ou ton loisir préféré (nom, lieu, partenaires) ;
\item dire à quelle fréquence tu le pratiques (utilise \esp{soler} + infinitif et au moins deux adverbes de fréquence : \esp{siempre}, \esp{normalmente}, \esp{a veces}, \esp{nunca}...) ;
\item exprimer tes goûts avec \esp{gustar} (au moins deux exemples) ;
\item donner au moins deux raisons pour lesquelles le sport est important pour la santé physique et mentale des jeunes ;
\item conclure en encourageant un ami à faire du sport.
\end{itemize}
\end{exercice}

\begin{methode}[Réussir ta production écrite]
\begin{enumerate}
\item \textbf{Prépare un plan} en trois parties : présentation du sport (2--3 phrases), fréquence et goûts (3--4 phrases), importance du sport et conclusion (3--4 phrases).
\item \textbf{Conjugue correctement} : vérifie chaque \esp{gustar} (singulier ou pluriel selon ce qui plaît) et chaque \esp{soler} (\esp{suelo}, \esp{sueles}, \esp{suele}, \esp{solemos}, \esp{soléis}, \esp{suelen}).
\item \textbf{Varie le vocabulaire} de la leçon : \esp{equipo}, \esp{entrenamiento}, \esp{partido}, \esp{ocio}, \esp{salud}, \esp{bienestar}.
\item \textbf{Relis-toi} : accents (\esp{fútbol}, \esp{atletismo}, \esp{campeón}), points d'interrogation et d'exclamation doubles (¿... ? ¡... !), articles devant les sports (\esp{el fútbol}, \esp{la natación}).
\end{enumerate}
\end{methode}

\begin{exemplebox}[Modèle de production (98 mots)]
\esp{Mi deporte favorito es el fútbol. Juego con mis amigos del barrio en el campo del liceo, en Yaoundé. Suelo entrenar dos veces por semana y normalmente jugamos un partido el sábado. Me gusta mucho correr con el balón y me gustan también los entrenamientos, aunque a veces son difíciles. El deporte es muy importante para los jóvenes : es bueno para la salud del cuerpo y también para la salud mental, porque da alegría y amigos. Además, enseña el respeto al otro equipo. ¡Anímate, amigo mío : ven a jugar con nosotros el próximo sábado!}

« Mon sport préféré est le football. Je joue avec mes amis du quartier sur le terrain du lycée, à Yaoundé. J'ai l'habitude de m'entraîner deux fois par semaine et d'habitude nous jouons un match le samedi. J'aime beaucoup courir avec le ballon et j'aime aussi les entraînements, bien qu'ils soient parfois difficiles. Le sport est très important pour les jeunes : il est bon pour la santé du corps et aussi pour la santé mentale, car il donne de la joie et des amis. De plus, il enseigne le respect de l'autre équipe. Allez, mon ami, motive-toi : viens jouer avec nous samedi prochain ! »
\end{exemplebox}

\begin{attention}[Erreurs fréquentes à éviter dans ta rédaction]
\begin{itemize}
\item Ne dis jamais \esp{Yo gusto el deporte} : dis \esp{Me gusta el deporte}.
\item N'oublie pas l'article devant le sport : \esp{el fútbol}, \esp{el atletismo}, \esp{la natación}.
\item \esp{Jugar} se conjugue avec diphtongaison : \esp{juego}, \esp{juegas}, \esp{juega}, mais \esp{jugamos}, \esp{jugáis}.
\item « Par semaine » se traduit \esp{por semana} (et non \esp{para semana}).
\item \esp{Hacer ejercicio} = « faire de l'exercice » ; ne traduis pas mot à mot par \esp{hacer ejercicios} au pluriel dans ce sens.
\end{itemize}
\end{attention}

\newpage

\coursec{Corrigés des exercices}

\begin{corrige}[Exercice 1 — QCM : le lexique du sport]
\begin{enumerate}
\item \textbf{b) un match.} \esp{El partido} est la rencontre sportive entre deux équipes : \esp{un partido de fútbol} « un match de football ».
\item \textbf{a) un supporter.} \esp{El aficionado} est la personne qui aime et suit un sport ou une équipe : \esp{los aficionados del equipo} « les supporters de l'équipe ».
\item \textbf{b) l'entraînement.} Attention au faux ami : \esp{entrenamiento} ne signifie pas « entretien » ; il vient du verbe \esp{entrenar} « s'entraîner ».
\item \textbf{b) gagne une compétition.} Féminin : \esp{la campeona}. Exemple : \esp{Las jugadoras sueñan con ser campeonas del mundo.} « Les joueuses rêvent d'être championnes du monde. »
\item \textbf{b) les loisirs.} \esp{El ocio} désigne le temps libre ; son contraire est \esp{el trabajo} ou \esp{las obligaciones}.
\end{enumerate}
\end{corrige}

\begin{corrige}[Exercice 2 — Vrai ou faux ?]
\begin{enumerate}
\item \esp{Falso}. \esp{Gustar} fonctionne comme « plaire à » : on dit \esp{Me gusta el fútbol} (« le football me plaît »). La forme \esp{Yo gusto el fútbol} est un calque du français, à proscrire absolument.
\item \esp{Verdadero}. \esp{Soler} exprime l'habitude et se construit obligatoirement avec un infinitif : \esp{Suelo correr.} « J'ai l'habitude de courir. »
\item \esp{Falso}. \esp{Suelo} est au présent de l'indicatif : la phrase signifie « j'ai l'habitude de m'entraîner le lundi » (ou « je m'entraîne d'habitude le lundi »), et non un futur. Le futur serait \esp{entrenaré el lunes}.
\item \esp{Verdadero}. Le verbe \esp{gustar} s'accorde avec la chose qui plaît : \esp{los deportes} est pluriel, donc \esp{gustan}. Compare : \esp{Me gusta el deporte} (singulier) / \esp{Me gustan los deportes} (pluriel).
\end{enumerate}
\end{corrige}

\begin{corrige}[Exercice 3 — Vocabulaire : trouve le mot]
\begin{enumerate}
\item \esp{Los aficionados animan a su equipo en el estadio de Yaoundé.} « Les supporters encouragent leur équipe au stade de Yaoundé. »
\item \esp{El sábado hay un partido de fútbol entre los dos liceos.} « Samedi, il y a un match de football entre les deux lycées. »
\item \esp{Samuel Eto'o fue campeón de África con los Leones Indomables.} « Samuel Eto'o a été champion d'Afrique avec les Lions Indomptables. »
\item \esp{El atletismo incluye la carrera, el salto y el lanzamiento.} « L'athlétisme comprend la course, le saut et le lancer. »
\item \esp{El deporte es mi actividad de ocio favorita.} « Le sport est mon activité de loisir préférée. »
\end{enumerate}
\textbf{Points de vigilance} : \esp{aficionados} au pluriel (accord avec \esp{los} et \esp{animan}) ; \esp{campeón} prend un accent sur la dernière syllabe ; \esp{atletismo} prend un accent sur la première syllabe (\emph{a-tlé-tis-mo}).
\end{corrige}

\begin{corrige}[Exercice 4 — Conjugaison : gustar et soler]
\begin{enumerate}
\item \esp{A mí me gusta mucho el baloncesto.} « J'aime beaucoup le basket-ball. » (\esp{gusta} : singulier, car ce qui plaît est \esp{el baloncesto}.)
\item \esp{A mis amigos les gustan los partidos de fútbol.} « Mes amis aiment les matchs de football. » (\esp{gustan} : pluriel, car ce qui plaît est \esp{los partidos} ; pronom \esp{les} pour \esp{a mis amigos}.)
\item \esp{¿A ti te gusta nadar ?} « Est-ce que tu aimes nager ? » (Devant un infinitif, \esp{gustar} est toujours au singulier.)
\item \esp{Nosotros solemos entrenar los martes.} « Nous avons l'habitude de nous entraîner le mardi. » (\esp{Soler} est un verbe en \emph{o $\rightarrow$ ue} comme \esp{volver}.)
\item \esp{Ella suele correr por la mañana.} « Elle a l'habitude de courir le matin. »
\item \esp{¿Vosotros soléis ver los partidos en la televisión ?} « Avez-vous l'habitude de regarder les matchs à la télévision ? » (Pas de diphtongaison à la 2\up{e} personne du pluriel : \esp{soléis}, avec accent.)
\end{enumerate}
\end{corrige}

\begin{corrige}[Exercice 5 — Texte à trous : gustar et soler]
Réponses : 1. \esp{gusta} \quad 2. \esp{gusta} \quad 3. \esp{suelo} \quad 4. \esp{solemos} \quad 5. \esp{gusta} \quad 6. \esp{suelen} \quad 7. \esp{gustan} \quad 8. \esp{Sueles} \quad 9. \esp{gustan} \quad 10. \esp{solemos}.

Texte complet : \esp{En mi familia, todos amamos el deporte. A mí me gusta mucho el fútbol y a mi hermana le gusta el atletismo. Yo suelo entrenar tres veces por semana con mi equipo. Los fines de semana, nosotros solemos jugar un partido en el barrio. A mis padres les gusta caminar por la mañana : ellos suelen levantarse muy temprano. ¿Y tú? ¿Te gustan los deportes de equipo? ¿Sueles hacer ejercicio? A mis amigos y a mí nos gustan también los videojuegos, pero nosotros no solemos jugar durante la semana.}

Traduction : « Dans ma famille, nous aimons tous le sport. Moi, j'aime beaucoup le football et ma sœur aime l'athlétisme. J'ai l'habitude de m'entraîner trois fois par semaine avec mon équipe. Le week-end, nous jouons d'habitude un match dans le quartier. Mes parents aiment marcher le matin : ils ont l'habitude de se lever très tôt. Et toi ? Tu aimes les sports d'équipe ? Tu as l'habitude de faire de l'exercice ? Mes amis et moi aimons aussi les jeux vidéo, mais nous n'avons pas l'habitude de jouer pendant la semaine. »

\textbf{Points de vigilance} :
\begin{itemize}
\item Questions 1, 2 et 5 : ce qui plaît est singulier ou infinitif (\esp{el fútbol}, \esp{el atletismo}, \esp{caminar}) $\rightarrow$ \esp{gusta}.
\item Questions 7 et 9 : ce qui plaît est pluriel (\esp{los deportes}, \esp{los videojuegos}) $\rightarrow$ \esp{gustan} ; à la question 9, le pronom est \esp{nos} car \esp{a mis amigos y a mí} = « nous ».
\item Question 6 : le sujet est \esp{ellos} (\esp{mis padres}) $\rightarrow$ \esp{suelen}.
\end{itemize}
\end{corrige}

\begin{corrige}[Exercice 6 — Thème : traduis en espagnol]
\begin{enumerate}
\item \esp{Me gusta mucho el deporte.} (Jamais \esp{Yo gusto} : calque du français.)
\item \esp{Mi equipo se entrena dos veces por semana.} ou \esp{Mi equipo entrena dos veces por semana.} (« par semaine » = \esp{por semana}, avec \esp{por}, jamais \esp{para}.)
\item \esp{Solemos jugar un partido los sábados.} ou \esp{Normalmente jugamos un partido los sábados.} (\esp{los sábados} au pluriel = « le samedi », habitude.)
\item \esp{A mi amigo no le gusta nada perder.} (Négation renforcée : \esp{no le gusta nada} ; ne pas oublier le pronom \esp{le} même avec \esp{a mi amigo}.)
\item \esp{El deporte es bueno para la salud mental.} (\esp{ser} + \esp{bueno para} : « bon pour » ; ici \esp{para} exprime le bénéficiaire.)
\end{enumerate}
\end{corrige}

\begin{corrige}[Exercice 7 — Reformulation : l'habitude avec soler]
\begin{enumerate}
\item \esp{Solemos jugar al fútbol los sábados.} « Nous avons l'habitude de jouer au football le samedi. »
\item \esp{Mi hermana suele nadar en la piscina los domingos.} « Ma sœur a l'habitude de nager à la piscine le dimanche. »
\item \esp{Los aficionados suelen ver el partido en la televisión.} « Les supporters ont l'habitude de regarder le match à la télévision. »
\item \esp{Suelo hacer ejercicio después de las clases.} « J'ai l'habitude de faire de l'exercice après les cours. »
\item \esp{Los lunes, mi equipo suele entrenar en el estadio.} « Le lundi, mon équipe a l'habitude de s'entraîner au stade. »
\end{enumerate}
\textbf{Méthode} : on supprime l'adverbe d'habitude (\esp{normalmente}, \esp{de costumbre}, \esp{generalmente}, \esp{habitualmente}, \esp{como siempre}), on conjugue \esp{soler} à la personne du sujet, et le verbe principal passe à l'infinitif. Attention à la diphtongaison : \esp{suelo}, \esp{suele}, \esp{suelen}, mais \esp{solemos}.
\end{corrige}

\begin{corrige}[Exercice 8 — Appariements : mots et définitions]
1 -- b ; 2 -- a ; 3 -- d ; 4 -- c ; 5 -- e ; 6 -- f.

Vérification en contexte : \esp{el entrenamiento} « la période de préparation physique » ; \esp{el aficionado} « la personne qui adore un sport » ; \esp{el campeón} « celui qui remporte la compétition » ; \esp{el partido} « la rencontre entre deux équipes » ; \esp{el ocio} « le temps libre » ; \esp{el equipo} « le groupe de joueurs ».
\end{corrige}

\begin{corrige}[Exercice 9 — Compréhension de l'écrit : le football féminin en Espagne]
\textbf{Barème indicatif (sur 20)} : A. 4 $\times$ 1,5 pt (0,5 pt pour \esp{Verdadero}/\esp{Falso} + 1 pt pour la citation exacte) = 6 pts ; B. 3 $\times$ 2 pts = 6 pts ; C. 2 + 4 + 2 = 8 pts.

\textbf{A. Comprensión}
\begin{enumerate}
\item \esp{Falso} : \esp{« Hace veinte años, pocas niñas españolas jugaban al fútbol »} (peu de filles jouaient, et non beaucoup).
\item \esp{Verdadero} : \esp{« la televisión emite cada vez más partidos femeninos »}.
\item \esp{Falso} : \esp{« Suelo entrenar cuatro veces por semana »} (elle s'entraîne quatre fois par semaine ; le samedi, ce sont les matchs).
\item \esp{Falso} : \esp{« Sin embargo, quedan problemas »} (il reste des problèmes : salaires plus bas, terrains insuffisants).
\end{enumerate}

\textbf{B. Preguntas abiertas} (réponses modèles)
\begin{enumerate}
\item \esp{Las tres razones son : las jugadoras españolas han ganado títulos importantes y son modelos para las jóvenes ; la televisión emite cada vez más partidos femeninos ; y muchas escuelas han creado equipos de fútbol para las alumnas.}
\item \esp{A Lucía le gustan los entrenamientos y le encantan los partidos del sábado. El fútbol le da amigas, salud y alegría.}
\item \esp{Las jugadoras profesionales ganan menos dinero que los hombres y algunos clubes no tienen buenos campos para las chicas.}
\end{enumerate}

\textbf{C. Lengua}
\begin{enumerate}
\item Imparfait : \esp{jugaban}, \esp{era}, \esp{practicaban}. Emploi : l'imparfait décrit ici une situation habituelle et durable dans le passé (« il y a vingt ans, les filles jouaient..., le football était... »), opposée à la situation présente introduite par \esp{hoy}.
\item \esp{Gustar} : \esp{« le gusta mucho su deporte »} « son sport lui plaît beaucoup » ; \esp{« Me gustan los entrenamientos »} « j'aime les entraînements ». \esp{Soler} : \esp{« los aficionados suelen verlos con entusiasmo »} « les supporters ont l'habitude de les regarder avec enthousiasme » ; \esp{« Suelo entrenar cuatro veces por semana »} « j'ai l'habitude de m'entraîner quatre fois par semaine ».
\item « les supporters » = \esp{los aficionados} ; « l'entraînement » = \esp{el entrenamiento} ; « championnes » = \esp{campeonas}.
\end{enumerate}

\textbf{Points de vigilance} : pour la partie A, recopie la citation exacte entre guillemets, sans la traduire ; pour la partie B, réponds par des phrases complètes en reprenant les mots de la question ; à la question C.1, ne confonds pas imparfait (\esp{jugaban}) et passé simple.
\end{corrige}

\begin{corrige}[Exercice 10 — Traduction : thème et version]
\textbf{Barème indicatif (sur 20)} : version 10 pts (2 pts par phrase) ; thème 10 pts (2 pts par phrase : 1 pt pour la structure, 1 pt pour la grammaire et l'orthographe).

\textbf{A. Versión} — « Diego est un adolescent mexicain de seize ans. Il aime beaucoup les sports, surtout le football. Il a l'habitude de s'entraîner avec son équipe trois fois par semaine, après les cours. D'habitude, le samedi, il joue un match avec ses amis au parc. Diego n'aime pas du tout perdre, mais il respecte toujours l'autre équipe. »

Pièges à éviter :
\begin{itemize}
\item \esp{Le gustan mucho los deportes} ne se traduit pas par « il leur plaît » mais par « il aime beaucoup les sports » : \esp{los deportes} est le sujet grammatical, d'où \esp{gustan} au pluriel.
\item \esp{respeta al otro equipo} : \esp{a} + \esp{el} se contracte en \esp{al} ; c'est la \emph{a} personnelle devant un complément de personne.
\item \esp{especialmente} = « surtout, en particulier », et non « spécialement » au sens de « exprès ».
\end{itemize}

\textbf{B. Thème}
\begin{enumerate}
\item \esp{Me gusta mucho el deporte.}
\item \esp{Mi equipo se entrena dos veces por semana.}
\item \esp{Solemos jugar un partido los sábados.} (ou \esp{Normalmente jugamos un partido los sábados.})
\item \esp{A mi amigo no le gusta nada perder.}
\item \esp{El deporte es bueno para la salud mental.}
\end{enumerate}

\textbf{Points de vigilance} : ne jamais écrire \esp{Yo gusto} ; « par semaine » = \esp{por semana} ; « bon pour » = \esp{bueno para} ; penser à l'article devant le sport (\esp{el deporte}).
\end{corrige}

\begin{corrige}[Exercice 11 — Expression écrite : tu deporte favorito]
\textbf{Barème indicatif (sur 20)} : respect de la consigne et du nombre de mots (80--100 mots) : 4 pts ; correction grammaticale (\esp{gustar}, \esp{soler}, conjugaisons, articles) : 6 pts ; richesse du lexique de la leçon : 4 pts ; organisation et cohérence du texte : 4 pts ; orthographe et accents : 2 pts.

\textbf{Grille d'attentes} : le texte doit contenir au moins deux emplois de \esp{gustar} correctement accordés, un emploi de \esp{soler} + infinitif, deux adverbes de fréquence, deux raisons liées à la santé physique et mentale, et une conclusion encourageante.

\textbf{Proposition de rédaction modèle (98 mots)} :

\esp{Mi deporte favorito es el fútbol. Juego con mis amigos del barrio en el campo del liceo, en Yaoundé. Suelo entrenar dos veces por semana y normalmente jugamos un partido el sábado. Me gusta mucho correr con el balón y me gustan también los entrenamientos, aunque a veces son difíciles. El deporte es muy importante para los jóvenes : es bueno para la salud del cuerpo y también para la salud mental, porque da alegría y amigos. Además, enseña el respeto al otro equipo. ¡Anímate, amigo mío : ven a jugar con nosotros el próximo sábado!}

Traduction : « Mon sport préféré est le football. Je joue avec mes amis du quartier sur le terrain du lycée, à Yaoundé. J'ai l'habitude de m'entraîner deux fois par semaine et d'habitude nous jouons un match le samedi. J'aime beaucoup courir avec le ballon et j'aime aussi les entraînements, bien qu'ils soient parfois difficiles. Le sport est très important pour les jeunes : il est bon pour la santé du corps et aussi pour la santé mentale, car il donne de la joie et des amis. De plus, il enseigne le respect de l'autre équipe. Allez, mon ami, motive-toi : viens jouer avec nous samedi prochain ! »

\textbf{Points de vigilance} :
\begin{itemize}
\item Jamais \esp{Yo gusto el deporte} : toujours \esp{Me gusta el deporte}.
\item Article obligatoire devant le sport : \esp{el fútbol}, \esp{el atletismo}, \esp{la natación}.
\item \esp{Jugar} se conjugue avec diphtongaison : \esp{juego}, \esp{juegas}, \esp{juega}, mais \esp{jugamos}, \esp{jugáis}.
\item « Par semaine » = \esp{por semana} (et non \esp{para semana}).
\item « Faire de l'exercice » = \esp{hacer ejercicio} (au singulier dans ce sens).
\item N'oublie pas les signes doubles : ¿... ? et ¡... ! ni les accents : \esp{fútbol}, \esp{campeón}, \esp{atletismo}.
\end{itemize}
\end{corrige}

\newpage

\coursec{Fiche bilan}

\begin{popfigure}
\begin{tikzpicture}[
  mindmap,
  grow cyclic,
  every node/.style={concept, align=center, font=\small},
  root concept/.append style={concept color=popblue, font=\bfseries\small, text=white, minimum size=3.2cm},
  level 1 concept/.append style={level distance=4.6cm, sibling angle=51.4, font=\footnotesize, minimum size=2.4cm},
  level 2 concept/.append style={level distance=2.6cm, font=\scriptsize, minimum size=1.9cm}
]
\node[root concept, align=center]{El deporte,\\el ocio y\\el bienestar}
  child[concept color=poporange]{ node[align=center]{Lexique\\el deporte}
    child[concept color=poporange!60]{ node[align=center]{equipo,\\partido}}
    child[concept color=poporange!60]{ node[align=center]{entrenar,\\campeón}}
  }
  child[concept color=popgreen]{ node[align=center]{Lexique\\el ocio}
    child[concept color=popgreen!60]{ node[align=center]{aficionado,\\pasatiempo}}
  }
  child[concept color=poppurple]{ node[align=center]{Gustar\\(le goût)}
    child[concept color=poppurple!60]{ node[align=center]{me gusta\\+ inf.}}
    child[concept color=poppurple!60]{ node[align=center]{a mí me,\\a ti te}}
  }
  child[concept color=poppink]{ node[align=center]{Soler + inf.\\(habitude)}
    child[concept color=poppink!60]{ node[align=center]{suelo,\\sueles...}}
  }
  child[concept color=popgold]{ node {Fréquence}
    child[concept color=popgold!60]{ node[align=center]{siempre,\\a menudo}}
    child[concept color=popgold!60]{ node[align=center]{una vez\\a la semana}}
  }
  child[concept color=popblue!70]{ node[align=center]{Santé et\\bienestar}
    child[concept color=popblue!40]{ node[align=center]{salud física\\y mental}}
  }
  child[concept color=popgreen!80!black]{ node[align=center]{Sport au\\Cameroun et\\en Hispanie}
    child[concept color=popgreen!50]{ node[align=center]{fútbol,\\atletismo}}
  };
\end{tikzpicture}
\legende{Carte mentale de la leçon : le sport, les loisirs et le bien-être}
\end{popfigure}

\begin{bilan}
\textbf{1. Lexique clé à connaître par cœur}

\begin{center}
\begin{tabularx}{\linewidth}{|X|X|X|X|}
\hline
\rowcolor{popblueL} \textbf{Español} & \textbf{Français} & \textbf{Español} & \textbf{Français} \\
\hline
el deporte & le sport & el partido & le match \\
\hline
el fútbol & le football & el equipo & l'équipe \\
\hline
el atletismo & l'athlétisme & el campeón / la campeona & le champion / la championne \\
\hline
el entrenamiento & l'entraînement & el aficionado & le supporter, l'amateur \\
\hline
entrenar(se) & s'entraîner & el ocio & les loisirs \\
\hline
practicar un deporte & pratiquer un sport & el bienestar & le bien-être \\
\hline
la salud & la santé & mantenerse en forma & se maintenir en forme \\
\hline
\end{tabularx}
\end{center}

\textbf{2. Grammaire : les règles à retenir}

\begin{center}
\begin{tabularx}{\linewidth}{|X|X|}
\hline
\rowcolor{popgreenL} \textbf{Règle} & \textbf{Exemple} \\
\hline
\esp{gustar} se construit à l'envers du français : la chose aimée est le sujet ; la personne est un complément (\esp{me, te, le, nos, os, les}). & \esp{Me gusta el fútbol.} « J'aime le football. » / \esp{Me gustan los deportes.} « J'aime les sports. » \\
\hline
Devant \esp{gustar} + infinitif, le verbe reste au singulier. & \esp{Me gusta correr y nadar.} « J'aime courir et nager. » \\
\hline
Pour insister ou préciser la personne : \esp{a mí, a ti, a él, a nosotros...} + pronom. & \esp{A mí me gusta el atletismo.} « Moi, j'aime l'athlétisme. » \\
\hline
Négation : \esp{no} se place avant le pronom. & \esp{No me gusta el boxeo.} « Je n'aime pas la boxe. » \\
\hline
\esp{soler} + infinitif = « avoir l'habitude de » ; verbe irrégulier (o $\rightarrow$ ue) : \esp{suelo, sueles, suele, solemos, soléis, suelen}. & \esp{Suelo entrenar los sábados.} « J'ai l'habitude de m'entraîner le samedi. » \\
\hline
Marques de fréquence : \esp{siempre, a menudo, normalmente, a veces, casi nunca, nunca, una vez a la semana, dos veces al mes}. & \esp{Juego al fútbol dos veces a la semana.} « Je joue au football deux fois par semaine. » \\
\hline
\end{tabularx}
\end{center}

\textbf{3. Méthodes express}
\begin{itemize}
  \item \textbf{Lire un texte sur le sport} : repérer le thème, les sports cités, les arguments (santé physique, santé mentale, vie sociale), puis reformuler en français avec ses propres mots.
  \item \textbf{Parler de ses loisirs} : structure en trois temps : goût (\esp{me gusta...}), fréquence (\esp{suelo... / dos veces a la semana}), bénéfice (\esp{es bueno para la salud}).
  \item \textbf{Traduire (thème)} : « j'aime » $\rightarrow$ \esp{me gusta} (jamais \esp{yo gusto}) ; « par semaine » $\rightarrow$ \esp{a la semana} ; « faire du sport » $\rightarrow$ \esp{hacer deporte} ou \esp{practicar deporte}.
  \item \textbf{Encourager la pratique sportive} : utiliser l'impératif ou \esp{hay que + infinitif} : \esp{¡Haz deporte! Hay que moverse todos los días.}
\end{itemize}

\textbf{4. Pièges des francophones}
\begin{itemize}
  \item \esp{Me gusta los deportes} est faux : accord avec la chose aimée $\rightarrow$ \esp{me gustan los deportes}.
  \item \esp{deporte} est masculin : \esp{el deporte}, \esp{un deporte}.
  \item Ne pas confondre \esp{partido} (match) et \esp{partida} (partie de cartes) ; « jouer à » = \esp{jugar a + article} : \esp{juego al fútbol}.
  \item \esp{soler} s'emploie toujours suivi d'un infinitif, jamais seul.
\end{itemize}
\end{bilan}

\begin{aretenir}[Checklist avant le Bac]
\begin{itemize}
  \item $\square$ Je connais le lexique du sport et des loisirs : \esp{deporte, equipo, partido, entrenamiento, campeón, aficionado, ocio, bienestar}.
  \item $\square$ Je sais conjuguer \esp{gustar} au présent avec tous les pronoms : \esp{me, te, le, nos, os, les}.
  \item $\square$ Je fais correctement l'accord de \esp{gustar} : \esp{me gusta el fútbol} mais \esp{me gustan los deportes}.
  \item $\square$ Je sais renforcer avec \esp{a mí, a ti, a él...} sans oublier le pronom obligatoire.
  \item $\square$ Je sais conjuguer \esp{soler} au présent : \esp{suelo, sueles, suele, solemos, soléis, suelen}.
  \item $\square$ J'emploie \esp{soler} uniquement devant un infinitif.
  \item $\square$ Je maîtrise les marques de fréquence : \esp{siempre, a menudo, a veces, nunca, dos veces a la semana}.
  \item $\square$ Je sais dire « jouer à » : \esp{jugar al fútbol}, avec la contraction \esp{al}.
  \item $\square$ Je sais parler de mes loisirs en trois phrases : goût, fréquence, bénéfice pour la santé.
  \item $\square$ Je sais traduire « faire du sport » par \esp{hacer deporte} et « se maintenir en forme » par \esp{mantenerse en forma}.
  \item $\square$ Je peux citer des arguments sur l'importance du sport pour la santé physique et mentale.
  \item $\square$ Je connais des exemples sportifs du Cameroun et du monde hispanophone pour illustrer un commentaire.
\end{itemize}
\end{aretenir}

\findecours

\end{document}
